\chapter{链接、装载与系统软件边界}

\begin{keyidea}
硬件执行的是确定地址上的机器码，软件工具链负责把名字、段和重定位变成这种状态。本章先沿目标文件、链接、装载和异常入口建立主线，再用固件镜像作为无操作系统环境的对照案例。
\end{keyidea}

\begin{longtable}{L{0.2\linewidth}L{0.4\linewidth}L{0.3\linewidth}}
\toprule
阶段 & 做什么 & 对应的硬件动作 \\
\midrule
编译 & 源码变目标文件 & 生成指令编码和重定位记录 \\\tablerowrule
链接 & 合并目标文件、解析符号 & 把符号名变成确定地址 \\\tablerowrule
装载 & 把程序映像放进内存 & 建立页表、映射段 \\\tablerowrule
运行 & 请求内核服务 & 系统调用陷入异常入口 \\
\bottomrule
\end{longtable}

\begin{keyidea}
链接和装载不是编译器内部魔法，而是“什么时候把哪个名字定成哪个地址”的时间表。读懂这张时间表，静态库、动态库、PLT、系统调用这些名词都会各就各位。
\end{keyidea}

\section{一、目标文件与符号}

本节回答一个问题：编译器交出来的目标文件，里面到底装
了什么。前面二十四章讲的都是硬件怎样执行指令，从本节开
始换一个视角：指令被执行之前，先作为字节躺在文件里。
目标文件就是这批字节的容器，读懂它只需要一张“哪段字
节是什么、给谁看”的地图。

\topic{ELF 是目标文件的通用容器}
x86-64 Linux 上，编译器每处理一个源文件就交出一个 ELF
目标文件，链接器拼出的可执行文件是 ELF，共享库也是
ELF。ELF（Executable and Linkable Format）是这一整条流
程的通用容器：同一套字节排布约定，既能让链接器读到“这
里有一段代码、那里有一张符号表”，也能让内核装载器读到
“把哪一段映射到哪个虚拟地址”——后者是第三节的事。

容器从固定格式的头部开始。前 16 字节是身份区：头 4 字
节是魔数 \code{0x7F} 加字母 \code{E}、\code{L}、
\code{F}，任何工具拿到文件先核对这个签名，对不上就立
刻拒收；随后两个字节说明类别（32 位还是 64 位）和字节
序（大端还是小端）——只凭这两个字段，解析程序就知道
后面所有多字节字段该按什么宽度、什么顺序读。头部其余
字段里最重要的几个列在下表。

\begin{longtable}{L{0.14\linewidth}L{0.4\linewidth}L{0.32\linewidth}}
\toprule
字段 & 内容 & 谁来读 \\
\midrule
\code{e\_type} & 文件类型：ET\_REL / ET\_EXEC / ET\_DYN & 一切工具先按它分流 \\\tablerowrule
\code{e\_machine} & 指令集编号，x86-64 为 62 & 链接器据此拒绝混链不同 ISA \\\tablerowrule
\code{e\_entry} & 进程入口的虚拟地址 & 装载器；可重定位文件恒为 0 \\\tablerowrule
\code{e\_phoff} & 程序头表（段表）的文件偏移 & 装载器 \\\tablerowrule
\code{e\_shoff} & 节头表（节表）的文件偏移 & 链接器与观察工具 \\
\bottomrule
\end{longtable}

三种类型对应程序出生过程的三个形态。可重定位文件
（\code{.o}）是半成品：代码和数据都在，但地址一律未
定，凡涉及地址的字段全部留空，另附一叠“这里以后要
填”的重定位记录。可执行文件是成品：每个节的虚拟地址
已定，入口字段指向第一条要执行的指令（现代发行版常把
可执行文件编成 ET\_DYN 形态的 PIE，把定址推迟到装载
时，第三节展开）。共享库（\code{.so}）介于两者之间：
代码数据齐全，却故意不绑定固定装载地址，好让多个进程
各自把它映射到不同位置。三种形态共用一套头部与节表格
式，差别只在几个字段的取值和洞里填没填。

头部之后是两份“目录”，给出同一批字节的两种看法。节
表（section header table）逐项列出每个节的名字、类
型、大小、文件偏移和属性标志，这是链接视图：链接器按
节合并、按节定址。程序头表（program header table，常
称段表）把属性相同的若干节归并成段，逐项给出段的虚拟
地址、文件范围和读写执行权限，这是执行视图：装载器只
看它，根本不关心 \code{.text} 叫什么名字。可重定位文
件通常只有节表——它离执行还早；可执行文件两份都有，
装载时只用段表。同一份字节，链接器和装载器各取所需，
这正是“通用容器”的含义。

两份目录的分工画成对视图更
直观：左边的节按名字与用途
分，右边的段按权限归并，箭
头就是链接器的归并动作：

\begin{pdfdiagram}
  % 左列：链接视图（节表）
  \node[hw label] at (1.5,3.3) {链接视图：节表};
  \node[hw compute, minimum width=2.6cm] (st) at (1.5,2.6) {.text};
  \node[hw compute, minimum width=2.6cm] (sr) at (1.5,1.8) {.rodata};
  \node[hw compute, minimum width=2.6cm] (sd) at (1.5,1.0) {.data};
  \node[hw memory, minimum width=2.6cm, dashed] (sb) at (1.5,0.2) {.bss};
  \node[hw state, minimum width=2.6cm, align=center] (ss) at (1.5,-0.8) {.symtab / .strtab\\符号与字符串};
  % 右列：执行视图（段表）
  \node[hw label] at (9.3,3.3) {执行视图：段表};
  \node[hw state, minimum width=3.2cm, align=center] (lx) at (9.3,2.6) {LOAD：读/执行};
  \node[hw state, minimum width=3.2cm, align=center] (lr) at (9.3,1.8) {LOAD：只读};
  \node[hw state, minimum width=3.2cm, align=center] (lw) at (9.3,0.6) {LOAD：读写};
  \node[hw label, align=center] (ln) at (9.3,-0.8) {不进段表\\装载器不看};
  % 归并箭头：同级水平直连
  \draw[hw bus] (st.east) -- (lx.west);
  \draw[hw bus] (sr.east) -- (lr.west);
  % .data 与 .bss 共享一小段干线，汇入同一个读写段
  \draw[hw bus] (sd.east) -- (4.9,1.0) -- (4.9,0.6) -- (lw.west);
  \draw[hw bus] (sb.east) -- (4.3,0.2) -- (4.3,0.6) -- (4.9,0.6);
  \fill[BookMuted] (4.9,0.6) circle (0.05);
  \node[hw label, fill=white, inner sep=1pt] at (6.2,0.32) {按属性归并};
  % 符号表与字符串表只供链接期
  \draw[hw return] (ss.east) -- (ln.west);
\end{pdfdiagram}
\diagramnote{同一份字节的两种看法。链接视图按节组织：名字、类型、大小与属性逐项登记在节表里，链接器按节合并、按节定址；执行视图把属性相同的节归并成 LOAD 段，只保留虚拟地址、文件范围与读写执行权限，装载器只看这一份。可重定位文件通常只有节表；可执行文件两份都有，装载时只用段表——\code{strip} 把节表整个剥掉，程序照样运行。}

把这份容器按文件偏移画成竖向排布，头部、两份目录与各节的位置关系便一目了然；\code{readelf} 的每个常用选项，恰好对应图中的一块：

\begin{pdfdiagram}
  % 主体：按文件偏移自上而下排布的字节布局
  \node[hw state, minimum width=3.4cm] (ehdr) at (1.6,4.5) {ELF 头（64 字节）\\魔数、e\_phoff、e\_shoff};
  \node[hw state, minimum width=3.4cm] (phdr) at (1.6,3.3) {程序头表（段表）\\执行视图，装载器读};
  \node[hw compute, minimum width=3.4cm] (sect) at (1.6,2.1) {.text：指令编码};
  \node[hw compute, minimum width=3.4cm] (srod) at (1.6,1.2) {.rodata：只读常量};
  \node[hw compute, minimum width=3.4cm] (sdat) at (1.6,0.3) {.data：已初始化数据};
  \node[hw memory, minimum width=3.4cm, dashed] (sbss) at (1.6,-0.6) {.bss：不占文件字节};
  \node[hw memory, minimum width=3.4cm] (ssym) at (1.6,-1.5) {.symtab：符号表};
  \node[hw state, minimum width=3.4cm] (shtab) at (1.6,-2.4) {节头表（节表）\\链接视图，链接器读};
  % 左侧：文件偏移增长方向
  \draw[BookMuted, line width=0.6pt, ->] (-0.6,4.5) -- (-0.6,-2.4);
  \node[hw label, anchor=east] at (-0.75,1.0) {文件偏移增大};
  % 右侧：readelf 各选项（虚线指出读取的部分）
  \node[hw label, anchor=west] (rh) at (4.3,4.5) {\code{readelf -h}：读 ELF 头};
  \node[hw label, anchor=west] (rl) at (4.3,3.3) {\code{readelf -l}：读段表};
  \node[hw label, anchor=west] (rs) at (4.3,-1.5) {\code{readelf -s}：读符号表};
  \node[hw label, anchor=west] (rS) at (4.3,-2.4) {\code{readelf -S}：读节表};
  \draw[hw return] (rh.west) -- (ehdr.east);
  \draw[hw return] (rl.west) -- (phdr.east);
  \draw[hw return] (rs.west) -- (ssym.east);
  \draw[hw return] (rS.west) -- (shtab.east);
\end{pdfdiagram}
\diagramnote{ELF 文件按字节偏移的排布（自上而下偏移增大）。ELF 头中的 \code{e\_phoff} 与 \code{e\_shoff} 分别给出段表与节表的文件偏移；\code{.bss} 是 NOBITS 类型，节表里只登记大小，文件里一个字节也不占，图中用虚线框表示。右侧标出 \code{readelf} 各选项读取的部分：同一份字节，链接器取节表与符号表，装载器只取段表。}

最后要戳破一层神秘感：ELF 没有任何魔法，它只是按约定
排布的字节。文件自己并不“知道”节表在哪，是工具按
\code{e\_shoff} 找过去的；头部里某个偏移写错，读出来
的就是垃圾，没有类型系统代为保证正确性。本书反复出现的主题
在这里再次出现：软与硬之间、工具与文件之间，一切协作
都靠约定，不靠神秘力量。

\topic{常见节的分工}
打开任意一个目标文件的节表，总能见到几张熟面孔。它们
的分工不是编译器的发明，而是把“用途与权限相同”的字
节归在一起：用途决定放在哪，权限决定内存页怎么保护。

\code{.text} 装指令字节，属性是分配、只读、可执行；
\code{.rodata} 装只读常量——字符串字面量、\code{const}
全局量、编译器为 \code{switch} 生成的跳转表，属性是分
配、只读；\code{.data} 装初值非零的全局变量和静态局部
变量，文件里存的就是初值本身，属性是分配、可写；
\code{.bss} 装未初始化或初值恰为零的全局量与静态量，
属性同样是分配、可写，但它在文件里一个字节都不占。

\begin{longtable}{L{0.12\linewidth}L{0.34\linewidth}L{0.16\linewidth}L{0.24\linewidth}}
\toprule
节 & 装什么 & 占文件字节 & 内存权限 \\
\midrule
\code{.text} & 指令编码 & 占 & 只读、可执行 \\\tablerowrule
\code{.rodata} & 字符串、const 常量、跳转表 & 占 & 只读 \\\tablerowrule
\code{.data} & 初值非零的全局/静态变量 & 占，内容即初值 & 可读写 \\\tablerowrule
\code{.bss} & 未初始化或初值为零的变量 & 不占 & 可读写，装载时清零 \\
\bottomrule
\end{longtable}

\code{.bss} 省空间的机制一句话就能说清：它的节头类型
是 NOBITS，节表里只有名字、大小和对齐要求，没有对应
的文件范围；装载器读到它时不复制任何字节，直接在内存
里划出大小相等的一片并填零。这与装载器的既有动作正好
合拍——内核给新映射提供的本来就是填零的页（第三节的
装载过程还会用到），所以
\code{.bss} 的清零几乎不花额外代价。

省多少可以算出来。假设程序里有一个 1 MiB 的未初始化
静态数组：若放进 \code{.data}，目标文件要多存 1 MiB
内容，而这 1 MiB 全是零字节，运载的信息量却是零；放
进 \code{.bss}，文件只多一行几十字节的节头记录，内存
里照样得到 1 MiB 全零空间。文件的职责是保存非零信
息；全零内容用“长度加填零”几个字就说完了。C 的规则
也顺着这个机制设计：显式初始化为零的变量同样进
\code{.bss}——\code{int x = 0;} 与 \code{int x;} 同等
待遇，只有初值非零才需要 \code{.data} 里的真实字节。

节的属性标志不是注释，它们会一路传到硬件。链接器按属
性把节归并成段，段头里的读写执行标志最终被装载器转成
页表项的权限位：于是 \code{.text} 所在页不可写、
\code{.data} 所在页不可执行。这些保护不是操作系统额
外的好意，而是目标文件里早就写明的属性，经链接、装载
两步原样传给 MMU。存储与设备事务相关章节问页表项权限位的来历，答案
有一半就在这里。

\topic{符号表是名字到位置的映射}
节回答了“字节是什么”，还没回答“名字在哪”。
\code{main.c} 里调用 \code{add}，编译器生成
\code{main.o} 时并不知道 \code{add} 的地址——它可能
根本不在本文件里。目标文件的处理办法是：把每个名字连
同它的状态登记进符号表 \code{.symtab}，名字与位置的
对应关系全部走这张表。

符号表每个条目记录：名字（字符串存在 \code{.strtab}，
条目里存偏移）、值（\code{st\_value}：可重定位文件里
是节内偏移，可执行文件里是虚拟地址）、大小、类型（函
数 FUNC、数据 OBJECT）、绑定属性，以及“定义在哪个
节”。最后一项是定义与引用的分界线：纯引用的条目这一
栏是 UND（未定义），意思是“这个名字本文件要用，定义
在别处”；有定义的条目这一栏是本文件某个节的编号。

绑定属性分三档。LOCAL 是文件内私有：\code{static} 修
饰的函数和变量都是 LOCAL，不参与跨文件解析，别的文件
有同名符号也互不干扰。GLOBAL 是默认档：普通函数与全
局变量，在全程序范围内参与解析。WEAK 是可以商量档：
弱定义遇到同名 GLOBAL 强定义时自动让位，没有强定义时
才用它——库常用弱符号提供可替换的默认实现；弱引用更
宽容，找不到定义也不报错，运行时取到 0，由程序自己判
空。

\begin{longtable}{L{0.2\linewidth}L{0.28\linewidth}L{0.38\linewidth}}
\toprule
类别 & 怎么产生 & 链接器怎么处理 \\
\midrule
LOCAL & \code{static} 函数/变量 & 不进全局解析，只供本文件与调试器 \\\tablerowrule
GLOBAL 定义 & 普通函数、全局变量 & 全程序只允许一个强定义 \\\tablerowrule
GLOBAL 引用（UND） & 调用或访问外部名字 & 必须在某个输入文件里找到定义 \\\tablerowrule
WEAK 定义 & \code{\_\_attribute\_\_((weak))} & 有强定义则让位，否则用它 \\
\bottomrule
\end{longtable}

链接器两个经典报错都来自这张表。未定义引用：扫完所有
输入文件，某个 UND 条目始终没找到对应的 GLOBAL 定义，
于是报告 \code{undefined reference to `foo'}——这并非不可解释的现象，而是符号表查询失败的直接结果。多重定义：同一名字
出现两个 GLOBAL 强定义，链接器无法裁决哪个算数，只能
拒绝。在 C 头文件里误定义非 \code{static}、非
\code{inline} 的变量，被多个源文件包含后各产生一个强
定义，正是这个错误的头号来源。

把符号表看成“接口的物理形态”很有用：目标文件对外提
供什么（GLOBAL 定义）、向外索取什么（UND 引用），表
上一目了然。第二节的符号解析、第三节的动态链接，处理
的始终是这三个问题：名字是谁、值是多少、定义在不在
场。

\topic{用 readelf/objdump 观察真实目标文件}
这些结构不必靠想象，工具可以直接翻开看。常用的四个用
法：\code{readelf -h} 读头部，核对类型、机器、入口和
两份目录的位置；\code{readelf -S} 列节表，看每个节的
类型、偏移、大小和属性标志；\code{readelf -s} 列符号
表；\code{objdump -d -M intel} 反汇编 \code{.text}，换
\code{-r} 则看重定位记录。拿一个最小源文件完整试一
遍：

\begin{lstlisting}
/* min.c */
int  counter = 42;        /* 初值非零：进 .data    */
const char msg[] = "hi";  /* 只读常量：进 .rodata  */
int  stash;               /* 未初始化：进 .bss     */

int bump(int x)           /* 函数体：进 .text      */
{
    counter += x;         /* 两条指令引用 counter  */
    return counter;
}
\end{lstlisting}

用 \code{gcc -c -O1 -o min.o min.c} 编译，然后逐个窗口
看进去。第一扇窗是头部，注意类型、机器、入口和两份目
录的位置：

\begin{lstlisting}
$ readelf -h min.o
ELF Header:
  Magic:   7f 45 4c 46 02 01 01 00 00 00 00 00 00 00 00 00
  Class:                             ELF64
  Data:                              2's complement, little endian
  Version:                           1 (current)
  OS/ABI:                            UNIX - System V
  Type:                              REL (Relocatable file)
  Machine:                           Advanced Micro Devices X86-64
  Entry point address:               0x0
  Start of program headers:          0 (bytes into file)
  Start of section headers:          472 (bytes into file)
  Size of this header:               64 (bytes)
  Number of program headers:         0
  Number of section headers:         9
  Section header string table index: 8
\end{lstlisting}

魔数之后，\code{02} 表示 64 位、\code{01} 表示小端，与
本机一致；类型 REL 说明这是半成品，所以入口地址是 0、
程序头表个数也是 0——离执行还早，段表还没必要存在。
第二扇窗是节表：

\begin{lstlisting}
$ readelf -S min.o
There are 9 section headers, starting at offset 0x1d8:

Section Headers:
  [Nr] Name       Type     Address          Off    Size   ES Flg Lk Inf Al
  [ 0]            NULL     0000000000000000 000000 000000 00      0   0  0
  [ 1] .text      PROGBITS 0000000000000000 000040 00000d 00  AX  0   0  1
  [ 2] .rela.text RELA     0000000000000000 0001a8 000030 18   I  6   1  8
  [ 3] .data      PROGBITS 0000000000000000 000050 000004 00  WA  0   0  4
  [ 4] .bss       NOBITS   0000000000000000 000054 000004 00  WA  0   0  4
  [ 5] .rodata    PROGBITS 0000000000000000 000054 000003 00   A  0   0  1
  [ 6] .symtab    SYMTAB   0000000000000000 000058 0000f0 18      7   6  8
  [ 7] .strtab    STRTAB   0000000000000000 000148 00001e 00      0   0  1
  [ 8] .shstrtab  STRTAB   0000000000000000 000166 00003f 00      0   0  1
Flags: W (write), A (alloc), X (execute), I (info)
\end{lstlisting}

三处细节值得停下来看。第一，所有节的 Address 一栏全
是 0——地址还不存在，这正是“半成品”最直观的迹象。
第二，\code{.bss} 的文件偏移（0x54）与 \code{.rodata}
相同：NOBITS 不占文件字节，所以“链接：重定位与地址形成”一节从同一偏移接着
排，而它的 4 字节大小照样登记在案，装载时按这个数划
出内存。第三，\code{.text} 只有 13 字节——一个函数
的全部代码就这么多，里面还留着两个待填的洞。第三扇
窗是符号表：

\begin{lstlisting}
$ readelf -s min.o
Symbol table '.symtab' contains 10 entries:
   Num:    Value          Size Type    Bind   Vis      Ndx Name
     0: 0000000000000000     0 NOTYPE  LOCAL  DEFAULT  UND
     1: 0000000000000000     0 FILE    LOCAL  DEFAULT  ABS min.c
     2: 0000000000000000     0 SECTION LOCAL  DEFAULT    1
     3: 0000000000000000     0 SECTION LOCAL  DEFAULT    3
     4: 0000000000000000     0 SECTION LOCAL  DEFAULT    4
     5: 0000000000000000     0 SECTION LOCAL  DEFAULT    5
     6: 0000000000000000    13 FUNC    GLOBAL DEFAULT    1 bump
     7: 0000000000000000     4 OBJECT  GLOBAL DEFAULT    3 counter
     8: 0000000000000000     4 OBJECT  GLOBAL DEFAULT    4 stash
     9: 0000000000000000     3 OBJECT  GLOBAL DEFAULT    5 msg
\end{lstlisting}

源码里的四个名字各就各位：\code{bump} 是 13 字节的
FUNC，定义在 1 号节（\code{.text}）；\code{counter} 是
4 字节 OBJECT，在 3 号节（\code{.data}）；\code{stash}
在 4 号节（\code{.bss}）；\code{msg} 是 3 字节（连结尾
的 \codetext{'\textbackslash 0'}），在 5 号节
（\code{.rodata}）。Value 一栏全是 0：它们目前都只是
节内偏移，三个变量恰好都位于各自节的起点。最后一扇窗
看代码本身和那两个洞：

\begin{lstlisting}
$ objdump -d -M intel min.o
Disassembly of section .text:

0000000000000000 <bump>:
   0:  01 3d 00 00 00 00   add  DWORD PTR [rip+0x0],edi   # 6
   6:  8b 05 00 00 00 00   mov  eax,DWORD PTR [rip+0x0]   # c
   c:  c3                  ret

$ objdump -r min.o
RELOCATION RECORDS FOR [.text]:
OFFSET           TYPE              VALUE
0000000000000002 R_X86_64_PC32     counter-0x4
0000000000000008 R_X86_64_PC32     counter-0x4
\end{lstlisting}

两条指令都带 \code{[rip+0x0]} 形式的 4 字节位移字段，
内容暂时是 0——这就是洞。重定位记录说清了两个洞的位
置（偏移 2 和 8，各在操作码之后）、填法
（R\_X86\_64\_PC32，PC 相对 32 位）和目标（符号
\code{counter}，加数 -4）。注意 \code{counter} 就在本
文件定义，照样要留洞：它在 \code{.data}，代码在
\code{.text}，两节的距离要等布局后才知道，怎么填正是
第二节的主题。

\topic{一个最小函数的完整解剖}
收尾前做一次最彻底的走查：把一个最小函数从源码到字
节、到符号、到重定位记录逐项对一遍。源码只有一行：

\begin{lstlisting}
int add(int a, int b) { return a + b; }
\end{lstlisting}

\code{gcc -O0} 生成的汇编如下（Intel 语法，小写寄存
器）：

\begin{lstlisting}
add:
        push    rbp
        mov     rbp, rsp
        mov     DWORD PTR [rbp-4], edi
        mov     DWORD PTR [rbp-8], esi
        mov     edx, DWORD PTR [rbp-4]
        mov     eax, DWORD PTR [rbp-8]
        add     eax, edx
        pop     rbp
        ret
\end{lstlisting}

这段代码汇编成 20 个字节。逐字节解剖如下，每一条指令
的编码都能用处理器核心相关章节的字段格式读出来：

\begin{longtable}{L{0.07\linewidth}L{0.17\linewidth}L{0.22\linewidth}L{0.4\linewidth}}
\toprule
偏移 & 字节 & 指令 & 编码要点 \\
\midrule
0x00 & \code{55} & \code{push rbp} & 单字节，寄存器号 5 编在操作码低 3 位（0x50+5） \\\tablerowrule
0x01 & \code{48 89 e5} & \code{mov rbp,rsp} & \code{48} 是 REX.W 前缀；ModRM=0xE5：mod=11 寄存器直寻，reg=100 是 rsp，rm=101 是 rbp \\\tablerowrule
0x04 & \code{89 7d fc} & \code{mov [rbp-4],edi} & ModRM=0x7D：mod=01 带 8 位位移，disp8=0xFC 即 -4，reg=111 是 edi \\\tablerowrule
0x07 & \code{89 75 f8} & \code{mov [rbp-8],esi} & 同一格式，reg=110 是 esi，disp8=0xF8 即 -8 \\\tablerowrule
0x0a & \code{8b 55 fc} & \code{mov edx,[rbp-4]} & 操作码 8B 方向相反：内存到寄存器，reg=010 是 edx \\\tablerowrule
0x0d & \code{8b 45 f8} & \code{mov eax,[rbp-8]} & reg=000 是 eax，位移仍是 -8 \\\tablerowrule
0x10 & \code{01 d0} & \code{add eax,edx} & ModRM=0xD0：rm=000（eax）是目的，reg=010（edx）是源 \\\tablerowrule
0x12 & \code{5d} & \code{pop rbp} & 单字节，0x58+5 \\\tablerowrule
0x13 & \code{c3} & \code{ret} & 弹栈顶作为返回地址 \\
\bottomrule
\end{longtable}

每一列都能和处理器核心相关章节对上。ModRM 字节的 mod、reg、rm 三
字段正是处理器核心相关章节“寻址方式加寄存器号”的同一套结构；
\code{[rbp-4]} 里的 -4 就是编码格式中的 disp8 位移字
段。参数与返回值的位置则来自调用约定——处理器核心相关章节“调用
约定把寄存器和栈的责任写成明确约定”的 x86-64 版本：
头两个整型参数依次在 edi、esi，返回值在 eax。还可以
看到 -O0 的字面意思：参数先原样存入栈、用时再取回，
20 个字节里真正做加法的只有 \code{01 d0} 两个字节。
若开 -O1，整个函数收缩成 \code{8d 04 37 c3}：
\code{lea eax,[rdi+rsi]} 借用“基址加变址”的求和路
径直接算出 a+b，外加一条 \code{ret}——地址形成硬件
被借来做算术，这是处理器核心相关章节 SIB 求和器的又一次登场。

再看符号与重定位两栏的对应。符号表里这个函数产生一个
条目：名字 \code{add}，类型 FUNC，绑定 GLOBAL，大小
20，Ndx 指向 \code{.text}，Value 为 0——函数正好位于
本文件 \code{.text} 的起点，最终虚拟地址要等链接器分
配，此刻只能记节内偏移。重定位记录一栏则是空的：这个
函数的每个操作数不是寄存器就是 rbp 相对的栈槽，位移
-4、-8 在编译时就是定数，整条指令流不引用任何名字，
自然没有要填的洞。这给出一条判据：重定位记录只在指令
里出现名字——外部函数、全局变量——时才产生。把函数
改成带外部调用的版本，记录立刻现身：

\begin{lstlisting}
extern int twice(int);

int add_twice(int a, int b) { return twice(a + b); }
\end{lstlisting}

\code{-O1} 生成的代码与字节如下（分号后是注释）：

\begin{lstlisting}
add_twice:
        sub     rsp, 8            ; 48 83 ec 08  栈对齐到 16 字节
        lea     edi, [rdi+rsi]    ; 8d 3c 37     a+b 送入第一个参数
        call    twice             ; e8 00 00 00 00  <- 洞在这里
        add     rsp, 8            ; 48 83 c4 08
        ret                       ; c3
\end{lstlisting}

\code{call} 的 4 字节相对偏移在编译时未知——
\code{twice} 在别的文件里，它的地址此刻不存在。汇编
器把字段清零，并登记一条重定位记录：

\begin{lstlisting}
$ objdump -r add_twice.o
RELOCATION RECORDS FOR [.text]:
OFFSET           TYPE              VALUE
0000000000000008 R_X86_64_PC32     twice-0x4
\end{lstlisting}

读法：\code{.text} 偏移 0x8 处（\code{e8} 操作码之后
那 4 个字节）要按 PC 相对 32 位规则、用符号
\code{twice} 的地址修补，加数 -4。三个要素——改哪
里、按什么规则、用谁的值——全部登记在案，但演算还做
不了：S 尚未存在。填洞需要符号的值，而符号的值来自链
接器的布局，这就是第二节的开场。

\section{二、链接：重定位与地址形成}

“目标文件与符号”一节的结论是：目标文件是半成品——字节齐备，地址未
定，凡涉及名字的字段都留着洞，洞的位置和填法记在重定
位记录里。本节看链接器怎样把若干半成品拼成一张完整的
地址地图。

\topic{链接器做三件事}
把链接器看成一台只做三件简单事的机器，它的行为就不再
神秘。第一件，合并同类节：把各输入文件的 \code{.text}
依次接成输出的 \code{.text}，\code{.data}、
\code{.rodata}、\code{.bss} 同样处理，再按读写执行属
性把节归并成段。第二件，符号解析与地址分配：给输出的
每个节定下虚拟地址，于是每个符号的值从“某节内偏移”
换算成“节基址加偏移”的具体数；同时把各文件的符号表
汇成一张全局表，为每个 UND 引用找到唯一的 GLOBAL 定
义。第三件，重定位：按重定位记录逐条改写输出文件里的
字节——填 \code{call} 的相对偏移、填全局变量的位移、
填数据区的指针初值。

\begin{longtable}{L{0.22\linewidth}L{0.24\linewidth}L{0.26\linewidth}L{0.14\linewidth}}
\toprule
步骤 & 输入 & 产出 & 改文件字节吗 \\
\midrule
合并同类节 & 各 \code{.o} 的同名节 & 输出的节与段布局 & 拼接拷贝 \\\tablerowrule
符号解析与定址 & 全部符号表 & 全局符号表、各节基址 & 不改，只算数 \\\tablerowrule
重定位 & 重定位记录加符号值 & 填好的指令与数据 & 改，逐条修补 \\
\bottomrule
\end{longtable}

三件事里前两件是簿记，第三件才真正动字节；但改哪里、
改多少，完全由前两件的结果决定。这就是“链接等于把多
个半成品拼成一张完整地址地图”的准确含义：地图由第二
件事算出，由第三件事誊进字节。次序也不能乱：只有合并
与定址完成，符号才有值；符号有了值，重定位的算术才能
做。所以第一节那两个经典报错都发生在第二件事——名字
查不到或重复定义，地图就画不出来，根本轮不到第三件事
动手。

把三件事放到一对最小输入上看一遍数值。设链接器只收到
\code{main.o} 与 \code{add.o}：合并把 \code{main.o} 的
\code{.text}（0x20 字节）铺在输出 \code{.text} 的基址
\code{0x401120} 处，\code{add.o} 的 \code{.text}（20
字节，即 0x14）紧随其后，从 \code{0x401140} 开始；数
据节同样依次排布、各按对齐要求取整。定址随即给出符号
值：\code{main = 0x401120}，\code{add = 0x401140}，
\code{main.o} 符号表里那条 UND 的 \code{add} 就此有了
定义；重定位拿着这两个值去填 \code{main.o} 里那条
\code{call} 留下的洞——具体演算正是下一个专题的内
容。

把合并与定址画成一张排布图，
“同类节依次拼接、符号值随
之确定”就是图上的两个动作：

\begin{pdfdiagram}
  % 输入：两个目标文件的同类节
  \node[hw state, minimum width=2.4cm, align=center] (mt) at (1.2,2.4) {main.o\\.text 0x20 字节};
  \node[hw state, minimum width=2.4cm, align=center] (at) at (1.2,1.2) {add.o\\.text 0x14 字节};
  \node[hw memory, minimum width=2.4cm, align=center] (md) at (1.2,-0.2) {main.o\\.data};
  \node[hw memory, minimum width=2.4cm, align=center] (ad) at (1.2,-1.4) {add.o\\.data};
  % 输出：合并后的节，地址已定
  \node[hw compute, minimum width=4.2cm, align=center] (ot) at (8.8,1.8) {输出 .text\\main @ 0x401120\\add @ 0x401140};
  \node[hw compute, minimum width=4.2cm, align=center] (od) at (8.8,-0.8) {输出 .data\\初值依次排布\\按对齐取整};
  % 拼接箭头：各自折入输出框
  \draw[hw bus] (mt.east) -- (3.4,2.4) -- (3.4,2.05) -- ($(ot.west)+(0,0.25)$);
  \draw[hw bus] (at.east) -- (3.0,1.2) -- (3.0,1.55) -- ($(ot.west)+(0,-0.25)$);
  \node[hw label, fill=white, inner sep=1pt] at (4.9,2.05) {从基址排起};
  \node[hw label, fill=white, inner sep=1pt] at (4.9,1.55) {紧随其后};
  \draw[hw bus] (md.east) -- (3.4,-0.2) -- (3.4,-0.55) -- ($(od.west)+(0,0.25)$);
  \draw[hw bus] (ad.east) -- (3.0,-1.4) -- (3.0,-1.05) -- ($(od.west)+(0,-0.25)$);
  % 定址结果
  \node[hw label, align=center] at (5.6,-2.4) {定址结果：main = 0x401120，add = 0x401140\\0x401140 = 0x401120 + 0x20，拼接次序即地址次序};
\end{pdfdiagram}
\diagramnote{静态链接的地址形成（对应正文 main.o 与 add.o 的最小例）。链接器把各输入文件的同类节依次拼接成输出的节：\code{main.o} 的 \code{.text}（0x20 字节）铺在基址 \code{0x401120}，\code{add.o} 的 \code{.text}（0x14 字节）紧随其后从 \code{0x401140} 开始；数据节同样处理。拼接完成，每个符号的值便从“节内偏移”换算成“节基址加偏移”的具体数——\code{main = 0x401120}，\code{add = 0x401140}，重定位所需的 S 至此全部就位。}

工程上，布局不是只能听默认值。链接脚本可以指定每个节
的位置、入口符号和段的对齐方式：本书末章的教学计算机项目的裸机镜像把
\code{.text} 固定在 ROM 起始地址、把入口固定在复位向量，
用的就是这套机制的手工模式；Linux 的默认链接脚本则把
\code{.text} 安排在 \code{0x400000} 附近的惯例地址。
同一台机器，自动挡与手动挡并存；学完本节和第三节，可
以回头把本书末章的教学计算机项目那个链接脚本逐行再读一遍。

\topic{重定位记录怎么读}
第一节末尾留下一条记录：\code{offset 0x8, type
R\_X86\_64\_PC32, symbol twice, addend -4}。现在把它读
透。一条重定位记录回答三个问题：位置（offset）——改
输出文件里的哪几个字节；类型（type）——按哪条公式
算；符号与加数（symbol、addend）——用谁的地址、再做
什么修正。x86-64 的 ELF 用 RELA 格式，加数显式存在记
录里，被修补字段的旧内容直接丢弃不看。

顺带一提，32 位 x86 的 ELF 用的是 REL 格式：记录里没
有加数字段，加数藏在被修补字段的旧内容里，链接器要先
读出旧值再代入计算。RELA 把加数显式化，字段旧内容一
律作废，逻辑更直白。这也解释了观察重定位时的一个现
象：在 RELA 文件里，洞里的旧字节是什么都无所谓，汇编
器统一清零——“目标文件与符号”一节目标文件里那些
\code{00 00 00 00} 的位移字段，正是这个约定的样子。

最常用的两种类型对应两条公式。绝对 32 位重定位
（R\_X86\_64\_32）：往字段写入 S+A 的低 32 位，S 是符
号的最终地址，A 是加数——填的是地址本身。PC 相对 32
位重定位（R\_X86\_64\_PC32）：写入 S+A−P，P 是被修补
字段自己的最终地址——填的是字段到目标的距离。下面把
一条 \code{call} 的修补从头算到尾。设 \code{main.o}
里有一处调用 \code{add}（定义在另一个文件），用
\code{objdump -dr} 可以把指令和记录放在一起看：

\begin{lstlisting}
$ objdump -dr -M intel main.o
Disassembly of section .text:

0000000000000000 <main>:
   ...
  10:  e8 00 00 00 00      call   15 <main+0x15>
                         11: R_X86_64_PC32  add-0x4
   ...
\end{lstlisting}

链接器做完合并与定址后：\code{main} 所在 \code{.text}
块的基址定为 \code{0x401120}，\code{add} 定在
\code{0x401140}。于是三个数全部就位：被修补字段的地
址 P = 0x401120 + 0x11 = 0x401131；符号值 S =
0x401140；加数 A = −4。代入公式，并站在 CPU 执行语义
一侧核对结果：

\begin{lstlisting}
写入值 = S + A - P
       = 0x401140 + (-4) - 0x401131
       = 0x401140 - 0x401135
       = 0x0000000B

核对（CPU 执行 call 的语义）：
目标 = 下一条指令地址 + rel32
     = (0x401131 + 4) + 0xB
     = 0x401135 + 0xB
     = 0x401140  = add    对上了
\end{lstlisting}

修补后，该处字节从 \code{e8 00 00 00 00} 变成
\code{e8 0b 00 00 00}——0xB 按小端排布，低字节在前，
这正是链接器必须通晓 ISA 字节序与字段位置的原因。加
数为什么是 −4 也值得说透：CPU 求 \code{call} 目标时
的基准是下一条指令的地址，即字段末尾之后；公式里的 P
却是字段的开头，两者恰好差 4 个字节，汇编器把这 4 字
节折进加数里补齐。整个演算没有一个新数：S 来自符号
表，P 来自节布局，A 来自记录本身，公式来自 ISA 对
\code{call} 语义的定义。四者各就各位，填洞就是一次减
法。

把这条演算画成修补流水线，四个数各就各位的动作就是四步：

\begin{pdfdiagram}
  % 左列：三个数的来源
  \node[hw memory, align=center, minimum width=3.0cm] (rec) at (0,1.6) {重定位记录\\offset 0x11，addend $-4$\\类型 PC32};
  \node[hw state, align=center, minimum width=3.0cm] (sym) at (0,0) {符号表\\S = add = 0x401140};
  \node[hw state, align=center, minimum width=3.0cm] (lay) at (0,-1.6) {节布局\\P = 0x401131};
  % 汇入干线
  \draw[draw=BookMuted, line width=0.62pt] (rec.east) -- (2.4,1.6);
  \draw[draw=BookMuted, line width=0.62pt] (sym.east) -- (2.4,0);
  \draw[draw=BookMuted, line width=0.62pt] (lay.east) -- (2.4,-1.6);
  \draw[draw=BookMuted, line width=0.62pt] (2.4,1.6) -- (2.4,-1.6);
  \fill[BookMuted] (2.4,1.6) circle (0.05);
  \fill[BookMuted] (2.4,0) circle (0.05);
  \fill[BookMuted] (2.4,-1.6) circle (0.05);
  % 计算与写回
  \node[hw compute, align=center, minimum width=3.6cm] (calc) at (5.6,0) {S + A $-$ P\\= 0x401140 $-$ 4 $-$ 0x401131\\= 0x0B};
  \node[hw memory, align=center, minimum width=3.0cm] (bytes) at (9.9,0) {call 字段\\e8 00 00 00 00\\→ e8 0b 00 00 00};
  \draw[hw bus] (2.4,0) -- node[midway, above=1pt, hw label, fill=white, inner sep=1pt]{三数就位} (calc.west);
  \draw[hw bus] (calc.east) -- node[midway, above=1pt, hw label, fill=white, inner sep=1pt]{写回} (bytes.west);
\end{pdfdiagram}
\diagramnote{重定位修补流水线（对应正文 main.o 调用 add 的演算）。链接器读出记录的位置、类型与加数，从符号表取 S、从节布局取 P，按类型公式代入：PC32 写入 S+A$-$P = 0x401140 $-$ 4 $-$ 0x401131 = 0x0B；结果按小端排布写回字段，e8 00 00 00 00 变成 e8 0b 00 00 00。RELA 格式下字段旧内容直接丢弃，整个修补就是一次减法加一次写字节。}

\topic{绝对重定位与 PC 相对重定位}
两种类型为什么并存？因为它们服务的字段性质不同。PC 相
对重定位填距离，适合代码内部的控制转移和就近访问的数
据：\code{call}、\code{jmp}、条件转移，以及 x86-64 的
RIP 相对数据寻址，字段 32 位，够到前后各 2 GiB。绝对
重定位填地址本身，适合数据区里的指针：函数指针数组、
\code{switch} 的跳转表、\code{char *p = msg;} 这类全
局指针的初值，以及 \code{movabs} 的 64 位立即数字段
——这些场合程序要的是值，不是距离。

差别要到镜像搬动时才显形。绝对字段把装载地址烧进了字
节：镜像每换一个装载位置，所有绝对字段都得按新基址重
填一遍——这正是第三节要讲的装载时重定位，也是 ASLR
每次随机布局都必须付的代价。PC 相对字段只记距离：整
段代码连同它引用的目标一起搬，距离不变，字节原样可
用，一处都不用改。

看一个绝对重定位的具体形状。\code{switch} 语句分支密
集时，编译器常在 \code{.rodata} 里生成跳转表：每个表
项是 8 字节的绝对地址，指向某个分支的代码。目标文件
里这些表项先是零，旁边一一对应 R\_X86\_64\_64 记录；
链接时按公式 S+A 填入各分支标签的最终地址。一旦镜像
换了装载基址，整张表每一项都要加上同一个差值——表
有多大，重填的工作量就有多大，绝对字段的代价在此现
出原形。

条件转移则是 PC 相对的纯血用户：x86-64 的 \code{jz}、
\code{jne} 只有 rel8 与 rel32 两种偏移形式，编码里根
本没有放绝对目标的位置，想跳远只能先跳到附近再做间
接转移。控制流指令因此天生位置无关，代码段整体搬移
时，从来不必为它们操心。

这就是位置无关代码（PIC）的雏形。共享库想被多个进程
映射到各自不同的虚拟地址、又共享同一批物理代码页，前
提是 \code{.text} 里几乎没有绝对字段：代码内部的转移
天然是 PC 相对的，剩下绕不开的绝对访问——调外部函
数、取外部数据——全部改道两张间接表（GOT 与 PLT），
需要按装载地址修补的绝对字段集中收进这两张表里，
\code{.text} 保持干净。编译共享库时加 \code{-fPIC} 的
本质，就是把绝对重定位从代码节里赶出去；两张表怎么工
作，是第三节的内容。

\begin{longtable}{L{0.2\linewidth}L{0.26\linewidth}L{0.24\linewidth}L{0.16\linewidth}}
\toprule
类型 & 写入内容 & 典型场合 & 镜像整体搬动后 \\
\midrule
R\_X86\_64\_32 / 64 & S+A，地址本身 & 数据指针、跳转表、\code{movabs} & 所有绝对字段要重填 \\\tablerowrule
R\_X86\_64\_PC32 & S+A−P，与字段的距离 & \code{call}/\code{jmp}、RIP 相对访存 & 字节原样可用 \\
\bottomrule
\end{longtable}

回头看历史能加深理解：x86-64 把 RIP 相对寻址设为默认
的数据寻址方式，正是操作系统全面转向 PIC 之后 ISA 做
出的配合。处理器核心相关章节讲寻址方式时给过这一模式的硬件结构，
现在能看到它另一半成因在软件侧：链接器和装载器喜欢距
离，讨厌地址本身。ISA 与系统软件在这里是互相塑造的。

\topic{静态库是按需求出的目标文件包}
静态库 \code{libfoo.a} 不是新的文件格式：用 \code{ar}
把一批 \code{.o} 捆成一个归档，再附一张符号索引
（\code{ar} 的 \code{s} 选项，旧系统里由 \code{ranlib}
生成），索引记录哪个名字在哪个成员里，仅此而已。库成
员仍是普普通通的可重定位文件，链接器对它们的处理也和
对普通 \code{.o} 一样——只多一条抽取规则。

抽取规则是理解静态库的全部关键。链接器从左到右扫描命
令行上的输入，沿途维护一个“当前未定义符号”集合：扫
到普通 \code{.o}，无条件整个收进，它定义的符号进全局
表，它引用的符号进未定义集合；扫到库，则翻看索引，只
把能消解当前未定义符号的成员抽出来——被抽出的成员自
己又可能带来新引用，于是对同一个库要反复查，直到不再
有新成员被抽出。扫描结束时，未定义集合必须为空，否则
报错。

把这条抽取规则画成数据流，“用多少、取多少”就是中间那个集合的进出平衡：

\begin{pdfdiagram}
  % 左列：两类输入
  \node[hw state, align=center] (maino) at (0,1.5) {main.o\\普通目标文件};
  \node[hw memory, align=center, minimum width=3.2cm] (lib) at (0,-1.1) {libfoo.a\\ar 归档＋符号索引\\成员 a.o b.o c.o};
  % 中列：未定义集合与被抽出的成员
  \node[hw control, align=center, minimum width=2.8cm] (undef) at (5.6,1.5) {当前未定义\\符号集合};
  \node[hw compute, align=center, minimum width=2.8cm] (ext) at (5.6,-1.1) {按索引抽出的\\成员};
  % 右列：输出
  \node[hw state, align=center] (out) at (10.3,0.2) {可执行文件};
  \draw[hw bus] (maino.east) -- node[midway, above=1pt, hw label, fill=white, inner sep=1pt]{无条件整个收进} (undef.west);
  \draw[hw bus] (lib.east) -- node[midway, above=1pt, hw label, fill=white, inner sep=1pt]{按需抽取} (ext.west);
  \draw[hw bus] (ext.east) -- (8.1,-1.1) -- (8.1,0.2) -- (out.west);
  % 集合与抽取之间的控制往返
  \draw[hw return] ($(undef.south)+(-0.6,0)$) -- node[midway, left=1pt, hw label, fill=white, inner sep=1pt]{查索引} ($(ext.north)+(-0.6,0)$);
  \draw[hw return] ($(ext.north)+(0.6,0)$) -- node[midway, right=1pt, hw label, fill=white, inner sep=1pt]{新引用回集合} ($(undef.south)+(0.6,0)$);
  \node[hw label, align=center] at (5.6,2.7) {扫描结束：集合必须为空};
\end{pdfdiagram}
\diagramnote{静态库的按需抽取。链接器从左到右扫描输入：普通目标文件无条件整个收进，它引用的符号进未定义集合；扫到库时翻符号索引，只把能消解当前未定义符号的成员抽出；被抽成员若带来新引用，就对同一个库再查一遍，直到没有新成员。扫描结束时集合必须为空，否则报未定义引用——库因此必须排在命令行上使用它的文件之后。}

顺序因此有了意义：库必须出现在使用它的文件之后。
\code{gcc main.o -lfoo} 里，扫到库时 \code{main.o} 的
引用已在未定义集合里等着，成员被正常抽出；
\code{gcc -lfoo main.o} 里，扫到库时集合还是空的，一
个成员也不抽；等 \code{main.o} 把引用放进来，库已经
扫过，链接器不会回头。

连锁抽取值得再看一眼：成员甲被抽出后又引用了成员乙定
义的符号，乙也得跟着出库，所以库内部的依赖会被引用链
自动拉齐。真正的麻烦在库与库互相引用：\code{liba.a}
需要 \code{libb.a}，\code{libb.a} 又回头需要
\code{liba.a}，单向扫描无论谁先谁后都可能漏抽。
\code{ld} 为此准备了 \code{--start-group} 与
\code{--end-group}：把互相循环的库圈进一组，链接器在
组内反复扫描，直到没有新成员被抽出。日常遇到的顺序问
题多半不是循环依赖，而是单纯把库摆错了位置。

\begin{lstlisting}
$ gcc -c main.c helper.c
$ ar rcs libhelper.a helper.o
$ gcc main.o -L. -lhelper -o app        # 正确：库在使用者之后
$ gcc -L. -lhelper main.o -o app        # 经典错误：库扫描早于未解析符号出现
/usr/bin/ld: main.o: in function `main':
main.c:(.text+0x1f): undefined reference to `helper'
collect2: error: ld returned 1 exit status
\end{lstlisting}

第一节那两个经典错误，在这里能看到完整成因。未定义引
用：扫描结束而未定义集合非空——可能忘了给库，可能顺
序错了，也可能库里确实没有这个名字；三重原因共用一个
症状，排查时就按这三条挨个问。多重定义：两个输入各交
出一个同名强定义，链接器无法裁决。边界情形也值得一
记：两个弱定义不报错，链接器任取一个；一强一弱取强
——这条规则正是库用弱符号提供可替换默认实现的依据。

“用多少、取多少”的抽取粒度带来两个直接的工程后果。
其一，只引用库里一个函数就只抽一个成员，所以
\code{libc.a} 这类库把每个函数单独成目标文件（成员），静态链接出的
可执行文件才不会把整个库背在身上。其二，抽取沿引用链
传递：抽出的成员又引用别的成员会连锁抽取，所以被依赖
的底层库（如 \code{-lm}）要放在命令行更靠后的位置。
第三节会看到，动态库换了一套哲学：整个库先映射进内存
再说，抽取问题让位给绑定问题。

\topic{与处理器核心相关章节 ISA 寻址方式的关系}
现在把镜头拉回处理器核心相关章节。那里讲指令编码时反复出现两类字
段：立即数字段和地址/偏移字段——教学 ISA 里
\code{JZ +3}、\code{JMP -5} 的偏移，x86-64 里
\code{call} 的 rel32、RIP 相对寻址的 disp32、
\code{movabs} 的 imm64。当时这些字段的值都由汇编器顺
手算出；本章能看到它们的另一半来历：凡字段里出现名
字，值就要等链接器。链接器填的正是处理器核心相关章节编码格式里那
些 imm 与地址字段，一分不多，一分不少。

两处约定遥相呼应。处理器核心相关章节的教学 ISA 规定“跳转目标等
于本条地址加 1 加偏移”，x86-64 规定“目标等于下一条
指令地址加 rel32”——同一个“相对下一条指令”的约
定，一个定长一个变长而已。本节的公式 S+A−P 就是这个
约定的代数形式：式中每个量都对应编码格式里一个具体的
字段位置。学会读重定位记录，等于把处理器核心相关章节的指令编码又
读了一遍，只是这次手里拿的是链接器的视角。

还有一个细节能巩固这层对应。处理器核心相关章节给教学 ISA 定编码
时，把立即数与偏移一律向低位对齐，理由是扩展逻辑单
一；链接器面对的正是同一批低位字段——rel32 从指令的
固定位开始占 4 字节，链接器只改写这 4 字节，绝不碰操
作码与 ModRM。字段的位置由 ISA 定死，字段的值由链接
器算定，双方各守一半约定，拼起来才是一条完整的指令。

更完整地说，地址形成是一条跨越四个时刻的流水线，每个
时刻各填各的字段：

\begin{longtable}{L{0.1\linewidth}L{0.24\linewidth}L{0.52\linewidth}}
\toprule
时刻 & 谁定值 & 定什么 \\
\midrule
编译时 & 编译器/汇编器 & 字段的位置与格式、寄存器与寻址方式、重定位记录 \\\tablerowrule
链接时 & 链接器 & 静态符号的地址：\code{call} 偏移、全局量位移、数据指针初值 \\\tablerowrule
装载时 & 装载器/动态链接器 & 运行前才定的地址：库基址、GOT 条目（第三节） \\\tablerowrule
运行时 & CPU 硬件 & 有效地址求和：基址+变址×比例+位移 \\
\bottomrule
\end{longtable}

链接器坐在这条流水线的倒数第二站：它不管语法语义（那
是编译前端的事），也不生成指令模板（那是编译后端的
事），只管“名字到地址”的最后定值。说链接是 ISA 地
址形成的最后一步，准确含义就在于此——处理器核心相关章节数据通路
里那个“基址加变址乘比例加位移”的求和器，它的位移输
入常常就是链接器在磁盘上写下的那个数。裸机世界同样成
立：本书末章的教学计算机项目的镜像没有操作系统，链接脚本把节固定在物理地
址上，重定位照算不误，S 与 P 只是换成了手工指定的
值。从教学板到 Linux，这套机制是同一件事的两种装法。

读懂这一层，链接错误就不再神秘：它是地址形成流水线倒
数第二站的报错——某个名字走到这一站还没有地址，或者
有两个地址。“装载与动态链接”一节进入最后一站：装载器把这张填好的地
图铺进内存。

\section{三、装载与动态链接}

“链接：重定位与地址形成”一节结束时，磁盘上已经躺着一个地址全部定好的可执行文件：代码在 \code{0x401000} 附近，数据在 \code{0x403dd8} 附近，入口点写在文件头里。但它此刻还只是一串字节，离“正在运行的进程”还差最后一步，这一步由内核完成：\code{execve} 把文件变成进程，动态链接器把库接进来。本节回答两件事：内核装载时做了什么，以及为什么其中相当大的一部分被推迟到“第一次用到”的那一刻。

先立一个总观点：装载不是搬运。直觉上“把程序装进内存”像一次大块复制——把文件读进 RAM，再跳过去执行；本书末章的教学计算机项目的面包板机把程序逐字节拨进 RAM，就是这种字面意义的装载。现代内核的做法不同：它只登记映射——“这段虚拟地址对应那个文件的这段字节”——真正的字节要等程序第一次访问那一页时，才由缺页异常从磁盘取回。把\hwkey{装载理解成建立映射而不是复制字节}，按需调页、写时复制、共享库就都是同一条逻辑的推论。

\topic{execve 之后内核做的事}
\code{execve} 是少数“成功后不返回”的系统调用：成功时原进程的地址空间被整个换掉，等它“返回”到用户态时，运行的已经是另一个程序。内核按固定顺序做四件事：读 ELF 头部与段表、按段表建立映射、准备栈、跳到入口。下面把它做完一遍。

先把四步按执行先后画成一条流水线，“只登记、不搬字节”写在第二步上：

\begin{pdfdiagram}
  \node[hw memory, minimum width=5.2cm, align=center] (src) at (2.8,3.3) {磁盘上的 ELF 文件\\头部、段表、各节字节};
  \node[hw compute, minimum width=5.2cm, align=center] (s1) at (2.8,2.1) {① 读头部：核对魔数与类型\\取入口 0x401040 与段表};
  \node[hw compute, minimum width=5.2cm, align=center] (s2) at (2.8,0.8) {② 建映射：按 LOAD 段折算\\登记 vm\_area，不读字节};
  \node[hw compute, minimum width=5.2cm, align=center] (s3) at (2.8,-0.5) {③ 备栈：字符串、argv/envp\\auxv，rsp 指向 argc};
  \node[hw control, minimum width=5.2cm, align=center] (s4) at (2.8,-1.8) {④ 跳入口：rip ← 入口\\动态程序先跳 ld.so};
  \draw[hw bus] (src.south) -- (s1.north);
  \draw[hw bus] (s1.south) -- (s2.north);
  \draw[hw bus] (s2.south) -- (s3.north);
  \draw[hw bus] (s3.south) -- (s4.north);
  % 右侧：装载产物
  \node[hw state, minimum width=3.0cm, align=center] (proc) at (9.1,-0.5) {新进程\\地址空间备好\\等待首次缺页};
  \draw[hw bus] (s4.east) -- (6.9,-1.8) -- (6.9,-0.5) -- (proc.west);
\end{pdfdiagram}
\diagramnote{\code{execve} 的四步流程（数值对应正文 \code{hello} 例）。内核核对魔数与类型后取出入口 \code{0x401040} 与段表，把每个 \code{LOAD} 段折算成一条映射登记下来——此时程序的一个字节都还没进内存；再把参数字符串、\code{argv}/\code{envp} 与辅助向量摆上栈顶，最后设置 \code{rip}：静态程序直指入口，动态程序先给动态链接器。四步几乎不搬数据，所以 \code{execve} 很快；失败时正常返回 \code{-1}，四步一件都没做。}

第一件事，读头部。内核核对 ELF 魔数与机器类型，确认文件类型（\code{EXEC} 或 \code{DYN}），取出入口地址和程序头表（段表）。段表里的每个 \code{LOAD} 段回答一个问题：文件的哪一段字节，要以什么权限，出现在虚拟地址的哪一段。用 \code{readelf} 看一个典型小程序（节选 \code{LOAD} 段；数值为模拟，格式与真实输出一致）：

\begin{lstlisting}
$ readelf -l hello
Elf file type is EXEC (Executable file)
Entry point 0x401040
There are 11 program headers, starting at offset 64

Program Headers:
  Type           Offset             VirtAddr           PhysAddr
                 FileSiz            MemSiz             Flags  Align
  LOAD           0x0000000000000000 0x0000000000400000 0x0000000000400000
                 0x00000000000004f0 0x00000000000004f0 R      0x1000
  LOAD           0x0000000000001000 0x0000000000401000 0x0000000000401000
                 0x0000000000000195 0x0000000000000195 R E    0x1000
  LOAD           0x0000000000002000 0x0000000000402000 0x0000000000402000
                 0x0000000000000110 0x0000000000000110 R      0x1000
  LOAD           0x0000000000002dd8 0x0000000000403dd8 0x0000000000403dd8
                 0x0000000000000248 0x0000000000000250 RW     0x1000
\end{lstlisting}

第二件事，按段表建映射。把四个 \code{LOAD} 段逐段算账。第一段从文件偏移 0 映到 \code{0x400000}，ELF 头和程序头表都在其中。第二段是代码：文件区间 \code{[0x1000, 0x1195)} 映到 \code{[0x401000, 0x401195)}，\code{0x195} 即 405 字节，不足一页；内核按整页映射，页内超出段长的字节来自文件同页的对齐填充。入口 \code{0x401040} 就在这段内偏移 \code{0x40} 处。第三段是只读数据。第四段是可写数据：\code{MemSiz} 比 \code{FileSiz} 多出 8 字节，多出的部分就是 \code{.bss}——文件里不占字节，内存里必须是零；内核映完文件部分后把这 8 字节清零，若 \code{.bss} 大到跨页，整页挂匿名零页。

\begin{longtable}{L{0.08\linewidth}L{0.26\linewidth}L{0.24\linewidth}L{0.32\linewidth}}
\toprule
段 & 文件偏移→虚拟地址 & 字节数（文件→内存） & 权限与内容 \\
\midrule
1 & \code{0x0000}→\code{0x400000} & \code{0x4f0}→\code{0x4f0} & 只读：ELF 头、段表、\code{.interp} \\\tablerowrule
2 & \code{0x1000}→\code{0x401000} & \code{0x195}→\code{0x195} & 读/执行：\code{.text}，入口在其中 \\\tablerowrule
3 & \code{0x2000}→\code{0x402000} & \code{0x110}→\code{0x110} & 只读：\code{.rodata} \\\tablerowrule
4 & \code{0x2dd8}→\code{0x403dd8} & \code{0x248}→\code{0x250} & 读写：\code{.data} 与 \code{.bss}，多出 8 字节清零 \\
\bottomrule
\end{longtable}

每个 \code{LOAD} 段还满足一条同余：虚拟地址与文件偏移模页长同余——\code{0x403dd8} 与 \code{0x2dd8} 模 \code{0x1000} 都余 \code{0xdd8}。这不是巧合：只有页内偏移一致，一个文件页才能原样映到一个虚拟页上，内核不必做任何搬移对齐。“链接：重定位与地址形成”一节的链接器在定址时就按这条约束摆段：链接不仅决定“放在哪个地址”，还决定“按页怎样对齐”。

映射在内核里的形态是一张映射表（\code{vm\_area\_struct} 的集合），每条记录由段表折算而来。注意这张表是登记簿而非内容：此时程序的一个字节都还没进内存。

\begin{longtable}{L{0.20\linewidth}L{0.38\linewidth}L{0.32\linewidth}}
\toprule
字段 & 含义 & 本例第四段的值 \\
\midrule
起止地址 & 虚拟地址区间 & \code{[0x403dd8, 0x404028)} \\\tablerowrule
权限 & 读/写/执行 & 读、写 \\\tablerowrule
来源 & 文件加偏移，或匿名 & \code{hello} 偏移 \code{0x2dd8} \\\tablerowrule
补齐规则 & 段尾不足一页如何处理 & 超出 \code{FileSiz} 的部分清零 \\
\bottomrule
\end{longtable}

第三件事，准备栈。内核把命令行参数和环境变量的字符串复制到新地址空间的栈顶，再向下摆辅助向量、\code{envp} 和 \code{argv} 指针数组，最后把 \code{rsp} 指向 \code{argc}：

\begin{lstlisting}
high  strings: "./hello\0" "world\0" "PATH=/usr/bin\0" ...
      alignment padding; rsp ends 16-byte aligned
      AT_NULL
      auxv pairs: AT_PAGESZ=4096, AT_ENTRY=0x401040,
                  AT_PHDR=0x400040, ...
      envp: NULL, ptr -> envp strings ...
      argv: NULL, ptr -> "world", ptr -> "./hello"
low   rsp -> argc = 2
\end{lstlisting}

辅助向量（auxiliary vector）是内核留给用户态的小抄：\code{AT\_PAGESZ} 告诉 libc 页长是 4096；\code{AT\_PHDR} 告诉动态链接器段表映在哪里——\code{0x400040} 的来历很直接，ELF 头长 64 字节（\code{0x40}），程序头表紧跟其后，第一段又从 \code{0x400000} 开始映射，所以表在 \code{0x400040}；\code{AT\_ENTRY} 就是文件头里的入口 \code{0x401040}。\code{argc}/\code{argv}/\code{envp} 的摆法则与处理器核心相关章节的栈约定一脉相承：指针数组以 \code{NULL} 结尾，字符串本体在栈顶，\code{rsp} 保持 16 字节对齐。

第四件事，跳到入口。静态链接程序的 \code{rip} 直接设为入口地址；动态链接程序则先映射 \code{.interp} 段指定的动态链接器 \code{/lib64/ld-linux-x86-64.so.2}，\code{rip} 指向它的入口，由它装载共享库、做完装载时重定位，最后才跳到 \code{0x401040}——本节第四、五个主题就发生在它身上。

装载完成后，这套映射从用户态可以直接核对：\code{/proc} 把它原样摆出来。逐项对照段表，能看到每个 \code{LOAD} 段变成一行，以及段表里没有的 libc 和栈：

\begin{lstlisting}
$ cat /proc/742/maps
00400000-00401000 r--p 00000000 08:01 131073    /home/user/hello
00401000-00402000 r-xp 00001000 08:01 131073    /home/user/hello
00402000-00403000 r--p 00002000 08:01 131073    /home/user/hello
00403000-00404000 r--p 00002000 08:01 131073    /home/user/hello
00404000-00405000 rw-p 00003000 08:01 131073    /home/user/hello
7f9a3c090000-7f9a3c22c000 r-xp 00000000 08:01 262146   /lib/x86_64-linux-gnu/libc.so.6
7ffffffde000-7ffffffff000 rw-p 00000000 00:00 0        [stack]
\end{lstlisting}

第四段的兑现方式值得多看一眼：它跨两页，第一页 \code{0x403000} 映的是文件页 \code{0x2000}——同一个文件页也映在第三段的 \code{0x402000}，两处只读映射共享同一个物理页；这页里混着 \code{.rodata} 的尾巴和 \code{.data} 的头——链接器把这段重定位之后不再改写的数据按 RELRO 归进只读段，映射是纯只读的：写入不会触发写时复制，只会吃到一个 SIGSEGV。需要运行期改写的 GOT 条目放在 \code{0x404000}，那页才是可写页。映射表把“一页一个权限、一页一个来源”表达得清清楚楚。

把这套映射从低地址到高地址画成一张布局图，段表到页表的折算链放在右侧：

\begin{pdfdiagram}
  % 左列：地址空间布局（上高下低）
  \node[hw label] at (1.5,5.7) {高地址};
  \node[hw memory, minimum width=3.4cm] (stk) at (1.5,5.0) {栈 [stack]，向下增长};
  \node[hw memory, minimum width=3.4cm] (lib) at (1.5,3.8) {共享库映射（libc）};
  \node[hw label] at (1.5,3.2) {（未映射空洞）};
  \node[hw memory, minimum width=3.4cm] (heap) at (1.5,2.6) {堆，向上增长（brk）};
  \node[hw compute, minimum width=3.4cm] (dseg) at (1.5,1.5) {.data/.bss：段 4，读写};
  \node[hw compute, minimum width=3.4cm] (rseg) at (1.5,0.5) {.rodata：段 3，只读};
  \node[hw compute, minimum width=3.4cm] (tseg) at (1.5,-0.5) {.text：段 2，读/执行};
  \node[hw state, minimum width=3.4cm] (hseg) at (1.5,-1.5) {ELF 头/段表：段 1，只读};
  \node[hw label] at (1.5,-2.2) {低地址 \code{0x400000}};
  % 右列：段表 LOAD 段 → 映射条目 → 页表项
  \node[hw state, minimum width=1.9cm] (l4) at (5.6,1.5) {LOAD 4：RW};
  \node[hw state, minimum width=1.9cm] (l3) at (5.6,0.5) {LOAD 3：R};
  \node[hw state, minimum width=1.9cm] (l2) at (5.6,-0.5) {LOAD 2：R+E};
  \node[hw state, minimum width=1.9cm] (l1) at (5.6,-1.5) {LOAD 1：R};
  \draw[hw bus] (l4.west) -- (dseg.east);
  \draw[hw bus] (l3.west) -- (rseg.east);
  \draw[hw bus] (l2.west) -- (tseg.east);
  \draw[hw bus] (l1.west) -- (hseg.east);
  % 共享竖直干线汇入映射条目
  \draw[BookMuted, line width=0.62pt] (l1.east) -- (6.85,-1.5);
  \draw[BookMuted, line width=0.62pt] (l2.east) -- (6.85,-0.5);
  \draw[BookMuted, line width=0.62pt] (l3.east) -- (6.85,0.5);
  \draw[BookMuted, line width=0.62pt] (l4.east) -- (6.85,1.5);
  \draw[BookMuted, line width=0.62pt] (6.85,-1.5) -- (6.85,1.5);
  \node[hw state, minimum width=1.7cm] (vma) at (8.15,0.0) {映射条目\\vm\_area};
  \node[hw control, minimum width=1.7cm] (pte) at (10.35,0.0) {页表项\\按需填入};
  \draw[hw bus] (6.85,0.0) -- (vma.west);
  \draw[hw return] (vma.east) -- node[hw label, above, yshift=10pt]{缺页时填} (pte.west);
\end{pdfdiagram}
\diagramnote{进程地址空间从低地址到高地址的布局（对应正文 \code{hello} 例）：主程序四个 \code{LOAD} 段从 \code{0x400000} 起依次排开，堆紧随其后向上增长，共享库映在 \code{0x7f...} 区间，栈在最高处向下增长，各区之间隔着未映射的空洞。右侧是同一张地图的折算链：装载时内核把 \code{LOAD} 段折算成映射条目登记下来，页表项则等某页第一次被访问、触发缺页时才填入——装载只建映射，不读字节。}

顺带把四件事摆成一张表，作为本主题的总览：

\begin{longtable}{L{0.14\linewidth}L{0.44\linewidth}L{0.32\linewidth}}
\toprule
步骤 & 内核动作 & 可观察处 \\
\midrule
读头部 & 核对魔数与类型，取入口地址与段表 & \code{readelf -h}、\code{readelf -l} \\\tablerowrule
建映射 & 按 \code{LOAD} 段登记映射，不读字节 & \code{/proc/PID/maps} \\\tablerowrule
准备栈 & 复制字符串，摆 \code{argc}/\code{argv}/\code{envp}/auxv & 调试器里查看栈顶 \\\tablerowrule
跳入口 & \code{rip} 指向动态链接器或入口地址 & \code{strace} 中 \code{execve} 后的 \code{mmap} \\
\bottomrule
\end{longtable}

还有两张表容易混淆，要分清：段表（程序头表）是给装载器看的，节表是给链接器和调试器看的。\code{strip} 把节表整个剥掉，程序照样运行——装载只需要段表；反过来，段表若损坏，内核连映射都建不起来。此外 \code{execve} 有一类“不换”的东西：文件描述符、进程号、工作目录跨 exec 保留，地址空间全换而它们不换——shell 的重定向正是靠“fork 之后、exec 之前先换描述符”实现的。

工程含义只有一句：\code{execve} 本身很快，因为它几乎不搬数据。等 \code{rip} 指向 \code{0x401040} 的第一次取指触发缺页，程序的第一页才真正进内存。把装载从搬运改成建映射，是本节后面全部机制的地基。

补一个反例帮助记忆：\code{execve} 失败时（文件不存在、格式不对、权限不足）正常返回 \code{-1}，原进程的地址空间分毫不动——四件事一件都没做。“成功后不返回”因此也可以反着读：只要它返回了，就一定没成功。

\begin{warningbox}
把装载想成“把程序读进内存”是最常见的误判。按这个模型，fork 之后 exec 前要复制全部分页、程序越大启动越慢、共享库无从谈起——每一条都与观察相反。修正办法是把单位想清楚：内核登记的对象是映射条目和页表项，不是字节；字节只在被访问到时才流动。
\end{warningbox}

\topic{写时复制与按需调页}
fork 的语义是复制整个地址空间，而几乎每个 shell 命令的出生都是 fork 紧跟 exec——若真复制，这将是全系统最昂贵的动作之一。两个机制把它变便宜，它们共用同一个驱动器：经典芯片与平台案例的页错误硬件。

写时复制（copy-on-write）：fork 时内核不复制任何页，只把父进程所有私有可写页在双方页表里都改成只读，父子页表指向同一批物理页。此后任何一方写入都会触发页错误——用经典芯片与平台案例的三问来读：页存在（P=1）、是写访问、权限不允许写——内核据此知道这不是真错误，而是该复制了：分配一个新页，把旧页内容抄过去，写入方的页表项指向新页并恢复可写；若某页只剩一个引用方，直接改回可写，连复制都省。读访问永远不触发复制。

把 fork 前后的页表画成三格，写时复制的全部动作就在其中：

\begin{pdfdiagram}
  % 三格：fork 后 / 写入触发 / 拆页后
  \node[hw label] at (1.6,2.62) {fork 后：共享只读};
  \node[hw label] at (5.4,2.62) {写入触发：页错误};
  \node[hw label] at (9.2,2.62) {拆页后：各自独立};
  \draw[BookMuted, line width=0.5pt, dashed] (3.7,-0.6) -- (3.7,2.35);
  \draw[BookMuted, line width=0.5pt, dashed] (7.5,-0.6) -- (7.5,2.35);
  % 格一：fork 后，双方只读共享同一物理页
  \node[hw state, minimum width=1.5cm, minimum height=0.75cm] (f1) at (0.75,1.5) {父页表项\\只读};
  \node[hw state, minimum width=1.5cm, minimum height=0.75cm] (c1) at (0.75,0.2) {子页表项\\只读};
  \node[hw memory, minimum width=1.15cm, minimum height=1.6cm] (p1) at (2.75,0.85) {物理页 A};
  \draw[hw bus] (f1.east) -| (p1.north);
  \draw[hw bus] (c1.east) -| (p1.south);
  % 格二：父写入，撞上只读权限，触发页错误
  \node[hw state, minimum width=1.5cm, minimum height=0.75cm] (f2) at (4.55,1.5) {父页表项\\只读};
  \node[hw state, minimum width=1.5cm, minimum height=0.75cm] (c2) at (4.55,0.2) {子页表项\\只读};
  \node[hw memory, minimum width=1.15cm, minimum height=1.6cm] (p2) at (6.55,0.85) {物理页 A};
  \draw[hw bus] (f2.east) -| (p2.north);
  \draw[hw bus] (c2.east) -| (p2.south);
  \draw[hw danger] (4.55,2.4) -- (f2.north);
  \node[hw label, anchor=west] at (4.72,2.15) {父写入};
  % 格三：拆页后，父指新页可写，子仍指旧页
  \node[hw state, minimum width=1.5cm, minimum height=0.75cm] (f3) at (8.35,1.5) {父页表项\\可写};
  \node[hw state, minimum width=1.5cm, minimum height=0.75cm] (c3) at (8.35,0.2) {子页表项\\只读};
  \node[hw memory, minimum width=1.15cm, minimum height=0.75cm] (p3n) at (10.35,1.5) {新页 A$'$};
  \node[hw memory, minimum width=1.15cm, minimum height=0.75cm] (p3o) at (10.35,0.2) {物理页 A};
  \draw[hw bus] (f3.east) -- (p3n.west);
  \draw[hw bus] (c3.east) -- (p3o.west);
\end{pdfdiagram}
\diagramnote{写时复制的三幕。fork 后：父子页表项指向同一批物理页，双方都标只读，复制量为零。写入触发：任何一方写共享页都会撞上只读权限，页错误硬件把这次写截住，内核据此知道该拆页了。拆页后：内核分配新页、抄好内容，写入方的页表项改指新页并恢复可写，另一方仍指旧页——复制只发生在被写的那一页上。}

算一笔账：父进程有 5120 个可写私有页，共 $5120\times4\ \text{KiB}=20\ \text{MiB}$；立即复制要搬 20 MiB 并改 5120 条页表项，写时复制只改 5120 条页表项。若子进程在 exec 前只写 200 页，实际复制 $200\times4\ \text{KiB}=800\ \text{KiB}$，约为 20 MiB 的 $4\%$；而 fork 后立刻 \code{execve} 的场合，子进程写的页屈指可数，复制量几乎为零。

历史上这条路走了三段：早期 Unix 的 fork 真的逐页复制，代价随地址空间线性增长；BSD 曾引入 \code{vfork} 作为权宜——子进程借用父进程的地址空间直到 exec，父进程全程冻结，语义脆弱；写时复制普及之后，\code{vfork} 这类绕道基本退场，fork 回到“语义最贵、实现最便宜”的状态。

按需调页（demand paging）是同一思路在装载一侧的应用：上一主题说装载只登记映射，映射的兑现靠缺页。CPU 访问一个存在位为 0 的页，触发页错误，CR2 保存故障地址；内核查映射表，发现该地址属于某文件的某段，就分配物理页、从磁盘读入对应文件页、填好页表项，再重新执行被中断的那条指令。对程序完全透明：它只看到访问成功，只是这一次慢了若干数量级。

再算一笔账：一个 200 MiB 映像共 $200\ \text{MiB}\div4\ \text{KiB}=51200$ 页的程序，启动到 \code{main} 实际触及约 500 页，即 $500\times4096=2\,048\,000$ 字节，约 2 MiB——整个映像的 $1\%$。这就是大程序启动快的原因：启动成本与实际触及的页数成正比，与文件总大小几乎无关。没碰的页不读，几百 MiB 的程序也能在一秒内开窗。

缺页的三种结局值得并排看，因为它们是同一条检查链的三个出口：

\begin{longtable}{L{0.16\linewidth}L{0.44\linewidth}L{0.30\linewidth}}
\toprule
结局 & 内核的发现与动作 & 程序看到什么 \\
\midrule
文件页 & 地址在某文件映射内，从文件读入该页 & 访问成功，仅这一次慢 \\\tablerowrule
匿名零页 & 地址在 \code{.bss}、堆或栈映射内，挂零页 & 访问成功，读到 0 \\\tablerowrule
非法访问 & 地址不属于任何映射 & 收到 SIGSEGV（第四节） \\
\bottomrule
\end{longtable}

按需调页还有一个对偶动作：换出。物理页紧张时，内核把久未访问的匿名页写到交换区，把对应页表项标为不存在；下次访问触发缺页，再把它读回来。于是“页错误”成了统一的内存入口——文件页、零页、交换页都从它进来，虚拟内存因此可以大于物理内存。每个物理页上的引用计数则回答另一个问题：这页被几方共享，写时复制要不要真拆。

\begin{longtable}{L{0.14\linewidth}L{0.28\linewidth}L{0.26\linewidth}L{0.22\linewidth}}
\toprule
方案 & 复制发生在何时 & 复制多少 & 页表改动 \\
\midrule
立即复制 & fork 当场 & 全部可写页 & 新建整套 \\\tablerowrule
写时复制 & 第一次写某页 & 仅被写的页 & 双方先改只读 \\\tablerowrule
按需调页 & 第一次访问某页 & 仅被访问的页 & 缺页时填一条 \\
\bottomrule
\end{longtable}

缺页成本可以按页的来源分档，这也是 \code{/usr/bin/time -v} 里两类缺页计数的含义：

\begin{longtable}{L{0.22\linewidth}L{0.40\linewidth}L{0.28\linewidth}}
\toprule
类型 & 页从哪来 & 典型量级 \\
\midrule
次要缺页（minor） & 页已在内存缓存，只需填页表项 & 亚微秒 \\\tablerowrule
主要缺页（major） & 从磁盘读入文件页 & 几十微秒（SSD） \\\tablerowrule
换入 & 从交换区读回 & 几十到几百微秒 \\
\bottomrule
\end{longtable}

\begin{lstlisting}
$ /usr/bin/time -v ./hello
hello, world
        ...
        Minor (reclaiming a frame) page faults: 152
        Major (requiring I/O) page faults: 3
\end{lstlisting}

内核还会顺势多走一步：缺页读盘时通常顺带读入后续若干页（预读），因为程序访问文件页往往有明显的顺序性——用一次磁盘事务换掉后面的若干次缺页，这正是存储与设备事务相关章节“突发传输比单拍省事务”的同一思想。观察侧的对应概念是常驻集（RSS）：进程实际占着的物理页数。一个 200 MiB 映像的程序启动后 RSS 只有几 MiB，这个差值就是按需调页最直白的观察。

两个机制是同一骨架的两端：都把工作推迟到第一次真正需要时，都由页错误驱动。工程推论同样直接：缺页对程序不可见但真实耗时——磁盘取页约几十微秒，页仍在内存缓存则亚微秒——于是有了预热、内存锁定、大页这类“提前把页备好”的手段；一张 2 MiB 大页把 $2\ \text{MiB}\div4\ \text{KiB}=512$ 次缺页合并成 1 次，存储与外设队列案例谈延迟与队列的视角在这里同样适用。

\topic{动态库的优势与代价}
“链接：重定位与地址形成”一节的链接把全部代码装订进可执行文件，那是静态链接。动态链接把一部分装订推迟到装载时：可执行文件里只记录“我需要 \code{libc.so.6} 里的 \code{printf}”，真正的代码躺在系统里所有进程共享的库文件中。用 \code{readelf -d} 能看到这份“我需要谁”的清单（动态段，输出为模拟节选）：

\begin{lstlisting}
$ readelf -d hello | head -4
Dynamic section at offset 0x2dd8 contains 25 entries:
  Tag        Type                         Name/Value
 0x0000000000000001 (NEEDED)             Shared library: [libc.so.6]
\end{lstlisting}

收益一，共享代码页：libc 的代码段在全部进程间共享一份物理页。算一笔账：libc 代码段约 1.6 MiB（上面 maps 里的 \code{0x19c000} 字节），系统里 300 个进程都用它；静态链接意味着 300 份副本共约 480 MiB，动态链接意味着内存里一份 1.6 MiB（每进程另有少量私有数据页），差约 300 倍。共享成立的条件是代码页只读：谁也不写它，写时复制就不会把共享页拆成私有页。第五个主题会看到，这条约束反过来决定了位置无关代码的形状。

把“共享一份代码页”画成页表指向图，三百倍的差距就写在两组指针上：

\begin{pdfdiagram}
  % 左列：两个进程的页表
  \node[hw state, align=center, minimum width=3.4cm] (pa) at (0,1.4) {进程 A 页表\\代码页项（只读）\\虚拟 0x7f9a3c090000};
  \node[hw state, align=center, minimum width=3.4cm] (pb) at (0,-1.4) {进程 B 页表\\代码页项（只读）\\虚拟地址另一基址};
  % 右列：物理页
  \node[hw memory, align=center, minimum width=3.6cm, minimum height=2.6cm] (phys) at (8.6,0) {物理内存\\libc 代码页一份\\约 1.6 MiB};
  \draw[hw bus] (pa.east) -- (5.2,1.4) -- (5.2,0.5) -- ($(phys.west)+(0,0.5)$);
  \draw[hw bus] (pb.east) -- (5.2,-1.4) -- (5.2,-0.5) -- ($(phys.west)+(0,-0.5)$);
  \node[hw label, fill=white, inner sep=1pt] at (4.3,0) {同一批物理页};
  % 私有数据页对照
  \node[hw compute, align=center, minimum width=3.6cm] (priv) at (8.6,-2.6) {各进程私有数据页\\每进程一份};
  \node[hw label, align=center] at (2.4,-2.9) {300 个进程：静态 300 份约 480 MiB，\\动态链接内存里一份 1.6 MiB};
\end{pdfdiagram}
\diagramnote{动态库的共享代码页。两个进程的页表各自把不同的虚拟地址映射到同一批 libc 物理页——代码段约 1.6 MiB，全系统只有一份；300 个进程时，静态链接要 300 份副本共约 480 MiB，动态链接只要一份（每进程另有少量私有数据页），差约 300 倍。共享成立的条件是页面只读：谁也不写它，写时复制就不会把共享页拆成私有页。}

收益二，升级不重链：修复 libc 的一个缺陷，只需替换库文件本身，下次启动的程序自动用新版；静态链接则必须重新链接并重新分发每一个程序。安全补丁的分发成本因此从“重发所有程序”变成“重发一个库”。文件名里的“\code{.6}”是主版本号：它承诺二进制接口（ABI）不变，只有不兼容时才进位——libc 的主版本号已经二十多年停在 6，于是一个安全更新可以直接替换同名文件。

主版本号管不了的演化，由符号版本接着管：同一个 \code{printf} 在 libc 里带着版本标签 \code{GLIBC\_2.2.5}，可执行文件记录的不仅是“要 \code{printf}”，还有“要哪个版本的 \code{printf}”。于是同一个 \code{libc.so.6} 可以在内部继续演化，老程序拿到的仍是当年那个语义的符号。

\begin{lstlisting}
$ readelf -sW /lib/x86_64-linux-gnu/libc.so.6 | grep -w printf
  1402: 00000000000645e0  137 FUNC   GLOBAL DEFAULT  15 printf@@GLIBC_2.2.5
\end{lstlisting}

还有一条更晚的定址时机值得一提：\code{dlopen}/\code{dlsym} 把“装载一个库”推迟到运行中的任意时刻，插件机制全靠它——浏览器加载扩展、文本编辑器加载语法插件，都是程序自己请动态链接器中途再接一个库进来。链接的三档时机（编译链接时、装载时、运行时）至此全部到齐。

代价同样明确。其一，装载时重定位：库内部对数据和函数的引用在基址确定之前是半成品，装载时要按“基址加偏移”把 GOT 等表项逐条填好，库越多表越长。其二，间接跳转：跨模块调用不能用直接 \code{call}，要经 PLT/GOT 间接一级，每次调用多一次内存间接。其三，依赖管理：“升级不重链”的另一面是“库变了程序没重编”，二进制接口兼容因此成为发行版甩不掉的长期负担。

\begin{longtable}{L{0.16\linewidth}L{0.37\linewidth}L{0.37\linewidth}}
\toprule
方面 & 静态链接 & 动态链接 \\
\midrule
代码份数 & 每进程一份副本 & 全系统共享一份代码页 \\\tablerowrule
升级 & 重新链接并分发每个程序 & 替换库文件即可 \\\tablerowrule
启动前工作 & 无 & 映射库、按基址修补重定位表 \\\tablerowrule
调用开销 & 直接 \code{call} & PLT/GOT 间接一次 \\\tablerowrule
文件体积 & 大，自带全部代码 & 小，只记依赖名 \\
\bottomrule
\end{longtable}

顺带说明，静态链接并没有被淘汰：容器镜像和嵌入式固件反而偏爱它——一个自包含、零依赖的二进制，拷到任何机器上都能跑，这正是用磁盘体积换部署确定性。推到极限，本书末章的教学计算机项目的裸机镜像就是“终极静态链接”：连操作系统都链在同一个映像里，没有任何运行期依赖可言。

\topic{PLT/GOT 的延迟绑定}
动态链接的程序调用 \code{printf} 时，目标地址在链接时还不在本文件里。链接器的办法是两级间接：\code{call} 先到本模块内一小段固定桩（PLT 项），桩里是一条间接跳转，跳的目标写在 GOT 表项里；GOT 是可写数据，真实地址由装载时或运行时填入。这样代码段保持只读可共享，所有需要修补的地址集中在一小张表里。

为什么不干脆在装载时改写 \code{call} 指令的目标？因为那要写代码页：代码页一变脏，写时复制立刻把共享页拆成私有页，共享库和下一主题的 ASLR 一起失效。PLT/GOT 的全部存在理由就是八个字：代码不动，数据可改——代码段一个字节不动，所有可写的定址信息集中到数据段的一张表。

\begin{lstlisting}
        call  printf@plt              ; call site in main

plt0:                                 ; 0x401020
        push  qword ptr [rip+0x2fe2]  ; GOT[1] = link_map
        jmp   qword ptr [rip+0x2fe4]  ; GOT[2] = _dl_runtime_resolve
printf@plt:                           ; 0x401030
        jmp   qword ptr [rip+0x2fe2]  ; GOT[3] at 0x404018
        push  0                       ; relocation index of printf
        jmp   plt0
\end{lstlisting}

把一次 \code{printf} 调用的两级跳走读一遍。第一级是静态的：\code{call printf@plt} 跳到 \code{0x401030}。第二级是动态的：桩里 \code{jmp qword ptr [rip+0x2fe2]} 是 rip 相对寻址，指令长 6 字节，下一条指令在 \code{0x401036}，所以表项地址是 \code{0x401036 + 0x2fe2 = 0x404018}，即 GOT[3]；\code{plt0} 的两条同理，分别指向 GOT[1] 与 GOT[2]，读者可自行核对 \code{0x401026 + 0x2fe2 = 0x404008}、\code{0x40102c + 0x2fe4 = 0x404010}。

第一次调用时，GOT[3] 里装的是 \code{0x401036}——桩里那条 \code{jmp} 的下一条指令自己，这次间接跳等于原地踏步；顺序执行 \code{push 0}（\code{printf} 的重定位下标）和 \code{jmp plt0}，\code{plt0} 把 \code{link\_map} 压栈后跳进动态链接器的解析例程。解析器按下标找到符号名 \code{printf}，在 libc 里查出真实地址（设为 \code{0x7f9a3c0f45e0}），写进 GOT[3]，再直接调进 \code{printf}。第二次调用只剩两条指令：\code{call} 进桩，\code{jmp [0x404018]} 直达 \code{0x7f9a3c0f45e0}，一次内存间接，没有任何解析。

把两次调用画成两条分开的路径，第一次的解析绕路与之后的一跳直达，差别一眼可见：

\begin{pdfdiagram}
  % 上路：首次调用，经解析例程并回写 GOT
  \node[hw label, anchor=west] at (0.05,3.6) {首次调用：解析并回写 GOT};
  \node[hw compute, minimum width=1.7cm] (callA) at (0.9,2.8) {call printf@plt};
  \node[hw compute, minimum width=1.5cm] (pltA) at (2.9,2.8) {PLT 桩};
  \node[hw memory, minimum width=2.2cm] (gotA) at (5.2,2.8) {GOT[3]：\\指回桩内};
  \node[hw compute, minimum width=1.6cm] (plt0) at (7.5,2.8) {plt0\\公共桩};
  \node[hw control, minimum width=1.9cm] (res) at (9.8,2.8) {动态链接器\\解析例程};
  \node[hw compute, minimum width=1.8cm] (fnA) at (9.8,1.3) {printf 本体};
  \draw[hw bus] (callA) -- (pltA);
  \draw[hw bus] (pltA) -- (gotA);
  \draw[hw bus] (gotA) -- (plt0);
  \draw[hw bus] (plt0) -- (res);
  \draw[hw bus] (res.south) -- node[hw label, left]{调入} (fnA.north);
  % 回写 GOT：从解析例程下方绕到 gotA.south，避开 plt0 与 fnA
  \draw[hw return] (res.west) -- (8.7,2.8) -- (8.7,2.0) -- (5.2,2.0) -- (gotA.south);
  \node[hw label, fill=white, inner sep=1pt] at (7.0,2.0) {写回真实地址};
  % 下路：第二次及以后，一次间接直达
  \node[hw label, anchor=west] at (0.05,0.8) {第二次及以后：一跳直达};
  \node[hw compute, minimum width=1.7cm] (callB) at (0.9,0.0) {call printf@plt};
  \node[hw compute, minimum width=1.5cm] (pltB) at (2.9,0.0) {PLT 桩};
  \node[hw memory, minimum width=2.2cm] (gotB) at (5.2,0.0) {GOT[3]：\\真实地址};
  \node[hw compute, minimum width=1.8cm] (fnB) at (8.0,0.0) {printf 本体};
  \draw[hw bus] (callB) -- (pltB);
  \draw[hw bus] (pltB) -- (gotB);
  \draw[hw bus] (gotB) -- node[hw label, below]{直达} (fnB);
\end{pdfdiagram}
\diagramnote{PLT/GOT 两级间接的两条路径。首次调用时 GOT 表项指回桩内，流程经公共桩进入动态链接器的解析例程，查出真实地址后先写回 GOT 表项（虚线）再调入函数本体；此后每次调用只剩 \code{call} 进桩、桩内经 GOT 一次间接跳直达。代码段自始至终不被改写，可写的定址信息全部集中在 GOT 这张表里。}

把首次调用绕进解析例程的那一段再放大一级，ld.so 被触发后的完整路径是“查、写、跳”三步：

\begin{pdfdiagram}
  % 上行：经两级桩进入解析例程
  \node[hw compute, minimum width=1.9cm, align=center] (n1) at (1.0,1.2) {printf@plt\\push 下标 0};
  \node[hw compute, minimum width=1.9cm, align=center] (n2) at (3.3,1.2) {plt0\\push link\_map};
  \node[hw control, minimum width=2.3cm, align=center] (n3) at (5.85,1.2) {动态链接器\\\_dl\_runtime\_resolve};
  \node[hw memory, minimum width=2.6cm, align=center] (n4) at (8.9,1.2) {按 .dynsym 查到\\printf 真实地址};
  % 下行：回写 GOT，再直接调入
  \node[hw memory, minimum width=3.2cm, align=center] (n5) at (8.9,-0.7) {GOT[3] $\leftarrow$ 0x7f9a3c0f45e0\\回写表项};
  \node[hw compute, minimum width=2.6cm, align=center] (n6) at (4.6,-0.7) {跳入 printf 本体\\解析不再发生};
  \draw[hw bus] (n1.east) -- (n2.west);
  \draw[hw bus] (n2.east) -- (n3.west);
  \draw[hw bus] (n3.east) -- (n4.west);
  \draw[hw return] (n4.south) -- node[midway, right=1pt, hw label, fill=white, inner sep=1pt]{写回} (n5.north);
  \draw[hw bus] (n5.west) -- node[midway, above=1pt, hw label, fill=white, inner sep=1pt]{直接调入} (n6.east);
\end{pdfdiagram}
\diagramnote{首次调用时动态链接器的解析路径（数值对应正文 \code{printf} 例）。桩先压入重定位下标 0，公共桩 \code{plt0} 再压入 \code{link\_map}，跳进解析例程；例程按下标找到符号名 \code{printf}，在 libc 中查出真实地址 \code{0x7f9a3c0f45e0}，写进 \code{0x404018} 处的 GOT[3]（虚线回写），最后直接调入函数本体。这张表项填好之后，第二次调用只剩桩内一次间接跳，解析路径被完全绕过。}

\begin{longtable}{L{0.16\linewidth}L{0.37\linewidth}L{0.37\linewidth}}
\toprule
检查点 & 第一次调用 & 第二次及以后 \\
\midrule
GOT[3] 内容 & \code{0x401036}（桩内下一条） & \code{0x7f9a3c0f45e0}（真实地址） \\\tablerowrule
经过路径 & 桩→\code{plt0}→解析器→\code{printf} & 桩→\code{printf} \\\tablerowrule
关键动作 & 查符号、写 GOT[3] & 无 \\\tablerowrule
预期开销 & 一次符号查找，微秒级 & 一次间接跳，纳秒级 \\
\bottomrule
\end{longtable}

PLT 与 GOT 的分工可以列成一张对照：

\begin{longtable}{L{0.16\linewidth}L{0.18\linewidth}L{0.26\linewidth}L{0.30\linewidth}}
\toprule
对象 & 所在段 & 可写性 & 内容由谁何时填 \\
\midrule
PLT 桩 & 代码段 & 只读 & 链接器静态生成，运行期不变 \\\tablerowrule
GOT 表项 & 数据段 & 可写（见下） & 动态链接器按基址与符号填 \\
\bottomrule
\end{longtable}

GOT 的“可写”本身是一个攻击面：谁改写了 GOT[3]，谁就把下一次 \code{printf} 调用拐到自己选的地方。对策叫全 RELRO：装载时用 \code{LD\_BIND\_NOW} 把全部符号解析完，然后把 GOT 所在页改回只读——用启动时间换“运行期不可覆写”。这与上一个主题“代码页只读才能共享”是同一条纪律的延伸：凡是可被写的定址信息，要么推迟到必要时再写，要么写完立刻封死。

这种“第一次才解析”的做法叫懒绑定（lazy binding）。它的收益可以用数字说话：一个中型程序导入约 1200 个动态符号，典型启动路径只用其中约 150 个，懒绑定让启动时省掉约 1050 次符号查找——每次查找都要算哈希、比字符串、写表项。对短生命进程（脚本里一闪而过的每个命令）省得更多：大量导入函数整个生命期一次都没被调用。代价是第一次调用某个函数时有一次不可预测的解析延迟；实时系统可以设 \code{LD\_BIND\_NOW=1}，装载时把全部符号解析完，用启动时间换运行期的确定性。还要记住分工：桩和 GOT 的骨架、GOT 的初值，都是“链接：重定位与地址形成”一节的链接器在静态时就排好的局，动态链接器运行期只做“填表”这一件事。

\topic{ASLR 与装载时重定位}
前面例子的地址全是固定的：代码在 \code{0x401000}，GOT 在 \code{0x404000}。固定地址对攻击者同样是固定的：栈溢出拿到 \code{rip} 控制权后，只要知道 libc 里 \code{system} 或某条指令片段的地址，就能拼接利用。地址空间布局随机化（ASLR）的对策是让每次运行的基址都不同：栈、堆、mmap 区和共享库各加一个随机偏移；编译成位置无关可执行文件（PIE）时，连主程序自己的基址也随机。同一个程序连跑两次：

\begin{lstlisting}
$ ./showaddr                       $ ./showaddr
main   = 0x555555555149            main   = 0x555555e8e149
printf = 0x7f9a3c0f45e0            printf = 0x7f2b8d1a35e0
&buf   = 0x7ffffffde2ac            &buf   = 0x7ffc2a1d59ac
\end{lstlisting}

三个地址全变。换算成攻击者的处境：若某区域基址带 28 位随机位，一次猜中的概率是 $1/2^{28}=1/268\,435\,456\approx3.7\times10^{-9}$；猜错通常意味着目标进程崩溃重来，而崩溃本身就可能暴露行踪。32 位时代有过反面教训：地址空间只有 4 GiB，随机位只有十几位，几千到几万次尝试量级就能爆破；64 位空间让随机位和猜测代价一起变大。ASLR 不是绝对防御——一次地址信息泄漏就能让它失效——但它把一大类利用的成本从“查一次表”抬高到“先泄漏、再猜、还不能错”。

把两次运行的地址空间并排画出，随机化的形状就是“布局不变、基址平移”：

\begin{pdfdiagram}
  \node[hw label] at (0,2.6) {第一次运行};
  \node[hw label] at (7.6,2.6) {第二次运行};
  % 左列（上高下低）
  \node[hw state, align=center, minimum width=4.4cm] (stk1) at (0,1.6) {栈：顶 0x7ffffffde000 附近};
  \node[hw memory, align=center, minimum width=4.4cm] (lib1) at (0,0.4) {libc：基址 0x7f9a3c090000};
  \node[hw compute, align=center, minimum width=4.4cm] (exe1) at (0,-0.8) {主程序：基址 0x555555555000};
  % 右列
  \node[hw state, align=center, minimum width=4.4cm] (stk2) at (7.6,1.6) {栈：顶 0x7ffc2a1d5000 附近};
  \node[hw memory, align=center, minimum width=4.4cm] (lib2) at (7.6,0.4) {libc：基址 0x7f2b8d13f000};
  \node[hw compute, align=center, minimum width=4.4cm] (exe2) at (7.6,-0.8) {主程序：基址 0x555555e8e000};
  % 对应区域的平移
  \draw[hw return] (stk1.east) -- (stk2.west);
  \draw[hw return] (lib1.east) -- node[midway, above=1pt, hw label, fill=white, inner sep=1pt]{各区域基址平移} (lib2.west);
  \draw[hw return] (exe1.east) -- (exe2.west);
  \node[hw label, align=center] at (3.8,-1.9) {布局与代码内容不变，代码页继续只读共享\\变的只是基址：选随机值，按“基址＋偏移”填 GOT};
\end{pdfdiagram}
\diagramnote{ASLR 两次运行对照（对应正文 showaddr 输出）。两次运行的区域顺序与相对布局完全相同，变的只是基址：主程序 0x555555555000 对 0x555555e8e000，libc 0x7f9a3c090000 对 0x7f2b8d13f000，栈顶 0x7ffffffde000 附近对 0x7ffc2a1d5000 附近。换基址不改代码一个字节——PIC 让模块内引用全是 rip 相对、跨模块引用全走 GOT，随机化的全部成本就是选基址加填那张可写的表。}

\begin{longtable}{L{0.18\linewidth}L{0.36\linewidth}L{0.36\linewidth}}
\toprule
区域 & 传统固定位置 & ASLR 之后 \\
\midrule
主程序 & \code{0x400000}（\code{-no-pie}） & PIE 下随机基址 \\\tablerowrule
共享库与 mmap 区 & 固定在 \code{0x7f...} 某处 & 每次运行换基址 \\\tablerowrule
栈 & \code{0x7ffffffff000} 附近 & 栈顶加随机偏移 \\\tablerowrule
堆（brk） & 紧接 \code{.bss} 之后 & 起点加随机偏移 \\
\bottomrule
\end{longtable}

ASLR 便宜的前提是位置无关代码（PIC）。若代码里固定绝对地址，换基址就要逐条修补代码；代码页一被写脏，写时复制立刻把共享页拆成私有页，“300 个进程共享一份 libc”就此作废。PIC 的做法是：模块内引用全部改成 rip 相对（\code{lea rdi, [rip+0x0e31]} 之类），跨模块引用全部走 GOT。换基址因此只剩两步——随机选一个基址，按“基址加偏移”修补那张可写的 GOT——代码页一个字节都不动，继续只读、继续共享。

装载时重定位也正因为集中在一张表里而便宜：表项密度高、缓存友好。一个链接几十个库的 GUI 程序，启动前动态链接器要处理几万条 RELA 记录，每条不过是“基址加符号偏移，写入 GOT[n]”的一次内存写，总量毫秒级。相比之下，若允许对代码段做重定位（老式的 text relocation），每个进程都要先把代码页复制再逐条改写——共享、ASLR、启动速度同时受损，现代工具链因此默认拒绝生成代码重定位。

随机化可以被直接观察，也可以被关掉做对照。\code{setarch -R} 关闭 ASLR 后再跑同一个程序，两次运行的地址逐位相同——这正是“关掉随机化，世界就回到固定地址”的实验：

\begin{lstlisting}
$ setarch -R ./showaddr          $ setarch -R ./showaddr
main   = 0x555555555149            main   = 0x555555555149
printf = 0x7ffff7d645e0            printf = 0x7ffff7d645e0
&buf   = 0x7fffffffe2ac            &buf   = 0x7fffffffe2ac
\end{lstlisting}

把四件事串成一条链：代码页只读，所以可共享；可共享，所以换基址不能改代码；不改代码，ASLR 的成本只剩选基址和填表。本书末章的教学计算机项目的面包板机是这条链的反面：程序拨在哪个地址就运行在哪个地址，没有随机化、没有重定位、没有共享——不是做不到，而是那台机器上根本没有第二个程序需要防。

\begin{classicnote}
经典教材把定址时机分成三档：编译链接时（静态链接，地址在链接后不再变）、装载时（动态链接器映射库并修补重定位表）、运行时（懒绑定在第一次调用时才解析符号）。本节给三档各配了一个可观察的数字：装载只建映射不读字节，所以 \code{execve} 快；懒绑定把上千次符号查找改成按需，所以启动快；PIC 把换基址缩成填一张 GOT，所以 ASLR 几乎免费。记住这张时机表，每个机制省的是哪一笔就都有数了。
\end{classicnote}

\section{四、系统调用与异常控制流}

程序一旦开始执行，就运行在用户态：能算数、能读写自己的内存，但碰不了设备、改不了页表、造不出新进程——这些事必须请内核代办，请代办的门就是系统调用。本节把系统调用、上下文切换和信号串成一条线：它们都是处理器核心相关章节那张“受控的非顺序 PC 更新”表的实例，与存储与设备事务相关章节的硬件中断共用同一副骨架——保存现场、按号定向、服务、返回。

\topic{syscall 指令触发的受控陷入}
用户程序与内核之间有一条硬约定——系统调用 ABI：服务编号放 \code{rax}，参数依次放 \code{rdi}、\code{rsi}、\code{rdx}、\code{r10}、\code{r8}、\code{r9}，然后执行 \code{syscall} 指令，返回值放在 \code{rax}。它与处理器核心相关章节的普通调用约定几乎同形，只有一个反常：第四个参数用 \code{r10} 而不是 \code{rcx}。原因马上会看到——\code{syscall} 指令本身要占用 \code{rcx}，参数只能绕道；这样设计的好处是包装层只需挪一个寄存器，其余位置原样沿用。参数至多六个，现实中系统调用的参数个数在设计时就被压在这条线内。

系统调用编号是系统调用表的稳定索引，可以类比异常向量中的选择号，但它并不是 IDT 向量。x86-64 Linux 的编号表由内核 ABI 定义；已经发布的编号需要保持兼容，即使某项服务废弃也通常保留相应编号或兼容处理，因为已编译程序会把编号写入机器码。

\begin{lstlisting}
; x86-64 Linux, arch/x86/entry/syscalls/syscall_64.tbl (excerpt)
  0   common  read
  1   common  write
  2   common  open
  3   common  close
  9   common  mmap
 57   common  fork
 59   common  execve
231   common  exit_group
\end{lstlisting}

一条 \code{write(1, msg, 13)} 按 ABI 写出来是这样（Intel 语法）：

\begin{lstlisting}
        mov  eax, 1              ; __NR_write
        mov  edi, 1              ; fd = stdout
        lea  rsi, [rip+msg]      ; buf
        mov  edx, 13             ; count
        syscall                  ; controlled trap into kernel
        ; rax = 13 on success, or -errno

msg:    db   "hello, world", 10
\end{lstlisting}

\code{syscall} 指令的硬件动作极少而精：把返回地址（下一条指令的 \code{rip}）装进 \code{rcx}，把 \code{rflags} 装进 \code{r11}，按 \code{IA32\_FMASK} MSR 的掩码清掉 \code{rflags} 的指定位，特权级从 3 升为 0，最后把 \code{IA32\_LSTAR} MSR 里的内核入口地址装进 \code{rip}。指令本身不访问内存，也不查询 IDT；入口地址早在系统启动时就由内核写入 MSR。\cite{intel-sdm}

为什么参数全走寄存器，而不像普通调用那样允许栈传参？因为内核不能信任用户栈指针：\code{syscall} 那一刻 CPU 还在用户栈上，而这个栈指针本身就可能指向不可写的页。六个寄存器是“内核不用解引用任何用户指针就能拿到的全部输入”；参数个数因此被设计时压在这条线内——历史上 32 位 ABI 只有五个寄存器可用，第六个参数被迫走内存，至今还能在老代码里看到那处别扭。

以 \code{write(1, msg, 13)} 为例，把寄存器摆参与内核栈上的 \code{pt\_regs} 画在一张图里，陷入与返回各是一次受控交接：

\begin{pdfdiagram}
  \node[hw state, minimum width=4.2cm, align=center] (ureg) at (2.1,1.0) {用户态摆参（write 例）\\rax ← 1（调用号）\\rdi ← 1，rsi ← msg\\rdx ← 13\\第 4–6 参数：r10、r8、r9\\指令占用：rcx ← rip，r11 ← rflags};
  \node[hw memory, minimum width=4.6cm, align=center] (ptregs) at (8.5,1.0) {当前线程内核栈：pt\_regs\\用户 rsp、rip、rflags\\由入口代码按布局保存\\orig\_rax ← 调用号\\rdi rsi rdx rcx rax r10\\r8 r9 r11 rbx rbp r12–r15};
  \draw[hw bus] (ureg.east) -- node[midway, above=1pt, hw label, fill=white, inner sep=1pt]{陷入：CPL 3→0} (ptregs.west);
  \draw[hw return] (ptregs.south) -- (8.5,-1.5) -- (2.1,-1.5) -- (ureg.south);
  \node[hw label, fill=white, inner sep=1pt] at (5.3,-1.22) {入口代码恢复通用状态；sysret 恢复 rip/rflags 与 CPL};
\end{pdfdiagram}
\diagramnote{系统调用的参数传递与现场保存（\code{write(1, msg, 13)} 例）。用户态按 ABI 摆好 \code{rax} 编号与 \code{rdi}/\code{rsi}/\code{rdx} 参数；\code{syscall} 指令把返回地址装进 \code{rcx}、标志装进 \code{r11}——这正是第四个参数改用 \code{r10} 的原因——并把 \code{rip} 换成 LSTAR 入口；它本身不保存用户 \code{rsp}，也不自动切换栈。入口代码随后切到当前线程的内核栈，把用户态现场连同全部通用寄存器存成 \code{pt\_regs}，内核 C 代码看到的“寄存器”就是这张内存表；最后由 \code{sysret} 恢复 \code{rip}/\code{rflags} 与用户特权级，返回值仍由 \code{rax} 带回，一进一出同一个寄存器。}

把普通调用和系统调用并排，ABI 的同形与差异一眼可见：

\begin{longtable}{L{0.24\linewidth}L{0.32\linewidth}L{0.34\linewidth}}
\toprule
方面 & 普通函数调用 & 系统调用 \\
\midrule
第四个参数 & \code{rcx} & \code{r10}（\code{rcx} 被指令占用） \\\tablerowrule
返回地址 & 由 \code{call} 压到栈上 & 装进 \code{rcx} 寄存器 \\\tablerowrule
特权级 & 不变 & 3→0，返回时 0→3 \\\tablerowrule
使用的栈 & 同一条用户栈 & 指令不换栈；入口代码切到当前线程内核栈 \\\tablerowrule
入口地址 & 固定在指令编码里 & \code{IA32\_LSTAR} MSR \\
\bottomrule
\end{longtable}

与老式的 \code{int 0x80} 对照最能看清这副骨架：\code{int} 要按处理器核心相关章节的“表基址+异常号×项宽”去内存里查 IDT，需要执行门描述符检查、硬件换栈和规定的现场保存；\code{syscall} 使用 MSR 中预置的专用入口，并只保存最小返回状态，把“编号到服务”的分发表 \code{sys\_call\_table} 留给内核软件去查。但两者只是同一副骨架的两种实现：指令显式请求进入特权服务，硬件保存返回位置、切换权限和入口——处理器核心相关章节异常表里“系统调用”那一行，写的就是它。

\begin{longtable}{L{0.14\linewidth}L{0.44\linewidth}L{0.32\linewidth}}
\toprule
骨架阶段 & \code{syscall} 路径上的动作 & 对应前章的谁 \\
\midrule
保存现场 & 硬件完成 \code{rip}→\code{rcx}、\code{rflags}→\code{r11}；入口代码切栈并建立 \code{pt\_regs} & 处理器核心相关章节“保存故障 PC 和原因” \\\tablerowrule
定向 & 硬件入口 = \code{IA32\_LSTAR}；内核再按 \code{sys\_call\_table[rax]} 分发 & 处理器核心相关章节向量表查表；\code{int 0x80} 走 IDT \\\tablerowrule
服务 & \code{sys\_read}/\code{sys\_write} 等操作设备与文件 & 存储与设备事务相关章节 MMIO/DMA 由内核代发 \\\tablerowrule
返回 & \code{sysret}：\code{rcx}→\code{rip}、\code{r11}→\code{rflags}、特权级 0→3 & 存储与设备事务相关章节中断返回恢复现场 \\
\bottomrule
\end{longtable}

这副“保存现场、按号定向、服务、返回”的骨架值得背下来，因为本节后面三个主题全是它的实例：上下文切换是“定时器中断加换一套现场”，信号是“返回用户态之前的检查点”，边界清单则是“哪些服务必须在这副骨架里完成”。

把一次系统调用的完整往返按时间自上而下画成两条泳道，特权级的两次切换就是图上两次越过边界的时刻：

\begin{pdfdiagram}
  % 泳道底框（先画，垫在节点之下）
  \node[hw lane, minimum width=4.4cm, minimum height=6.6cm] at (2.2,1.7) {};
  \node[hw lane, minimum width=4.4cm, minimum height=6.6cm] at (8.0,1.7) {};
  \node[hw label] at (2.2,5.35) {用户态（CPL 3）};
  \node[hw label] at (8.0,5.35) {内核态（CPL 0）};
  % 特权级边界
  \draw[BookRule, line width=0.45pt, dashed] (5.1,-1.6) -- (5.1,5.0);
  \node[hw label] at (5.1,-1.95) {特权级边界};
  % 用户态泳道
  \node[hw compute, minimum width=3.2cm] (u1) at (2.2,4.2) {用户程序摆好编号与参数\\rax、rdi、rsi、rdx};
  \node[hw control, minimum width=3.2cm] (u2) at (2.2,2.9) {syscall 指令\\rip→rcx，rflags→r11};
  \node[hw compute, minimum width=3.2cm] (u3) at (2.2,-1.0) {继续执行下一条指令\\rax = 返回值};
  % 内核态泳道
  \node[hw state, minimum width=3.2cm] (k1) at (8.0,2.9) {LSTAR 入口汇编\\保存用户 rsp、切线程内核栈\\现场存为 pt\_regs};
  \node[hw compute, minimum width=3.2cm] (k2) at (8.0,1.6) {按号分发\\sys\_call\_table[rax]};
  \node[hw compute, minimum width=3.2cm] (k3) at (8.0,0.3) {服务例程\\sys\_write 等};
  \node[hw state, minimum width=3.2cm] (k4) at (8.0,-1.0) {入口代码恢复寄存器与用户 rsp\\sysret：rcx→rip，r11→rflags};
  \draw[hw bus] (u1.south) -- (u2.north);
  \draw[hw bus] (u2.east) -- (k1.west);
  \node[hw label] at (5.2,3.9) {陷入：特权级 3→0};
  \draw[hw bus] (k1.south) -- (k2.north);
  \draw[hw bus] (k2.south) -- (k3.north);
  \draw[hw bus] (k3.south) -- (k4.north);
  \draw[hw bus] (k4.west) -- (u3.east);
  \node[hw label] at (5.1,-0.1) {返回：特权级 0→3};
\end{pdfdiagram}
\diagramnote{一次系统调用往返的泳道图：左为用户态（CPL 3），右为内核态（CPL 0），中间虚线是特权级边界，两次越界即 CPL 3→0 与 0→3 的切换点。\code{syscall} 指令把返回地址与标志装进 \code{rcx}/\code{r11}、把 \code{rip} 换成 MSR 中的内核入口；入口代码切到当前线程内核栈，把全部通用寄存器存成 \code{pt\_regs}，再按 \code{rax} 中的编号查 \code{sys\_call\_table} 分发；服务完成后，入口代码先恢复通用寄存器和用户栈指针，再由 \code{sysret} 完成最后的返回状态切换。内核全程不必解引用用户栈指针。}

\topic{陷入与返回的上下文}
\code{syscall} 指令只切换了最小状态，完整的上下文交接由内核入口代码完成：\code{swapgs} 换出内核数据基址；先把用户 \code{rsp} 保存到每 CPU 的入口临时状态，再从当前任务信息取得该线程的内核栈指针。每个线程都有独立内核栈，具体大小由内核配置决定，专门承接它进入内核后的现场。TSS 的 \code{RSP0} 用于中断门等发生特权级切换时的硬件换栈，并不是 \code{SYSCALL} 自动换栈的来源；然后把用户态的 \code{rsp}、\code{rcx}（用户 \code{rip}）、\code{r11}（用户 \code{rflags}）和全部通用寄存器按固定布局存成一张 \code{pt\_regs}。此后内核 C 代码看到的“寄存器”，其实是这张内存表。返回时逆序恢复，\code{sysret} 用 \code{rcx}/\code{r11} 还原 \code{rip}/\code{rflags}，并把特权级降回 3。

\begin{lstlisting}
; pt_regs on the current thread's kernel stack, high to low address:
;   rsp(user), rip(user), rflags(user)           <- arranged by entry code
;   orig_rax                                     <- syscall number
;   rdi rsi rdx rcx rax r8 r9 r10 r11 rbx rbp
;   r12 r13 r14 r15                              <- saved by entry code
\end{lstlisting}

为什么每个线程需要一条独立内核栈？因为线程可能在内核里被切走：它服务到一半被定时器打断，现场得有个各自安放的地方。内核栈保存该线程进入内核后的调用链与现场，栈顶指针保存在当前线程的任务控制结构中——下一个主题会看到，上下文切换保存的全部东西，最后都归结为这个指针。若当前内核配置使用 16 KiB 栈，这一容量同时形成明确约束：内核态不能深递归，也不能在栈上放大数组，否则栈溢出会破坏当前线程的现场。

把一次 \code{read(0, buf, 4096)} 的完整往返走一遍。\code{strace} 里它只有一行：

\begin{lstlisting}
$ strace -e read ./hello
read(0, "hello\n", 4096)               = 6
\end{lstlisting}

第一步，用户态摆好 ABI：\code{rdi=0}、\code{rsi=buf}、\code{rdx=4096}、\code{rax=0}，执行 \code{syscall}。到这一步为止，一切和普通函数调用的准备没有区别——区别从下一步开始。

第二步，陷入：特权级升为 0，\code{rip} 换成 LSTAR 入口，入口代码建 \code{pt\_regs}，按 \code{rax=0} 查 \code{sys\_call\_table} 分发到 \code{sys\_read}。注意返回地址此刻在 \code{rcx} 里、被存进 \code{pt\_regs}，而不在用户栈上——内核从始至终不需要解引用用户栈指针。

第三步，服务：经虚拟文件层到终端驱动，键盘缓冲里已有 6 字节，复制到 \code{buf}——复制前内核先检查 \code{buf} 所在页可写，不可写就不复制，直接返回 \code{-14}。设备寄存器、DMA 描述符这些存储与设备事务相关章节的对象，只有在这一步、只有内核态才能碰。

第四步，返回：\code{rax=6}，\code{sysret} 回用户态，libc 包装看到非负值，直接把它作为返回值交出。返回值用 \code{rax} 不是巧合：编号进内核时用的就是 \code{rax}，一进一出，同一个寄存器完成交接，调用者连第二个寄存器都不用读。

\begin{lstlisting}
before:  rax=0  rdi=0  rsi=0x7ffffffde2a0  rdx=4096
trap:    rcx=rip(next)  r11=rflags  CPL 3->0  rip=LSTAR
return:  rax=6
\end{lstlisting}

错误约定值得单独说清。内核用 \code{-errno} 表示失败：\code{rax} 落在 $[-4095,-1]$ 区间即出错，例如 \code{-14=-EFAULT}（\code{buf} 不可写）、\code{-9=-EBADF}（描述符无效）。判断只需一次无符号比较：\code{-14} 的无符号形式是 \code{0xFFFFFFFFFFFFFFF2}，不小于 \code{-4095} 的无符号形式 \code{0xFFFFFFFFFFFFF001}，一条 \code{cmp} 就分出成败。glibc 的包装函数负责翻译：检出该区间就把 \code{errno} 置为 14、返回 \code{-1}——C 程序看到的是 \code{read} 返回 \code{-1} 且 \code{errno==EFAULT}，内核内部的 \code{-14} 由包装层转换，不直接暴露给程序。

为什么连 \code{buf} 都要检查？因为内核不能信任用户态递来的任何指针：它可能指向不可写的页，也可能故意指向内核地址。每个系统调用开头的参数检查，是内核自保的第一课，也是“边界”二字在服务一侧的具体形态。

常见错误码可以列成一张小表，以后读 \code{strace} 输出时会反复遇到它们：

\begin{longtable}{L{0.10\linewidth}L{0.18\linewidth}L{0.62\linewidth}}
\toprule
数值 & 名字 & 典型触发场景 \\
\midrule
\code{-1} & \code{EPERM} & 操作不被允许，如向别人的进程发信号 \\\tablerowrule
\code{-2} & \code{ENOENT} & 路径不存在，\code{open} 失败 \\\tablerowrule
\code{-9} & \code{EBADF} & 描述符无效，如 \code{read} 一个已关闭的 fd \\\tablerowrule
\code{-11} & \code{EAGAIN} & 非阻塞读暂时无数据 \\\tablerowrule
\code{-14} & \code{EFAULT} & 指针指向不可访问的页 \\
\bottomrule
\end{longtable}

\code{pt\_regs} 里的 \code{orig\_rax} 还藏着一条与信号的衔接：慢系统调用（比如等键盘的 \code{read}）若被信号打断，服务例程会返回一个内部码 \code{ERESTART}；处理函数若按 \code{SA\_RESTART} 安装，内核就在 \code{sigreturn} 后用 \code{orig\_rax} 里的编号把这次调用重新发起，用户态毫无察觉，否则 libc 把 \code{-EINTR} 交给程序，由程序自己决定要不要重试。系统调用能被“重启”，前提正是完整现场都在 \code{pt\_regs} 里。

工程含义是一笔时间账：一次系统调用往返的典型开销在百纳秒级，比普通函数调用贵约两个数量级——贵在特权切换、\code{pt\_regs} 保存、查表分发和参数检查。后世的优化几乎全部围绕“少进几次内核”：stdio 缓冲、\code{readv}/\code{writev} 批量提交、\code{mmap} 免复制、vDSO 和 \code{io\_uring}，最后一个主题还会回到这张清单。

\topic{上下文切换的完整步骤}
系统调用是程序主动进内核，定时器中断则是内核定期不请自来——走存储与设备事务相关章节那条链路：定时器拉起 IRQ，中断控制器判优，CPU 在指令边界接收，查向量表进服务程序。服务程序发现当前进程的时间片用完，就做一次上下文切换：把整套程序员可见状态换成另一个进程的。完整步骤可以写成一条 trace：

\begin{lstlisting}
// 注释（推进）：按行追踪数据来源、经过的部件和最终写入目标。
// 注释（边界）：只有标出的时钟沿、阶段或异常入口会改变体系结构可见状态。
timer interrupt arrives while A runs
  CPU:   push rip/cs/rflags/rsp/ss of A; enter ISR
  ISR:   push A's general registers to A's kernel stack
  A.state = READY
  B = pick_next()                        ; scheduler decision
  if (B.mm != A.mm) CR3 = B.pgd          ; switch address space
  rsp = B.saved_kernel_sp                ; switch kernel stack
  pop B's general registers
  iretq                                  ; resume B at its saved rip
\end{lstlisting}

逐步核对。第一步保存 A：中断已经自动压好 \code{rip}/\code{cs}/\code{rflags}/\code{rsp}/\code{ss}，入口代码补齐通用寄存器——A 的全部可见状态此刻躺在它自己的内核栈上，进程控制块里只要存一个栈顶指针。第二步选下一个：调度器从就绪队列挑出 B——“换谁”是软件策略，“怎么换”才是这里的硬件约定，两者泾渭分明。第三步切地址空间：把 B 的页表基址写进 CR3（经典芯片与平台案例：CR3 指向顶级页表）——从这一刻起，同一个虚拟地址 \code{0x401000} 指向的是 B 的代码；副作用随之而来，TLB 里 A 的翻译全部作废，除非启用 PCID（PCID 标签打在 TLB 项上，取自 CR3 低 12 位），B 恢复后的前几千次访存都会慢一拍。第四步恢复 B 的寄存器并 \code{iretq}——B 从自己的视角看，只是“上一次陷入返回得特别久”。

把这条 trace 画成泳道图，“换一套现场”就是三条道之间的一次往返：

\begin{pdfdiagram}
  % 泳道标签与分隔虚线
  \node[hw label] at (0,2.9) {进程 A（用户态）};
  \node[hw label] at (5.2,2.9) {内核};
  \node[hw label] at (10.4,2.9) {进程 B（用户态）};
  \draw[draw=BookMuted, dashed, line width=0.5pt] (2.9,2.55) -- (2.9,-2.6);
  \draw[draw=BookMuted, dashed, line width=0.5pt] (7.5,2.55) -- (7.5,-2.6);
  % 事件序列（时间向下）
  \node[hw state, align=center] (arun) at (0,1.9) {A 运行中};
  \node[hw control, align=center, minimum width=3.0cm] (save) at (5.2,1.9) {保存 A 现场\\压上 A 的内核栈};
  \node[hw compute, align=center, minimum width=3.0cm] (pick) at (5.2,0.65) {pick\_next() 选出 B\\A.state = READY};
  \node[hw compute, align=center, minimum width=3.0cm] (cr3) at (5.2,-0.6) {CR3 $\leftarrow$ B 页表基址\\rsp $\leftarrow$ B 内核栈顶};
  \node[hw control, align=center, minimum width=3.0cm] (rest) at (5.2,-1.85) {恢复 B 现场\\iretq};
  \node[hw state, align=center] (brun) at (10.4,-1.85) {B 从断点继续};
  \draw[hw bus] (arun.east) -- node[pos=0.25, above=1pt, hw label, fill=white, inner sep=1pt]{定时器 IRQ} (save.west);
  \draw[hw bus] (save.south) -- (pick.north);
  \draw[hw bus] (pick.south) -- (cr3.north);
  \draw[hw bus] (cr3.south) -- (rest.north);
  \draw[hw bus] (rest.east) -- node[pos=0.65, above=1pt, hw label, fill=white, inner sep=1pt]{返回用户态} (brun.west);
  \node[hw label, align=center] at (0,-0.6) {A 冻结：全部状态\\凝为一个栈指针};
\end{pdfdiagram}
\diagramnote{上下文切换泳道（对应正文 trace）。定时器中断在指令边界把 CPU 从 A 带进内核：A 的现场压上它自己的内核栈，调度器选出 B 后写入 B 的页表基址（CR3）、把栈指针切到 B 的内核栈，再弹出 B 的寄存器并 iretq——返回落进 B 的用户态，B 只觉得上次陷入返回得特别久。被切走的 A，全部状态归结为进程控制块里的一个栈指针，等下次被选中时原路恢复。}

定时器不是唯一的切换入口：进程阻塞在 \code{read} 上等键盘、主动 \code{sched\_yield} 让出、等待互斥锁，走的都是同一条“保存、选下一个、恢复”的路径。入口有多个，出口只有一个——沿目标进程的内核栈恢复现场并返回。

\begin{longtable}{L{0.22\linewidth}L{0.38\linewidth}L{0.30\linewidth}}
\toprule
要素 & 保存在哪、切到哪里 & 检查点 \\
\midrule
通用寄存器与 \code{rip} & 各自内核栈顶；进程控制块只存栈指针 & 恢复后 \code{rip} 指向被切走时的下一条 \\\tablerowrule
地址空间 & CR3 写入新页表基址 & 同一虚拟地址从此属于新进程 \\\tablerowrule
内核栈 & 栈指针切到目标栈 & 返回路径沿目标栈展开 \\
\bottomrule
\end{longtable}

算一笔账。保存加恢复的寄存器约 20 个 64 位值，内存流量 $2\times20\times8=320$ 字节，纳秒级；真正的成本在 CR3 切换之后——TLB 和 cache 变冷，使一次进程切换的典型总成本在微秒级，约是一次系统调用往返的十倍。再看频率：设时间片 3 ms，一个核每秒至多切换约 330 次，每次 2 微秒，总开销约 660 微秒，占可用时间的千分之一不到——切换很便宜，而“看起来便宜”是因为时间片把它摊薄了。

有一个边角案例能检验这条 trace 是否真的成立：新进程从来没有“上次被切走”，它的第一次运行从哪来？答案是 fork 时内核手工搭出来的——按上一主题的格式伪造一份 \code{pt\_regs} 和内核栈，把 \code{rip} 设成父进程 fork 返回后的下一条，把 \code{rax} 设成 0。于是子进程第一次被调度时，走的完全是“恢复现场并返回”的老路，只是现场是造的。同一份代码在父子俩身上跑出两个返回值（父得 pid、子得 0），差别就写在这份手工现场里。

定时器中断接手调度，本身是一段历史进步：早期系统是协作式的，进程靠自觉定期让出，一个死循环就能拖死整机；抢占式调度把主动权收归内核——存储与设备事务相关章节那个“除了报时还会发起中断”的定时器，另一个身份就是操作系统的心跳。切换的频繁程度可以直接观察：\code{vmstat} 的 \code{cs} 列是每秒上下文切换次数，\code{in} 列是每秒中断次数，两者同量级，因为大部分切换正是中断带进来的。

\begin{lstlisting}
$ vmstat 1 3
procs -----------memory---------- ---swap-- -----io---- -system-- ------cpu-----
 r  b   swpd   free   buff  cache   si   so    bi    bo   in   cs us sy id wa st
 1  0      0 6123456 123456 2345678    0    0    10    20  210  512  5  3 92  0  0
 0  0      0 6123300 123456 2345680    0    0     0    12  198  487  4  2 94  0  0
\end{lstlisting}

TLB 作废的代价也在被硬件压缩：带 PCID 的 CPU 给每条 TLB 项打上地址空间标签，切 CR3 时不必清空 TLB，新旧进程的翻译按标签各归各——这正是存储与设备事务相关章节 cache 用 tag 区分归属的思想，在地址转换缓存上的重演。

\begin{longtable}{L{0.20\linewidth}L{0.40\linewidth}L{0.30\linewidth}}
\toprule
切换入口 & 触发原因 & 典型例子 \\
\midrule
被动切换 & 定时器中断，时间片用完 & 计算密集的循环被切走 \\\tablerowrule
主动切换 & 系统调用内阻塞，等资源 & \code{read} 等键盘输入 \\\tablerowrule
让出切换 & 显式请求调度器换人 & \code{sched\_yield} \\
\bottomrule
\end{longtable}

两个推论。其一，\hwkey{进程切换就是把整套可见状态换一遍}：寄存器、地址空间、内核栈，一个不能少；而文件描述符表、信号处理表这类“属于进程但不属于寄存器”的状态挂在进程控制块上，换进程时自然跟着换，不需要逐条保存。其二，同进程的两个线程互切不用写 CR3——它们共享地址空间，TLB 不必作废——这正是线程比进程“轻”的硬件根源。

\topic{信号是软件级中断}
中断让设备打断 CPU，信号让内核（或另一个进程）打断一个进程。两者同构：信号编号对应中断号，\code{sigaction} 安装的处理表对应 IDT，\code{sigprocmask} 对应中断屏蔽字，pending 位图对应中断挂起寄存器，\code{sigreturn} 对应 \code{iret}。差别只在发起者和投递点：中断由硬件在指令边界投递，信号由内核在“即将从内核返回用户态”这个固定检查点投递——每个进程从系统调用或中断返回前，都要过这一关。

Ctrl-C 是最好的一生案例，四步走完。第一步，按键：终端把 \code{0x03} 字节经存储与设备事务相关章节那条链路送进内核——设备锁存状态、拉起 IRQ、中断控制器判优、CPU 在指令边界接收、键盘服务程序读入扫描码。第二步，行律解释：tty 驱动发现 \code{0x03} 是控制字符 INTR，不把它放进输入缓冲，而是向前台进程组的每个进程置 SIGINT 的 pending 位。第三步，投递：目标进程下次返回用户态前，内核发现 SIGINT 未被屏蔽且注册过处理函数，就在用户栈上搭一帧信号帧（保存被打断时的 \code{rip}、\code{rsp}、\code{rflags} 和通用寄存器），把 \code{rip} 改为处理函数入口、\code{rsp} 指向新帧，然后才执行那次“返回”——返回去的地方已经是处理函数。第四步，收尾：处理函数执行完调用 \code{sigreturn}，内核按帧恢复现场，进程从被打断处继续；若进程从未注册处理函数，默认动作是终止——内核在投递点直接回收进程。Ctrl-C“杀掉”前台程序，走的就是这完整的四步。

把 Ctrl-C 的四步画成一条投递路径，“软件级中断”五个字就写在图上：

\begin{pdfdiagram}
  % 第一行：发起与置位
  \node[hw state, align=center] (key) at (0,1.3) {键盘\\0x03};
  \node[hw compute, align=center, minimum width=2.6cm] (tty) at (3.1,1.3) {tty 行律\\认出 INTR};
  \node[hw control, align=center, minimum width=3.0cm] (pend) at (6.7,1.3) {前台进程组各进程\\置 SIGINT pending 位};
  % 第二行：投递
  \node[hw control, align=center, minimum width=2.6cm] (chk) at (3.1,-1.0) {返回用户态前检查\\未屏蔽且已注册才投递};
  \node[hw memory, align=center, minimum width=2.6cm] (frame) at (6.7,-1.0) {用户栈上\\搭信号帧};
  \node[hw state, align=center, minimum width=2.6cm] (hdl) at (10.0,-1.0) {rip $\leftarrow$ 处理函数\\rsp $\leftarrow$ 新帧};
  \draw[hw bus] (key.east) -- node[midway, above=1pt, hw label, fill=white, inner sep=1pt]{扫描码} (tty.west);
  \draw[hw bus] (tty.east) -- node[midway, above=1pt, hw label, fill=white, inner sep=1pt]{控制字符} (pend.west);
  \draw[hw bus] (pend.south) -- (6.7,0.15) -- (3.1,0.15) -- (chk.north);
  \node[hw label, fill=white, inner sep=1pt] at (4.9,0.42) {下次返回前被检查};
  \draw[hw bus] (chk.east) -- (frame.west);
  \draw[hw bus] (frame.east) -- (hdl.west);
  \draw[hw return] (hdl.south) -- (10.0,-2.3) -- (6.7,-2.3) -- (frame.south);
  \node[hw label, fill=white, inner sep=1pt] at (8.35,-2.56) {sigreturn：按帧恢复现场};
\end{pdfdiagram}
\diagramnote{信号投递路径（Ctrl-C 例）。按键字节 0x03 经键盘服务程序送进内核，tty 行律认出控制字符 INTR，不放进输入缓冲，而是向前台进程组每个进程置 SIGINT（2 号）的 pending 位；目标进程下次返回用户态前，内核检出该位未被屏蔽且处理函数已注册，就在用户栈上搭信号帧、把 rip 改为处理函数入口、rsp 指向新帧——那次“返回”落进处理函数；处理函数执行完 sigreturn，内核按帧恢复现场，进程从被打断处继续。}

四步里搭在用户栈上的那帧信号帧值得单独画出来，处理函数能返回被打断处，全靠帧里保存的现场：

\begin{pdfdiagram}
  \node[hw label] at (2.3,3.5) {用户栈（高地址在上，向下增长）};
  \node[hw compute, minimum width=3.6cm, align=center] (old) at (2.3,2.6) {被打断处的栈帧\\原 rsp 所指};
  \node[hw memory, minimum width=4.6cm, align=center] (frame) at (2.3,0.5) {信号帧（内核搭建）\\保存的 rip、rsp、rflags\\全部通用寄存器\\返回地址 → sigreturn 跳板};
  \draw[BookMuted, line width=0.62pt, ->] (old.south) -- node[midway, right=1pt, hw label, fill=white, inner sep=1pt]{内核压入} (frame.north);
  \node[hw control, minimum width=3.2cm, align=center] (hdl) at (8.7,0.5) {信号处理函数\\rip ← 处理函数入口};
  \draw[hw bus] (frame.east) -- node[midway, above=1pt, hw label, fill=white, inner sep=1pt]{投递} (hdl.west);
  \draw[hw return] (hdl.south) -- (8.7,-1.2) -- (2.3,-1.2) -- (frame.south);
  \node[hw label, fill=white, inner sep=1pt] at (5.5,-1.2) {sigreturn：按帧恢复现场，回到被打断处};
\end{pdfdiagram}
\diagramnote{用户栈上的信号帧结构（对应 Ctrl-C 四步的第三步）。内核在目标进程的用户栈上搭一帧：保存被打断时的 \code{rip}、\code{rsp}、\code{rflags} 与全部通用寄存器，帧里的返回地址指向 \code{sigreturn} 跳板；随后 \code{rip} 改为处理函数入口、\code{rsp} 指向新帧，那次“返回用户态”便落进处理函数。处理函数执行完经跳板发起 \code{sigreturn}，内核按帧恢复现场，进程从被打断处继续。若用户栈本身已不可用（栈溢出招来 SIGSEGV 的情形），处理函数需要一条事先用 \code{sigaltstack} 备好的备用栈。}

两个信号被故意排除在约定之外：SIGKILL（9）和 SIGSTOP（19）不可捕获、不可屏蔽、不可忽略。理由很硬：若失控进程能把所有信号都拦下，系统就没有回收它的最后手段——\code{kill -9} 之所以总是有效，靠的就是这条“保留一扇不可锁的门”的约定。

\begin{longtable}{L{0.16\linewidth}L{0.37\linewidth}L{0.37\linewidth}}
\toprule
概念 & 硬件中断（第四、五章） & 信号 \\
\midrule
编号 & 中断号/向量号 & 信号编号（SIGINT=2，SIGSEGV=11） \\\tablerowrule
入口表 & IDT：基址+号×项宽 & 每进程 \code{sigaction} 表 \\\tablerowrule
屏蔽 & IF 位/中断控制器屏蔽位 & \code{sigprocmask} 屏蔽字 \\\tablerowrule
挂起 & IRR 挂起寄存器 & pending 位图 \\\tablerowrule
返回 & \code{iret} & \code{sigreturn} 系统调用 \\\tablerowrule
投递点 & 指令边界 & 内核返回用户态之前 \\
\bottomrule
\end{longtable}

常见信号的编号和默认动作可以各看一张表。编号表用 \code{kill -l} 就能列出来：

\begin{lstlisting}
$ kill -l (excerpt)
 1) SIGHUP    2) SIGINT    9) SIGKILL  11) SIGSEGV
13) SIGPIPE  14) SIGALRM  15) SIGTERM  18) SIGCONT
19) SIGSTOP  20) SIGTSTP
\end{lstlisting}

\begin{longtable}{L{0.20\linewidth}L{0.24\linewidth}L{0.46\linewidth}}
\toprule
默认动作 & 代表信号 & 含义 \\
\midrule
终止 & SIGINT、SIGTERM & 内核直接回收进程 \\\tablerowrule
终止并转储 & SIGSEGV、SIGILL & 回收，另留下 core 文件供调试 \\\tablerowrule
忽略 & SIGCHLD、SIGURG & 默认什么也不做 \\\tablerowrule
停止 & SIGSTOP、SIGTSTP & 冻结进程，等 SIGCONT 唤醒 \\
\bottomrule
\end{longtable}

Ctrl-C 的四步只换了发起者，就能解释一整族终端控制：Ctrl-Z 发出的是 SIGTSTP，前台进程被“停止”而非终止——它的全部状态原样冻结，shell 把它挂进作业表；之后 \code{fg} 或 \code{bg} 让它收到 SIGCONT，从冻结处继续。这就是 shell 作业控制的全部硬件基础：同一套信号投递路径，不同的编号与默认动作。

还有两个工程细节补全这幅图。其一，标准信号的 pending 是位图，只能记“有”，记不了“几次”——处理函数运行期间再来的同号信号最多再记一位，这就是“信号不排队、可能合并丢失”的原因；需要计数的场合得用可排队的实时信号。其二，信号帧搭在用户栈上，若程序正因栈溢出而被 SIGSEGV 投递，用户栈已经不可用，处理函数需要一条事先用 \code{sigaltstack} 备好的备用栈——又一个“不能信任用户指针”的实例。

同构最漂亮的一处是段错误：经典芯片与平台案例那台检查 P/W/U 三个原因位的页错误硬件，在内核判定“故障地址不属于任何映射”之后，并不直接杀死进程，而是向它投递 SIGSEGV；程序崩溃时打印的“Segmentation fault”，就是这枚信号走完上述投递路径的结果。信号因此不需要任何新硬件——它全部用已有的异常返回路径搭出来，这正是“软件级中断”五个字的含义。

处理函数一侧还有一条对应的工程约束：信号可能在任意指令边界到来，处理函数里只能调用“异步信号安全”的函数，否则可能打断 \code{malloc} 到一半的链表——这与中断服务程序“快进快出、不碰半成品数据”的纪律同出一源。

最后注意一个归属问题：信号处理表和上一主题的页表一样，是每进程状态，挂在进程控制块上。上下文切换换走地址空间的同时，pending 位图和处理表也一起被换走——信号因此总是投递给“当前这个进程”的语义，而不会串到别的进程身上。

\topic{用户态与内核态的边界清单}
本章反复出现的那条线，现在可以写到底：边界不是软件风俗，而是特权级检查。页表项里的 U/S 位决定哪些页用户态可碰——经典芯片与平台案例页错误三问的第三问——检查由硬件在每次访存时完成，成本为零；\code{syscall}/\code{sysret} 是跨越特权级的唯一受控正门，中断门是另一扇，但钥匙在内核手里。判断一件事在哪一侧，只需问一句：它要不要碰特权状态？

\begin{longtable}{L{0.26\linewidth}L{0.12\linewidth}L{0.52\linewidth}}
\toprule
操作 & 在哪侧 & 为什么 \\
\midrule
加法、比较、分支 & 用户态 & 只动通用寄存器，不碰特权资源 \\\tablerowrule
\code{strlen}/\code{memcpy} 等库函数 & 用户态 & 只读写本进程已映射的页 \\\tablerowrule
\code{malloc} 小对象 & 用户态 & 在已分配的堆区内管理空闲块；堆要长大才 \code{brk}/\code{mmap} 进内核 \\\tablerowrule
\code{read}/\code{write}/\code{open} & 内核态 & 设备寄存器与 DMA 由内核独占（存储与设备事务相关章节） \\\tablerowrule
\code{fork}/\code{execve}/\code{exit} & 内核态 & 建/销地址空间要改页表与 CR3，属特权操作 \\\tablerowrule
\code{mmap}/映射变更 & 内核态 & 页表页标了 U/S=0，用户态写不了 \\\tablerowrule
\code{gettimeofday}（vDSO） & 用户态 & 内核把只读时间页映给用户，读它不算越界 \\
\bottomrule
\end{longtable}

清单最后一行值得展开。\code{gettimeofday} 传统上是系统调用，现代内核把它和校准参数写进一页只读共享页，用户态直接读——一次跨越都省。这是边界工程的一般做法，按优先级排成三条：能不跨就不跨（vDSO）；必须跨就合并（stdio 缓冲、\code{readv} 批量、\code{io\_uring} 一次提交一组）；非要高频跨越，就干脆修一座桥（共享页、门铃寄存器、用户态轮询）。

把跨越的代价并排摆出来，三条策略的动机就一目了然：

\begin{longtable}{L{0.26\linewidth}L{0.18\linewidth}L{0.46\linewidth}}
\toprule
动作 & 典型量级 & 边界含义 \\
\midrule
普通函数调用 & 纳秒级 & 不跨边界 \\\tablerowrule
系统调用往返 & 百纳秒级 & 陷入加返回各一次 \\\tablerowrule
进程上下文切换 & 微秒级 & 陷入加更换整套可见状态 \\\tablerowrule
SSD 读 4 KiB 页 & 几十微秒 & 内核加设备加缺页服务 \\
\bottomrule
\end{longtable}

\code{strace -c} 能把一个程序一生跨越边界的次数统计出来——清单对应的是上面那个最小程序：\code{write} 只出现一次；即使把它改成打印十次，当输出重定向到文件或管道时（全缓冲），\code{write} 通常也只出现一次（stdio 已把十次合并），其余是装载期的 \code{mmap} 和 \code{openat}。数一数这些行，就是把本节所有的账再对一遍：

\begin{lstlisting}
$ strace -c ./hello
hello, world
% time     seconds  usecs/call     calls    errors syscall
------ ----------- ----------- --------- --------- ----------------
 35.42    0.000017          17         1           execve
 18.75    0.000009           1         9           mmap
 16.67    0.000008           2         4         2 openat
 12.50    0.000006           6         1           write
  8.33    0.000004           1         3         1 read
  4.17    0.000002           1         2           close
  4.17    0.000002           2         1           arch_prctl
------ ----------- ----------- --------- --------- ----------------
100.00    0.000048                    21         3 total
\end{lstlisting}

边界同时是观察边界：\code{strace} 看不见函数调用，只能看见系统调用，因为陷入点是用户态唯一的“海关”。调试器、审计、沙箱全都挂在这同一个点上——拦住边界，就拦住了一个进程对外能做的全部事情；这也是为什么容器隔离的实施点是“过滤系统调用”，而不是“检查程序代码”。

边界的形态本身也是设计选择。宏内核把文件系统、驱动、协议栈都放在内核态，跨越成本是一次陷入；微内核把它们挪进用户态服务进程，跨越变成进程间消息——边界没有消失，只是挪了位置，而每次跨越的特权检查是物理事实，挪到哪里都要付。\code{time} 的输出把这笔开销记成两列：\code{user} 是用户态指令时间，\code{sys} 是内核代为执行的时间；\code{sys} 占比高的程序，多半就是边界跨越密集型。

\begin{lstlisting}
$ time ./hello
hello, world
real    0m0.002s
user    0m0.001s
sys     0m0.001s
\end{lstlisting}

本书末章的教学计算机项目的面包板机是对照的另一极：程序直接拨进 RAM，每个 MMIO 地址都均可直接访问，没有第二个特权级——那是“程序即内核”。多用户、多进程、要隔离要共享的系统必须有边界，而边界的全部成本，就是每次跨越时那一次受控陷入。

\begin{warningbox}
“库函数=系统调用”是最常见的混淆。\code{printf} 不是系统调用：它是用户态函数，把字符攒进缓冲区，只在缓冲满或换行刷新时才发起一次 \code{write}。用 \code{strace} 观察一个打印十次 \code{hello} 的程序：当输出重定向到文件或管道时（全缓冲），看不到十次 \code{write}——通常只有一次。判断依据不变：只有碰特权资源才需要陷入，拼字符串始终发生在用户态。
\end{warningbox}

\begin{classicnote}
CSAPP 把异常控制流分成四类：中断、陷阱、故障、中止。本章的系统调用是陷阱（trap）——主动、同步、返回下一条；定时器中断是中断——被动、异步；缺页是故障——可修复、返回重试；段错误走到尽头是中止。一张表管住四种入口，剩下的全是各入口里的具体约定。
\end{classicnote}

\section{五、从源码到进程的完整案例}

本节跟踪一个最小的 C 程序的一生——
从源码，到进程退出、shell 打印出退出码为止。程序故意
不用 \code{printf} 而直接调 \code{write}，这样它一生向内
核提出的请求最少，看得最清楚：

\begin{lstlisting}
#include <unistd.h>

int main(void)
{
        write(1, "hi\n", 3);
        return 42;
}
\end{lstlisting}

为了让地址看得见、对得上，本节统一用 \code{gcc -no-pie} 编
出传统的固定地址可执行文件；位置无关可执行（PIE）只是把同
一套机制整体搬到装载时才选定的基址上，步骤一步不少。

\topic{编译链的四个阶段}\code{gcc} 自己并不编译任何东西，
它是调度员：按固定顺序调用预处理器、编译器本体、汇编器和
链接器，每一道程序只负责一种表示到下一种表示的翻译。预处
理是文本层的翻译：宏展开、头文件拼接、条件编译裁剪，产物
仍然是 C——这个最小的 \code{hello.c} 经过它变成几千行的
\code{hello.i}，多出来的几乎全是头文件里的声明。编译是语
义层的翻译：词法语法分析、类型检查、优化，最后把 C 语句变
成 x86-64 汇编序列，产物还是文本。汇编是编码层的翻译：汇
编器把每条助记符翻成机器码字节，装进 ELF 可重定位目标文
件；从这一站起产物不再是文本，外部符号只有名字没有地址，
机器码里留着空位和重定位记录（“链接：重定位与地址形成”一节）待补。链接是定址层的
翻译：合并目标文件、解析符号、修补空位，再拼入启动文件与
动态链接器路径，产出内核能装载的可执行文件。

四把开关各让流水线停一站：

\begin{lstlisting}
$ gcc -E hello.c -o hello.i   # 停在预处理后：宏与头文件已展开
$ gcc -S hello.i -o hello.s   # 停在编译后：x86-64 汇编文本
$ gcc -c hello.s -o hello.o   # 停在汇编后：可重定位目标文件
$ gcc hello.o -o hello        # 走完全程：可执行文件
\end{lstlisting}

\begin{longtable}{L{0.09\linewidth}L{0.12\linewidth}L{0.14\linewidth}L{0.14\linewidth}L{0.35\linewidth}}
\toprule
阶段 & 停站开关 & 输入 & 输出 & 这一层消去的东西 \\
\midrule
预处理 & \code{-E} & \code{hello.c} & \code{hello.i} & 宏、条件编译、头文件嵌套 \\\tablerowrule
编译 & \code{-S} & \code{hello.i} & \code{hello.s} & 表达式、类型、控制流结构 \\\tablerowrule
汇编 & \code{-c} & \code{hello.s} & \code{hello.o} & 助记符，只剩机器码与符号名 \\\tablerowrule
链接 & 无（默认） & \code{hello.o} & \code{hello} & 一切未定的符号地址 \\
\bottomrule
\end{longtable}

编译一站的产物值得亲眼过目（\code{-masm=intel} 输出，删
去伪指令与标号细节）：

\begin{lstlisting}
main:
        push    rbp
        mov     rbp, rsp
        mov     edx, 3              ; 第三参数：长度
        lea     rsi, .LC0[rip]      ; 第二参数：字符串地址
        mov     edi, 1              ; 第一参数：文件描述符 1
        call    write@PLT
        mov     eax, 42             ; 返回值放进 eax
        pop     rbp
        ret
\end{lstlisting}

C 的三条基本约定在这七条指令里全部现形：参数按
\code{rdi}、\code{rsi}、\code{rdx} 依次传递，返回值放
\code{eax}，调用与返回靠 \code{call} 与 \code{ret}——正
是处理器核心相关章节讲的 ISA 约定在编译产物里的样子；\code{@PLT} 的
来历，要等“装载与动态链接”一节讲完延迟绑定再回头认。工程上，四把开关是四
扇观察窗：排错的第一问永远是“错在哪一层”。宏不对看
\code{.i}，代码生成有疑问看 \code{.s}，符号未定义是链接层
的事，装载期才报的错则与运行时找库的路径有关。错误发生在
哪一层，决定该打开哪扇窗。

\topic{\code{\_start} 与 \code{main} 之间的约定}内核装载
完可执行文件，把 \code{rip} 置为 ELF 头里的入口地址。这个
地址指向的不是 \code{main}——\code{main} 只是一个普通函
数，没有任何硬件地位。入口指向链接时由启动文件
（\code{crt1.o} 一族）提供的 \code{\_start}：

\begin{lstlisting}
$ readelf -h hello
ELF Header:
  Class:                             ELF64
  Data:                              2's complement, little endian
  Type:                              EXEC (Executable file)
  Machine:                           Advanced Micro Devices X86-64
  Entry point address:               0x401040
\end{lstlisting}

\code{\_start} 的处境比任何 C 函数都原始：它不是被谁“调
用”的，栈上没有返回地址，只有内核按约定摆好的
\code{argc}、\code{argv} 和环境变量。它的全部工作是把这堆
原料整理成 C 世界的样子（下为化简后的指令序列）：

\begin{lstlisting}
_start:
        xor     ebp, ebp            ; 帧指针清零：这是最外层帧
        pop     rsi                 ; 新栈顶第一项是 argc（内核在用户栈上摆好 argc/argv）
        mov     rdx, rsp            ; argv 数组紧随其后
        and     rsp, -16            ; 调用前 rsp 按 16 字节对齐
        mov     rdi, main           ; main 的地址作为第一参数
        call    __libc_start_main   ; 进 C 运行时，不再返回
        hlt                         ; 永远执行不到
\end{lstlisting}

三条细节值得停下来看。\code{and rsp, -16} 把栈指针低四位
清零，满足 ABI 的硬约定：\code{call} 执行的前一刻
\code{rsp} 必须是 16 的倍数，被调函数里才能安全使用要求
对齐的 SSE 指令。\code{pop rsi} 取出 \code{argc} 之后，
\code{rsp} 正好指向 \code{argv[0]}，于是一条
\code{mov rdx, rsp} 就得到了 \code{argv}。
\code{\_\_libc\_start\_main} 拿到 \code{main} 的地址后，
先初始化 C 运行时（环境变量、线程局部存储、栈保护值），再
按参数约定调用 \code{main}；\code{main} 返回后，它把返回
值交给 \code{exit}，由 \code{exit} 跑完 \code{atexit} 登
记的函数、冲洗标准库缓冲，最后发出 \code{exit\_group} 系
统调用。

所以“程序从 \code{main} 开始”是 C 运行时送的一份礼物：
控制到达 \code{main} 时，栈已对齐、参数已按约定摆好、退出
路径已挂好钩子。\code{return 42} 不是返回给内核，而是返
回给 \code{\_\_libc\_start\_main}，由它转译成
\code{exit\_group(42)}。退出码只有低 8 位会传给父进程：
本例在 shell 里查看 \codetext{\$?} 看到 42；若让 \code{main}
返回 256，看到的将是 0——256 的低 8 位全是零。本书末章的教学计算机项目的
裸机上没有这份礼物：复位向量后面直接就是你的代码，栈要自
己设，“退出”是一条停机指令，没有任何东西替你收尾。

\topic{strace 看一次进程出生}“系统调用与异常控制流”一节说过，系统调用是特殊的
异常入口。\code{strace} 借 \code{ptrace} 在每次陷入与返回
时拦一下，把编号翻成名字、把寄存器翻成参数——它打印的就
是进程一生向内核提出的全部请求的时间表。对这个最小程序，
时间表短到可以全列出来（模拟输出，地址与布局随机器而
异）：

\begin{lstlisting}
$ strace ./hello
execve("./hello", ["./hello"], 0x7ffd9e5a2400 /* 18 vars */) = 0
brk(NULL)                               = 0x1c56000
mmap(NULL, 8192, PROT_READ|PROT_WRITE, MAP_PRIVATE|MAP_ANONYMOUS, -1, 0) = 0x7ffff7fca000
access("/etc/ld.so.preload", R_OK)      = -1 ENOENT (No such file or directory)
openat(AT_FDCWD, "/etc/ld.so.cache", O_RDONLY|O_CLOEXEC) = 3
openat(AT_FDCWD, "/lib/x86_64-linux-gnu/libc.so.6", O_RDONLY|O_CLOEXEC) = 3
read(3, "\177ELF\2\1\1\3\0\0\0\0\0\0\0\0"..., 832) = 832
mmap(NULL, 2125824, PROT_READ, MAP_PRIVATE|MAP_DENYWRITE, 3, 0) = 0x7ffff7d00000
mprotect(0x7ffff7d28000, 1851392, PROT_NONE) = 0
close(3)                                = 0
write(1, "hi\n", 3)                     = 3
exit_group(42)                          = ?
+++ exited with 42 +++
\end{lstlisting}

真实输出里 \code{libc} 初始化还有几条调用，此处从略；剩下
的每一行都能和本章的机制对上号：

\begin{longtable}{L{0.17\linewidth}L{0.4\linewidth}L{0.29\linewidth}}
\toprule
系统调用 & 它在代表谁执行 & 对应本章机制 \\
\midrule
\code{execve} & 内核建地址空间、映射段、把控制交给动态链接器 &“装载与动态链接”一节装载；成功时永不返回 \\\tablerowrule
\code{brk}/\code{mmap} & 动态链接器登记堆起点、申请工作区 &“装载与动态链接”一节虚拟内存区与按需调页 \\\tablerowrule
\code{openat}/\code{read} & 动态链接器查缓存、找到并打开 \code{libc.so.6} &“装载与动态链接”一节动态库的查找路径 \\\tablerowrule
\code{mmap}/\code{mprotect} & 把 \code{libc} 各段映射进来、按段表修正权限 &“装载与动态链接”一节段映射与页权限 \\\tablerowrule
\code{write} & 你的 \code{main}：\code{rax} 为 1，经 \code{syscall} 陷入 &“系统调用与异常控制流”一节异常入口与参数约定 \\\tablerowrule
\code{exit\_group} & \code{exit} 的最后一站，注销整个进程 &“系统调用与异常控制流”一节与本节的退出路径 \\
\bottomrule
\end{longtable}

这张表里藏着两个反直觉的事实。其一，\code{execve} 是唯一
一个“成功就不回来”的系统调用：行尾的 \code{= 0} 发生在
一个崭新的地址空间里，原来那个调用者（shell 用
\code{fork} 复制出的副本）已经被整体替换——\code{fork}
返回两次，\code{execve} 一次也不返回，这对组合正是 shell
启动每条命令的标准动作。其二，\code{write} 之前那一串
\code{openat} 与 \code{mmap} 没有一条是你的代码发出的：
动态链接的可执行文件，第一条用户态指令属于动态链接器
\code{ld.so}；它跑完找库、映射、准备 GOT（“装载与动态链接”一节）之后，才
把控制交给 \code{\_start}，然后才轮到 \code{main}。

数字也都能闭合：\code{write} 返回 3，即实际写出的字节
数；\code{exit\_group} 行尾是 \code{= ?}，因为它根本不返
回，下一行 \code{+++ exited with 42 +++} 是 \code{strace}
自己补的旁白，42 正是 \code{main} 的返回值。从 C 的
\code{return} 到 shell 的退出码，整条链在一张表里走完。
若改用 \code{-static} 静态链接重跑一遍，\code{execve} 与
\code{write} 之间几乎一片空白——那串调用全是动态链接为
“共享库”三个字付的装载期成本，也是“装载与动态链接”一节说的装载时重定
位代价里最显眼的部分。

拿其中最短的一行做字段分解，strace 每行打印的几段内容都能指回陷入时的寄存器：

\begin{pdfdiagram}
  \node[hw label] at (5.6,2.2) {strace 一行 ＝ 调用名 ＋ 参数 ＋ 返回值（＋ 耗时）};
  \node[hw compute, minimum width=1.8cm, align=center] (f1) at (1.2,1.0) {\code{write}\\调用名\\由 rax 编号译出};
  \node[hw compute, minimum width=3.4cm, align=center] (f2) at (4.3,1.0) {\codetext{(1, "hi\textbackslash n", 3)}\\参数\\由 rdi/rsi/rdx 译出};
  \node[hw compute, minimum width=2.2cm, align=center] (f3) at (7.5,1.0) {\code{= 3}\\返回值\\由 rax 带回};
  \node[hw compute, minimum width=2.4cm, align=center] (f4) at (10.1,1.0) {\code{<0.000012>}\\耗时\\\code{-T} 选项追加};
  % 字段之间的衔接短横（非数据流，只表示同行相连）
  \draw[BookMuted, line width=0.62pt] (f1.east) -- (f2.west);
  \draw[BookMuted, line width=0.62pt] (f2.east) -- (f3.west);
  \draw[BookMuted, line width=0.62pt] (f3.east) -- (f4.west);
  \node[hw label, align=center] at (5.6,-0.4) {编号与寄存器值都被 ptrace 拦下后译成可读形式\\失败时返回值一段写作 \code{= -1 ENOENT} 这样的形式};
\end{pdfdiagram}
\diagramnote{典型 strace 行的字段分解（调用行与正文 \codetext{write(1, "hi\textbackslash n", 3) = 3} 一致）。\code{strace} 借 \code{ptrace} 在每次陷入与返回时拦一下：陷入时读 \code{rax} 译出调用名、读 \code{rdi}/\code{rsi}/\code{rdx} 译出参数，返回时读 \code{rax} 译出返回值——出错则印成 \code{= -1 ENOENT} 加一行说明；耗时一段只有加 \code{-T} 才出现在行尾（图中数值为示例）。一行的四段，全部来自“系统调用与异常控制流”一节那条 ABI 约定里的寄存器。}

读这张表还有几个顺手的工具用法。\code{strace -f} 会把
\code{fork} 出的子进程一并跟踪，shell 管道的每一环都能看
到；\code{strace -c} 不逐行打印，结束时给一张按耗时排序的
系统调用统计表，先找“哪类调用最多”再回头细看，是性能排
查的常规第一步；\code{strace -e trace=openat} 这类过滤开
关只保留关心的调用，用来对付输出以万行计的大程序。时间表
是本章所有机制的公共观察窗：装载、动态链接、系统调用、退
出清理，每一节的内容都能在这张表里找到对应的行。

\topic{程序退出时内核回收什么}\code{exit\_group}（编号
231）是进程的终点，也是看清“进程是什么”的最好窗口。对
用户来说，进程是一个正在运行的程序；对内核来说，进程是一
组挂在 \code{task\_struct} 上的数据结构：\code{mm\_struct}
描述地址空间（页表加虚拟内存区链表），
\code{files\_struct} 是文件描述符表，再加信号表、子进程链
表、时钟统计。退出，就是按固定顺序注销这组结构。

严格说，\code{exit\_group} 之前还有一层分工要辨清：
\code{exit} 是 C 库函数，负责用户态的收尾——按注册的逆
序跑 \code{atexit} 函数、冲洗 stdio 缓冲区；\code{\_exit}
才是薄包装，直接把控制交进内核。差别在缓冲上最明显：若
\code{main} 里用的是 \code{printf} 且输出不含换行，数据可
能还躺在用户态缓冲里，直接调 \code{\_exit} 会把这些字节
丢掉；\code{exit} 则先把它们冲刷出去，再走
\code{exit\_group}。本节程序用 \code{write} 不用
\code{printf}，也正是为了绕开这层缓冲，让输出与系统调用
一一对应。

地址空间最先退场：\code{mm\_struct} 的引用计数归零后，页
表与虚拟内存区链表被释放，匿名页和私有映射页回空闲链表；
共享映射（比如 \code{libc} 的代码页）只减一次引用计数，别
的进程还在用就一页不还——这正是“装载与动态链接”一节“多个进程共享一份
库代码”的镜像。文件描述符表随后：每个表项对内核文件对象
减一次计数，计数归零才真正关闭；这解释了 \code{fork} 之后
管道的经典现象——父子各持一份写端，必须两边都关闭，读端
才看得到 EOF。最后是亲属关系：退出者的子进程被过继给 1 号
进程，保证每个“孤儿”将来都有人为它收尸。

把这三项注销画成一张清单图，“进程是什么”就成了挂在 \code{task\_struct} 上的三笔：

\begin{pdfdiagram}
  \node[hw control, minimum width=2.8cm, align=center] (task) at (1.6,0.6) {task\_struct\\进程控制块};
  \node[hw state, minimum width=3.8cm, align=center] (i1) at (8.0,2.2) {地址空间 mm\_struct\\页表与映射释放\\共享页只减一次计数};
  \node[hw state, minimum width=3.8cm, align=center] (i2) at (8.0,0.6) {描述符表 files\_struct\\逐项减计数\\归零才真正关闭};
  \node[hw state, minimum width=3.8cm, align=center] (i3) at (8.0,-1.0) {子进程链表\\过继给 1 号进程};
  % 共享干线扇出到三项
  \draw[hw bus] (task.east) -- (4.6,0.6);
  \draw[hw bus] (4.6,0.6) -- (4.6,2.2) -- (i1.west);
  \draw[hw bus] (4.6,0.6) -- (i2.west);
  \draw[hw bus] (4.6,0.6) -- (4.6,-1.0) -- (i3.west);
  \fill[BookMuted] (4.6,0.6) circle (0.05);
  \node[hw label, fill=white, inner sep=1pt] at (4.6,1.6) {逐项注销};
  % 注销之后留下的骨架
  \node[hw state, minimum width=2.8cm, align=center] (zomb) at (1.6,-1.7) {僵尸态：仅剩\\退出码与统计\\等父进程 wait};
  \draw[hw return] (task.south) -- node[midway, right=1pt, hw label, fill=white, inner sep=1pt]{留骨架} (zomb.north);
\end{pdfdiagram}
\diagramnote{\code{exit\_group} 之后内核的三项注销。地址空间最先退场：\code{mm\_struct} 引用计数归零后页表与映射释放，共享映射只减一次计数——别的进程还在用就一页不还；描述符表随后：每个表项对内核文件对象减一次计数，归零才真正关闭；最后是亲属关系：退出者的子进程过继给 1 号进程。三项办完，\code{task\_struct} 仍留一条骨架记录——退出码还没交到父进程手里，进程以僵尸态等 \code{wait} 来取。}

但 \code{task\_struct} 本身不能立刻删——退出码还没交到
父进程手里。于是 \code{exit\_group} 之后进程进入僵尸态：
只剩一条骨架记录（退出码与资源统计），不占地址空间，永不
再被调度，等父进程用 \code{wait} 一族来取。父进程拿到的
是 32 位状态字，正常退出时退出码放在第 8 到 15 位：本例
的 42 对应原始状态值 42 左移 8 位，即 10752（十六进制
\code{0x2a00}），宏 \code{WEXITSTATUS} 做的就是右移 8 位
再取低 8 位。若父进程先死，过继来的子进程由 1 号进程负责
收取，系统里不会留下无人认领的僵尸；真正会留僵尸的情形只
有一种——父进程活着却忘了 \code{wait}，那是 \code{ps}
输出里状态列 \code{Z} 的来历。

“进程结束”因此不是一个硬件事件，而是内核数据结构的注销
顺序；僵尸也不是故障，是“退出码必须可靠传给父进程”这条
约定的必然代价。判断一个进程是否死透了，等价于问三个问
题：地址空间还了吗，描述符表清空了吗，最后那条
\code{task\_struct} 记录有人领走吗。

\topic{与本书末章的教学计算机项目裸机路径的对照}把本节的托管路径与本书末章的教学计算机项目的
裸机路径并排摆，会看到同一件事——“程序出生”——在两
个世界的两种做法。本书末章的教学计算机项目的机器上，链接地址在构建时定死：
复位向量的值就是程序的物理地址；装载器就是你自己——烧写
器把镜像写进 ROM，或拨码开关逐字节拨进 RAM；复位后的第一
条指令就是全部启动代码，没有 \code{crt}，更没有动态链接
器；输出一个字符，是往 I/O 端口或内存映射的显存写一个字
节，走存储与设备事务相关章节的 I/O 路径；“退出”是停机指令或原地跳转，
没有内核，自然也没有谁来回收什么。本节的世界相反：链接时
地址可以待定，装载由内核 \code{execve} 完成，\code{main}
之前已有 \code{ld.so} 和 \code{\_start} 跑了一整圈，输出
要经 \code{write} 请内核代劳，退出由内核逐项注销。

\begin{longtable}{L{0.14\linewidth}L{0.36\linewidth}L{0.36\linewidth}}
\toprule
维度 & 裸机镜像（本书末章的教学计算机项目） & 托管进程（本章） \\
\midrule
链接地址 & 构建时定死的固定地址 & \code{-no-pie} 时链接时定址（本章）；PIE 时只定相对偏移，基址由内核装载时定 \\\tablerowrule
装载者 & 烧写器或拨码开关，即你自己 & 内核 \code{execve} 加动态链接器 \\\tablerowrule
入口 & 复位向量直指你的第一条指令 & 入口指向 \code{\_start}，\code{main} 前还有一层 \\\tablerowrule
输出一个字符 & 直接写 I/O 端口或显存（存储与设备事务相关章节） & \code{write} 系统调用，内核驱动硬件 \\\tablerowrule
退出 & 停机或死循环，无人回收 & \code{exit\_group}，内核逐项注销数据结构 \\
\bottomrule
\end{longtable}

两个世界的分界线只有一条：有没有内核代为执行。有内核，地
址空间是借来的，文件是借来的，碰硬件要经“系统调用与异常控制流”一节的异常入口
申请；作为交换，装载、符号收尾、设备驱动、退出清理全有人
包办。没有内核，全部物理地址都归你，全部中断（存储与设备事务相关章节）也
都冲你来，自由与负担同比例放大。经典芯片与平台案例讲的启动链——固
件、引导程序、内核接力——正是从前一个世界走进后一个世
界的那座桥：引导程序活在裸机世界，它做的最后一件事是把内
核请进内存；内核一旦接管，桥的这边就改按托管世界的规矩
办。

把两条路径画成双列对照，中间那条虚线就是“有没有内核”的分界：

\begin{pdfdiagram}
  \node[hw label] at (2.6,2.9) {裸机路径（经典芯片与平台案例）};
  \node[hw label] at (9.1,2.9) {托管路径（本章）};
  \draw[BookMuted, line width=0.5pt, dashed] (5.85,-2.2) -- (5.85,2.55);
  \node[hw label] at (5.85,-2.55) {分界线：有没有内核};
  % 左列：裸机
  \node[hw compute, minimum width=4.0cm, align=center] (b1) at (2.6,2.0) {链接脚本限定地址\\入口＝复位向量};
  \node[hw compute, minimum width=4.0cm, align=center] (b2) at (2.6,0.9) {烧写器/拨码装载\\镜像进 ROM/RAM};
  \node[hw compute, minimum width=4.0cm, align=center] (b3) at (2.6,-0.2) {输出：直写 I/O 端口\\或显存（存储与设备事务相关章节）};
  \node[hw compute, minimum width=4.0cm, align=center] (b4) at (2.6,-1.3) {退出：停机或死循环\\无人回收};
  \draw[hw bus] (b1.south) -- (b2.north);
  \draw[hw bus] (b2.south) -- (b3.north);
  \draw[hw bus] (b3.south) -- (b4.north);
  % 右列：托管
  \node[hw state, minimum width=4.0cm, align=center] (h1) at (9.1,2.0) {链接定相对偏移\\基址装载时才定};
  \node[hw state, minimum width=4.0cm, align=center] (h2) at (9.1,0.9) {execve＋ld.so 装载\\建映射、填 GOT};
  \node[hw state, minimum width=4.0cm, align=center] (h3) at (9.1,-0.2) {输出：write 系统调用\\内核驱动硬件};
  \node[hw state, minimum width=4.0cm, align=center] (h4) at (9.1,-1.3) {退出：exit\_group\\内核逐项注销};
  \draw[hw bus] (h1.south) -- (h2.north);
  \draw[hw bus] (h2.south) -- (h3.north);
  \draw[hw bus] (h3.south) -- (h4.north);
\end{pdfdiagram}
\diagramnote{裸机路径与托管路径的双列对照（与正文表格同一内容）。同一行是同一件事的两种做法：定址、装载、输出、退出。左列没有内核，链接脚本把地址限定、装载靠自己、输出直写硬件、退出无人收尾；右列有内核，每一步都有人代办，代价是每次碰硬件都要经系统调用申请。两列的分界线只有一条：有没有内核代为执行。}

\begin{classicnote}
经典教材把 \code{fork}、\code{execve}、\code{exit} 与信号
一并归入“异常控制流”：它们和系统调用一样，都是控制从用
户态到内核态的异常转移。本节的 \code{strace} 输出可以读成
这种转移留下的时间表——每一行都是一次让渡与一次归还。
把程序的一生读成一张系统调用时间表，是理解进程最不容易过
时的方式。
\end{classicnote}




























\section{六、从目标文件、固件到可启动机器}

\topic{目标文件里的符号最终要变成地址}
汇编器或编译器产生的目标文件通常还不知道所有最终地址。函数名、全局变量名、外部符号和跳转目标先以符号或重定位项表示。链接器把多个目标文件和库合并，决定代码段、只读数据段、数据段、BSS 段等放到什么地址，并把需要修补的位置改成最终地址或相对偏移。

\begin{longtable}{L{0.18\linewidth}L{0.38\linewidth}L{0.34\linewidth}}
\toprule
对象 & 系统含义 & 硬件落点 \\
\midrule
代码段 & 机器指令编码 & 取指路径从这些地址读取 opcode \\\tablerowrule
只读数据 & 常量、跳转表、字符串 & 普通加载读取，权限通常禁止写 \\\tablerowrule
数据段 & 已初始化全局变量 & 启动时从镜像复制到 RAM \\\tablerowrule
BSS & 零初始化全局变量 & 启动时清零，镜像中不必保存全量零字节 \\\tablerowrule
重定位项 & 需要链接器或装载器修补的地址位置 & 改写指令立即数字段、地址表或指针字 \\
\bottomrule
\end{longtable}

这说明“函数地址”不是抽象标签。取指硬件需要 PC 中的数值地址；CALL 需要目标地址或相对偏移；全局变量加载需要数据地址；中断向量需要入口地址。链接器和装载器的工作，就是把这些名字和段布局变成 CPU 能取指、加载、存储和跳转的地址。

\begin{lstlisting}
source idea:
  call uart_putc
  load global_counter

after link/load:
  CALL 0x00102480
  MOV R1, [0x00203010]
\end{lstlisting}

\topic{装载器建立的是第一张可执行地址地图}
在有操作系统的机器上，装载器把可执行文件映射进进程地址空间，建立代码、数据、堆、栈和共享库区域；在裸机或固件场景中，启动代码负责把镜像从 ROM/Flash 复制或映射到合适位置，清 BSS，设置栈指针，初始化全局数据，然后跳到入口函数。两者规模不同，但硬件问题一致：哪些地址可取指，哪些地址可读写，栈在哪里，异常入口在哪里，设备寄存器在哪里。

\begin{longtable}{L{0.18\linewidth}L{0.42\linewidth}L{0.3\linewidth}}
\toprule
启动/装载动作 & 硬件意义 & 失败表现 \\
\midrule
设置栈指针 & CALL/RET、局部变量和异常入口有可用内存 & 第一次函数调用或中断即破坏内存 \\\tablerowrule
复制数据段 & RAM 中全局变量得到初始值 & 程序读到未初始化数据 \\\tablerowrule
清 BSS & 零初始化对象符合语言和固件约定 & 状态机、计数器、指针初值随机 \\\tablerowrule
建立向量表 & 中断和异常有入口地址 & 外部事件或错误无法处理 \\\tablerowrule
初始化时钟/总线 & 后续 RAM、设备、cache 路径满足时序 & 访问慢器件超时或读写失败 \\\tablerowrule
跳入口 & PC 改到主程序第一条指令 & 程序真正进入可维护代码路径 \\
\bottomrule
\end{longtable}

裸机启动常见伪代码如下。它表面像软件初始化，实质上每一步都在建立后续硬件事务的前提。

\begin{lstlisting}
// 注释（推进）：按步骤依次执行，前一步的稳定结果是后一步的前提。
// 注释（边界）：每一步都保留可观察读回；不满足预期就停在当前层排错。
reset_handler:
  set SP to top_of_stack
  copy .data from flash image to SRAM
  zero .bss in SRAM
  initialize clocks and memory controller
  install vector table base
  initialize essential MMIO devices
  jump main
\end{lstlisting}

\topic{异常向量表是硬件失败路径的地址表}
只要系统会处理中断、页错误、非法指令或设备事件，就需要某种入口表。早期处理器可能用固定地址或固定向量；保护模式处理器使用更复杂的描述符表；微控制器常在启动区域放置向量表。表的形式不同，但作用相同：当硬件事件发生时，CPU 能把事件编号转成处理程序入口，并保存足够状态以便恢复或诊断。

\begin{longtable}{L{0.18\linewidth}L{0.36\linewidth}L{0.36\linewidth}}
\toprule
向量项 & 保存内容 & 硬件使用场景 \\
\midrule
复位向量 & 第一条启动代码地址或跳转指令 & 上电或复位后装入 PC \\\tablerowrule
非法指令向量 & 无法译码或不允许执行时的入口 & opcode 不合法或特权不允许 \\\tablerowrule
页错误/总线错误向量 & 地址转换或外部事务失败入口 & 加载/存储/取指不能完成 \\\tablerowrule
外设中断向量 & 设备完成、错误或数据到达入口 & 中断控制器投递事件 \\\tablerowrule
系统调用入口 & 用户代码请求特权服务入口 & 权限切换和受控进入内核或监控程序 \\
\bottomrule
\end{longtable}

向量表本身也必须位于可访问地址。若页表或地址译码没有覆盖异常入口，CPU 在处理第一个错误时会遇到第二个错误；若栈没有准备好，保存 PC/Flags 或错误码时会破坏未知内存；若中断处理程序所在代码被错误标成不可执行，事件到达后系统会在入口处再次失败。启动顺序必须先保证最小异常路径可靠，再打开复杂设备和中断。

\topic{页表把装载结果提升为可保护的运行时地图}
链接器决定镜像中的相对布局，装载器把段放进内存，页表则把运行时地址空间变成可检查的约定。代码页可以标为只读和可执行，数据页可读写但不可执行，内核页禁止用户访问，设备 MMIO 页设置为不可缓存或强顺序。这样，同一条加载/存储/取指路径在真正访问 cache 或总线前，会先经过地址转换和权限检查。

\begin{longtable}{L{0.2\linewidth}L{0.36\linewidth}L{0.34\linewidth}}
\toprule
区域 & 典型页属性 & 硬件效果 \\
\midrule
代码 & 只读、可执行 & 对代码页的写操作触发保护异常 \\\tablerowrule
只读数据 & 只读、不可执行或按平台规定 & 防止常量被误写或执行数据 \\\tablerowrule
普通数据/堆/栈 & 可读写、通常不可执行 & 数据写入允许，取指可被阻止 \\\tablerowrule
内核映射 & 特权访问，用户禁止 & 用户态越权访问触发异常 \\\tablerowrule
MMIO 窗口 & 不可缓存、设备顺序属性 & 读写按设备事务执行，不被普通 cache 延迟或合并 \\
\bottomrule
\end{longtable}

现代系统的许多安全和可靠性机制，本质上就是让硬件在每次取指和访存前检查这些约定。

\topic{从固件到操作系统是平台所有权转移}
固件启动阶段拥有硬件控制权：它配置时钟、内存控制器、基本 I/O、启动介质和必要安全检查。操作系统接管后，会重新建立页表、中断控制、设备驱动和调度机制。这个交接是一段平台所有权转移：谁拥有中断表，谁拥有页表，谁管理内存，谁接管设备，谁定义错误处理策略。

\begin{longtable}{L{0.18\linewidth}L{0.4\linewidth}L{0.32\linewidth}}
\toprule
阶段 & 固件职责 & 交给后续系统的约定 \\
\midrule
早期启动 & 电源、时钟、复位、最小取指路径 & CPU 能从可信入口执行 \\\tablerowrule
内存初始化 & DRAM 控制器、cache/TLB 初始策略 & 可用内存范围和保留区域 \\\tablerowrule
设备发现 & 基本总线枚举、启动介质选择 & 设备地址窗口、中断和 DMA 能力 \\\tablerowrule
镜像验证/加载 & 读取、校验、解压或复制内核/程序 & 入口地址、参数和内存布局 \\\tablerowrule
控制权转移 & 设置寄存器、栈、页表或模式后跳转 & 后续系统接管异常、中断和设备 \\
\bottomrule
\end{longtable}

在微控制器上，固件和“应用”可能就是同一个镜像；在 PC 或服务器上，固件、引导加载器、内核和用户程序分层明显；在 SoC 上，还可能有 Boot ROM、一级引导、可信固件、安全监控和普通操作系统。层数不同，但每一层都要把可执行地址、栈、异常入口、内存属性和设备所有权交代清楚。

\topic{安全启动把第一条指令也纳入信任链}
现代 SoC 常把 Boot ROM 固化在芯片内部，复位后先执行不可修改的 ROM 代码。Boot ROM 读取下一阶段固件，验证签名、版本和配置，再决定是否跳转。这样做是硬件启动约定的一部分：如果第一段可变固件不可信，后续页表、设备配置、密钥和操作系统都可能被篡改。

\begin{longtable}{L{0.18\linewidth}L{0.4\linewidth}L{0.32\linewidth}}
\toprule
安全启动对象 & 硬件/固件动作 & 失败处理 \\
\midrule
Boot ROM & 从固定复位入口执行，不可现场改写 & 提供根信任和最小验证代码 \\\tablerowrule
公钥/哈希 & 存在熔丝、ROM 或安全存储中 & 用于验证下一阶段镜像 \\\tablerowrule
镜像签名 & 覆盖代码、配置和版本信息 & 验证失败则停止、降级或进入恢复模式 \\\tablerowrule
回滚保护 & 记录允许的最小版本 & 防止加载旧漏洞固件 \\\tablerowrule
调试锁定 & 限制 JTAG、串口下载或内存读出 & 防止启动后被外部接口改写控制流 \\
\bottomrule
\end{longtable}

整机一旦能取指、访存和访问设备，就必须回答谁能改代码、谁能访问内存、谁能发起 DMA、谁能接收中断、启动镜像是否可信、错误路径是否可恢复。

\topic{案例：一个裸机镜像从链接地址跑到 main}
假设一块教学板把 Flash 映射到 \code{0x00000000}，SRAM 映射到 \code{0x20000000}，GPIO 映射到 \code{0x40000000}。链接脚本规定向量表和代码放在 Flash，已初始化数据运行在 SRAM，但初始内容存放在 Flash 镜像中，BSS 在 SRAM 中清零，栈顶放在 SRAM 高地址。这个布局看起来是软件文件，实际上直接决定复位后每一次取指、加载和存储的地址。

\begin{lstlisting}
memory map:
  0x00000000-0x0003FFFF  Flash: vectors, code, const, data image
  0x20000000-0x20007FFF  SRAM: data, bss, stack
  0x40000000-0x40000FFF  GPIO MMIO
\end{lstlisting}

\begin{longtable}{L{0.18\linewidth}L{0.38\linewidth}L{0.34\linewidth}}
\toprule
镜像区域 & 链接/装载含义 & 硬件事务 \\
\midrule
向量表 & 复位入口和异常入口地址 & 复位时取入口，异常时取向量 \\\tablerowrule
\code{.text} & 指令编码放在 Flash & PC 取指访问 Flash 或映射后的代码区 \\\tablerowrule
\code{.rodata} & 常量和只读表 & 加载读取，通常禁止写 \\\tablerowrule
\code{.data} 镜像 & 初始值存放在 Flash & 启动代码复制到 SRAM 运行地址 \\\tablerowrule
\code{.bss} & 运行时应为 0 的对象 & 启动代码对 SRAM 执行写操作清零 \\\tablerowrule
栈 & 函数调用和异常保存上下文 & SP 指向 SRAM 可写区域 \\
\bottomrule
\end{longtable}

复位后第一段代码不应依赖尚未初始化的全局变量，也不应调用需要完整栈和运行时的复杂函数。它要先建立最小可运行边界。

\begin{lstlisting}
// 注释（推进）：按行追踪数据来源、经过的部件和最终写入目标。
// 注释（边界）：只有标出的时钟沿、阶段或异常入口会改变体系结构可见状态。
reset trace:
  PC = vector_table.reset_handler
  SP = _stack_top
  copy bytes from _sidata in Flash to _sdata in SRAM
  zero bytes from _sbss to _ebss
  configure clocks and essential wait states
  set vector table base if architecture supports it
  call main
\end{lstlisting}

\begin{longtable}{L{0.16\linewidth}L{0.42\linewidth}L{0.32\linewidth}}
\toprule
启动动作 & 不能省略的原因 & 常见错误 \\
\midrule
设置 SP & 后续 CALL、PUSH、异常保存需要栈 & SP 指到 Flash 或未映射区 \\\tablerowrule
复制 \code{.data} & 全局初值要出现在 RAM 运行地址 & 变量初值全错或仍为 Flash 内容 \\\tablerowrule
清 \code{.bss} & 语言和固件约定零初始化 & 状态机初值随机，指针非零 \\\tablerowrule
配置等待状态 & Flash、SRAM 或外设满足时序 & 高频下取指或访问偶发失败 \\\tablerowrule
设置向量表 & 中断和异常能进入正确入口 & 第一场中断跳到错误地址 \\\tablerowrule
跳 \code{main} & 进入应用主逻辑 & 入口地址或模式错误导致程序失控 \\
\bottomrule
\end{longtable}

\topic{案例：保护模式最小启动清单}
以 80386 风格机器为例，若要从实模式进入带分页的保护环境，必须先建立一组硬件可读的表和寄存器状态。这个过程常被写成启动代码技巧，但本质是把 CPU 的地址解释约定从旧模式切换到新模式。

\begin{longtable}{L{0.16\linewidth}L{0.42\linewidth}L{0.32\linewidth}}
\toprule
步骤 & 硬件目的 & 预期行为 \\
\midrule
建立 GDT & 给代码段、数据段、栈段描述符 & 描述符 base/limit/type/present 正确 \\\tablerowrule
建立 IDT & 给异常和中断入口描述符 & 页错误、非法指令等入口可达 \\\tablerowrule
加载 GDTR/IDTR & CPU 知道描述符表位置 & 表所在内存可读，界限正确 \\\tablerowrule
打开保护模式位 & 改变段选择子解释方式 & 远跳转后 CS 使用新描述符 \\\tablerowrule
重装段寄存器 & DS/SS/ES 等进入新约定 & 数据和栈访问按新描述符解释 \\\tablerowrule
建立页目录/页表 & 线性地址可转换到物理地址 & 取指、栈、异常代码均被映射 \\\tablerowrule
加载 CR3 并开启分页 & CPU 开始查页表和 TLB & 缺失映射会触发页错误而非随机访问 \\
\bottomrule
\end{longtable}

最小保护模式 bring-up 应先保守映射。把当前代码、栈、页表、GDT、IDT 和 VGA/串口调试 MMIO 映射好，再打开分页；确认能打印或点亮调试输出后，再逐步收紧权限。若一开始就建立复杂高半区映射、多个页大小、用户/内核分离和设备属性，任何一个地址错误都可能导致连续异常，难以定位。

\begin{lstlisting}
// 注释（推进）：按步骤依次执行，前一步的稳定结果是后一步的前提。
// 注释（边界）：每一步都保留可观察读回；不满足预期就停在当前层排错。
protected-mode checklist:
  GDT covers code/data/stack selectors
  IDT covers at least fault handlers
  stack is writable after mode switch
  current instruction stream is mapped
  page tables themselves are mapped
  debug output MMIO is mapped uncached
  fault handlers do not fault recursively
\end{lstlisting}

\topic{案例：从异常号定位启动失败}
启动阶段最可怕的现象是“黑屏”或“无输出”。更有效的做法是把失败分层。若能进入非法指令异常，说明取指至少到了某条不合法编码；若能进入页错误，说明保护模式和 IDT 至少部分工作；若连续进入双重错误或复位，说明异常处理路径本身可能失败；若完全没有总线活动，可能是复位、时钟或复位向量映射问题。

\begin{longtable}{L{0.22\linewidth}L{0.38\linewidth}L{0.3\linewidth}}
\toprule
现象 & 可能边界 & 优先观察 \\
\midrule
无取指总线活动 & 复位、时钟、复位向量、ROM 片选 & 示波器看时钟/复位/地址线 \\\tablerowrule
取指后立即异常 & 编码、模式、段描述符或权限错误 & IR/PC、异常号、错误码 \\\tablerowrule
页错误循环 & 页表缺映射、权限错误、异常入口未映射 & 故障地址、PDE/PTE、CR3 \\\tablerowrule
中断后死机 & 栈、IDT、EOI 或清中断顺序错误 & SP、向量号、中断控制器状态 \\\tablerowrule
DMA 后数据错 & cache 可见性、IOMMU、缓冲区所有权 & 描述符、completion、cache 维护记录 \\
\bottomrule
\end{longtable}

\topic{案例：把教学板扩展成一台小系统}
最后给一个可以落地的整机扩展路线。第一阶段只跑 ROM 中的 LED 循环；第二阶段加入 SRAM 和栈，能调用函数；第三阶段加入定时器中断，能周期性切换状态；第四阶段加入简化 DMA，把一段内存复制到另一段；第五阶段加入保护或至少地址错误检测，让未映射访问进入错误入口。每一阶段都对应书中一个硬件层次。

\begin{longtable}{L{0.16\linewidth}L{0.42\linewidth}L{0.32\linewidth}}
\toprule
阶段 & 新增硬件/固件 & 预期行为 \\
\midrule
ROM LED & 复位向量、ROM、GPIO MMIO & 上电后固定节奏改写 LED 寄存器 \\\tablerowrule
SRAM 栈 & SRAM 译码、SP 初始化、CALL/RET & 函数调用和局部变量不破坏全局数据 \\\tablerowrule
定时器中断 & 定时器、向量表、中断清除 & 周期事件不丢失，主循环可继续运行 \\\tablerowrule
简化 DMA & 描述符、源/目标地址、completion & DMA 完成后 CPU 读到正确内存内容 \\\tablerowrule
错误入口 & 未映射检测、总线超时或简化异常 & 错误地址进入可观察处理程序 \\
\bottomrule
\end{longtable}

\topic{综合项目：写一份最小整机规格}
规格不需要一开始很复杂，但必须让别人能按它搭出同一台机器，并能判断每个访问是成功、等待还是错误。下面给出一个建议模板，读者可以用分立逻辑、FPGA、微控制器外设模拟或软件仿真来实现。

\begin{longtable}{L{0.18\linewidth}L{0.42\linewidth}L{0.3\linewidth}}
\toprule
规格项目 & 必须写清楚 & 预期行为 \\
\midrule
地址空间 & ROM、RAM、MMIO、未映射区范围 & 任意地址最多一个目标响应 \\\tablerowrule
复位行为 & PC 初值、SP 初值、控制线安全值 & 上电第一条取指可重复观察 \\\tablerowrule
指令集 & 编码、寄存器字段、立即数扩展和标志规则 & 每条指令能写成控制 trace \\\tablerowrule
内存访问 & 字节序、对齐、字节使能、等待和错误 & 加载/存储不破坏相邻字节 \\\tablerowrule
设备寄存器 & 每个位的读写、副作用和清除规则 & GPIO/定时器/UART 状态可验证 \\\tablerowrule
中断/异常 & 向量入口、保存状态、清除和返回规则 & 外部事件和错误能进入可诊断路径 \\\tablerowrule
DMA/队列 & 描述符格式、所有权、completion 和 cache 规则 & CPU 与设备不同时拥有缓冲区 \\
\bottomrule
\end{longtable}

规格写完后，建议按下面的顺序实现，而不是同时接上所有模块。

\begin{longtable}{L{0.14\linewidth}L{0.42\linewidth}L{0.34\linewidth}}
\toprule
里程碑 & 要实现的最小能力 & 预期行为 \\
\midrule
M1 & 复位后从 ROM 取指，PC 顺序前进 & 逻辑分析仪能看到连续取指地址 \\\tablerowrule
M2 & ADD/SUB/JNZ 正确执行计数循环 & 寄存器值和分支方向符合 trace \\\tablerowrule
M3 & SRAM 加载/存储和栈可用 & 函数调用能返回，局部变量不乱 \\\tablerowrule
M4 & GPIO MMIO 可读写 & LED/按键通过地址事务控制 \\\tablerowrule
M5 & 定时器中断可进入和返回 & 中断不破坏主循环寄存器和栈 \\\tablerowrule
M6 & 简化 DMA 或块复制控制器可搬数据 & completion 后内存内容正确 \\\tablerowrule
M7 & 未映射访问进入错误入口 & 错误地址和原因可观察 \\
\bottomrule
\end{longtable}

每个里程碑都应保留一个最小测试程序。不要只保留最终大程序，因为最终程序一旦失败，无法判断是取指、译码、ALU、存储器、MMIO、中断还是 DMA 出错。小测试程序就是硬件书里的单元测试。

\begin{lstlisting}
// 注释（推进）：每轮只推进一个元素或一个控制步，并保持中间状态的一致性。
// 注释（边界）：退出条件成立后才提交结果；空集合、越界和失败要单独处理。
M2 counter loop:
  MOV  R1, 0
  MOV  R2, 10
loop:
  ADD  R1, 1
  SUB  R2, 1
  JNZ  R2, loop
  MOV [GPIO_OUT], R1
\end{lstlisting}

\begin{longtable}{L{0.2\linewidth}L{0.38\linewidth}L{0.32\linewidth}}
\toprule
测试失败现象 & 优先怀疑 & 观察信号 \\
\midrule
PC 不变 & 复位、时钟、PC 写使能或取指等待 & \code{PC}、\code{reset}、\code{clock}、\code{ready} \\\tablerowrule
循环次数错 & 立即数扩展、SUB、Zero 标志或 JNZ & \code{R2}、\code{Zero}、\code{pc_sel}、\code{branch_target} \\\tablerowrule
GPIO 不变 & 地址译码、MMIO 写使能、字节使能 & \code{address}、\code{io_cs}、\code{mem_write}、\code{write_data} \\\tablerowrule
函数返回错 & SP、CALL/RET 返回地址、栈 RAM & \code{SP}、栈内容、\code{PC}、\code{writeback} \\\tablerowrule
中断后错乱 & 保存/恢复寄存器、EOI、清状态顺序 & \code{vector}、\code{SP}、saved PC、irq status \\\tablerowrule
DMA 数据错 & 描述符、地址权限、cache 可见性、completion & \code{desc}、\code{buffer}、\code{owner}、\code{done} \\
\bottomrule
\end{longtable}

\topic{综合项目的实验报告模板}
为了避免实验结果只停留在“好像能跑”，每个里程碑都应留下实验报告。报告不需要长篇叙事，但必须保存足够的观察记录，让另一个人能复现同样结论。建议每个实验至少记录以下字段：

\begin{longtable}{L{0.18\linewidth}L{0.42\linewidth}L{0.3\linewidth}}
\toprule
报告字段 & 要记录的内容 & 为什么重要 \\
\midrule
硬件版本 & 原理图版本、器件型号、时钟频率、电源电压 & 不同器件和频率会改变时序结论 \\\tablerowrule
固件版本 & ROM 镜像、链接地址、入口地址、编译选项 & 同一硬件跑不同镜像可能表现不同 \\\tablerowrule
地址地图 & ROM/RAM/MMIO/未映射区域范围 & 后续访问错误可直接对照 \\\tablerowrule
测试程序 & 最小汇编或机器码列表 & 能定位具体指令和控制 trace \\\tablerowrule
观察信号 & PC、IR、控制字、地址、数据、ready、写使能 & 证明结果来自正确硬件路径 \\\tablerowrule
预期结果 & 预期寄存器、内存、设备和异常状态 & 避免把偶然现象当成成功 \\\tablerowrule
失败记录 & 失败现象、假设、修改和复测结果 & 保留完整调试记录 \\
\bottomrule
\end{longtable}

例如 M3“SRAM 栈可用”的实验报告应包含：复位后 SP 的值，CALL 前后的 PC，栈上保存的返回地址，函数内部局部变量地址，RET 后 PC 是否回到调用点，R0 或指定返回值寄存器是否正确。若只写“函数调用成功”，后续出现异常时无法判断是栈、返回地址、寄存器保存还是取指路径出了问题。

\begin{lstlisting}
// 注释（推进）：按行追踪数据来源、经过的部件和最终写入目标。
// 注释（边界）：只有标出的时钟沿、阶段或异常入口会改变体系结构可见状态。
M3 trace example:
  reset PC = 0x00000000
  initial SP = 0x20008000
  CALL writes return address 0x00000124 to [SP-4]
  callee uses [SP-8] for local variable
  RET loads PC = 0x00000124
  R0 = expected return value
\end{lstlisting}

硬件工程的核心是每次改动后仍能解释：哪个状态边界改变了，哪个事务完成了，哪个错误路径被覆盖了，哪些观察结果足以支持结论。

\topic{专题：从 8086 到 SoC 的“同一问题不同边界”}

\begin{longtable}{L{0.18\linewidth}L{0.34\linewidth}L{0.38\linewidth}}
\toprule
平台问题 & 早期芯片形态 & 现代 SoC 形态 \\
\midrule
复位入口 & 外部 ROM 响应固定复位向量 & 片内 Boot ROM、熔丝、启动介质选择和安全检查 \\\tablerowrule
地址译码 & 板级译码器产生片选 & 片上互连、地址防火墙、IOMMU 和外设桥 \\\tablerowrule
等待状态 & READY/WAIT 拉长总线周期 & valid/ready、AXI response、QoS 和超时 \\\tablerowrule
中断 & 外部中断控制器或简单向量 & GIC/APIC、优先级、亲和性、虚拟化中断 \\\tablerowrule
DMA & 简单外设总线主控 & 多主设备、cache 一致性、SMMU/IOMMU 权限 \\\tablerowrule
调试 & 逻辑分析仪看总线引脚 & JTAG/SWD、trace macrocell、性能计数器和固件日志 \\
\bottomrule
\end{longtable}

\topic{专题：BIOS、Boot ROM 和 loader 的最小移交清单}
启动链路不是“跳到下一段代码”这么简单。每一级启动代码都必须向下一级交出可用的执行环境：当前特权级、栈、内存可用范围、异常入口、时钟、串口或日志设备、镜像位置和校验结果。缺少任何一项，下一阶段可能运行一段时间后才以难以定位的方式失败。

\begin{longtable}{L{0.2\linewidth}L{0.42\linewidth}L{0.28\linewidth}}
\toprule
移交项目 & 必须说明的内容 & 失败现象 \\
\midrule
入口地址 & 下一阶段第一条指令位置和执行状态 & 跳入错误模式或错误对齐地址 \\\tablerowrule
栈 & 栈顶、增长方向、可用范围 & CALL/异常保存破坏数据或 ROM 区 \\\tablerowrule
内存地图 & RAM、ROM、MMIO、保留区、安全区 & loader 覆盖设备寄存器或固件数据 \\\tablerowrule
异常向量 & 向量基址、异常栈、默认处理器 & 早期错误无法诊断，直接死机 \\\tablerowrule
时钟/电源 & CPU、总线、外设时钟是否打开 & 访问外设永久等待或超时 \\\tablerowrule
日志通道 & UART、半主机、内存日志或调试口 & 失败时没有可观察的输出 \\\tablerowrule
镜像验证 & 长度、哈希、签名、加载地址 & 执行损坏或不可信代码 \\
\bottomrule
\end{longtable}

安全启动只是把这张移交清单加上信任约束：Boot ROM 信任芯片内固化公钥或熔丝状态，验证下一阶段签名，再把控制权交给经过验证的镜像。若验证失败，平台不能继续执行任意代码，而应进入恢复、锁定或错误报告路径。

\topic{案例：一次 \code{int 13h} 磁盘读取的流程}
实模式下的 PC 固件不只负责启动，还向上层软件提供一组设备服务。磁盘读写服务挂在中断号 \code{0x13} 上：调用者把功能号和参数放进寄存器，执行 \code{int 0x13}，BIOS 中的磁盘服务例程接管 CPU，完成与磁盘控制器的全部交互，再把结果按约定交还。以经典的 CHS 读扇区为例，调用约定如下表。

\begin{longtable}{L{0.12\linewidth}L{0.36\linewidth}L{0.38\linewidth}}
\toprule
寄存器 & 调用时放入的值 & 补充约定 \\
\midrule
\code{AH} & 功能号，\code{0x02} 表示读扇区 & 返回时变为状态码，0 表示成功 \\\tablerowrule
\code{AL} & 要读的扇区数 & 一次调用不宜越过当前磁道剩余扇区 \\\tablerowrule
\code{CH} & 柱面号低 8 位 & 10-bit 柱面号的高 2 位挤在 \code{CL} 高 2 位 \\\tablerowrule
\code{CL} & 低 6 位为扇区号 & 扇区从 1 开始编号，没有 0 号扇区 \\\tablerowrule
\code{DH} & 磁头号 & 双面软盘取 0 或 1 \\\tablerowrule
\code{DL} & 驱动器号 & \code{0x00} 为第一张软盘，\code{0x80} 为第一块硬盘 \\\tablerowrule
\code{ES:BX} & 缓冲区段地址与偏移 & 必须落在 1 MiB 实模式空间内，软盘 DMA 时不可跨越 64 KiB 边界 \\
\bottomrule
\end{longtable}

一次完整的调用写成汇编是这样的。

\begin{lstlisting}
; read one sector from the first hard disk, CHS = (0, 0, 1)
; buffer at ES:BX = 0x0000:0x7E00
mov ah, 0x02        ; function: read sectors
mov al, 1           ; sector count
mov ch, 0           ; cylinder low 8 bits
mov cl, 1           ; sector 1, cylinder high bits 0
mov dh, 0           ; head 0
mov dl, 0x80        ; first hard disk
mov bx, 0x7E00      ; buffer offset, ES already 0
int 0x13
jc  disk_error      ; CF=1 means failure, AH holds status code
\end{lstlisting}

\code{int} 指令把标志和返回地址压栈、按向量表跳到 BIOS 例程之后，剩下的工作由固件完成：它把 CHS 地址翻译成磁盘控制器的命令序列，对软盘借助 DMA 控制器把 512 字节直接搬进内存，对硬盘按当时平台的能力用 PIO 或总线主控方式传输；传送结束后，BIOS 检查控制器状态寄存器，把结果写成“进位标志 CF 为 0 表示成功、为 1 表示失败且 \code{AH} 保存状态码”的约定，再用 \code{iret} 返回调用者。对调用者而言，整件事只表现为一条中断指令和一段被填好的缓冲区：控制器端口地址、DMA 通道编程、出错重试全部藏在 ROM 里。

这正是“固件提供设备驱动”的早期形态：硬件细节被封装在一组寄存器约定接口后面，启动扇区代码不必知道磁盘控制器长什么样也能读盘。代价同样明显：BIOS 例程工作在实模式，缓冲区被限制在 1 MiB 以下的常规内存，传输方式以当时的 DMA 和 PIO 为上限，接口能力跟不上后来的大容量硬盘和高速控制器。操作系统接管平台后，会用自带的驱动替换这层固件服务——内核直接编程磁盘控制器，从 IDE、AHCI 一路演进到 NVMe，用自己的缓冲区和队列调度 I/O；\code{int 0x13} 只在启动早期、内核驱动就绪之前承担读盘任务。这条替换关系把“从固件到操作系统的所有权转移”具体到了每一个设备服务入口。

\topic{专题：平台 bring-up 日志怎样定位到具体层}
bring-up 的第一原则是每一步只引入一个新假设，并记录足够的观察结果。若一开始就运行完整系统，失败时无法知道是时钟、复位、内存、cache、异常、串口、设备树、驱动还是文件系统出错。更稳妥的办法是从复位向量开始逐层扩大可用机器。

\begin{longtable}{L{0.16\linewidth}L{0.42\linewidth}L{0.32\linewidth}}
\toprule
阶段 & 建议日志 & 证明的边界 \\
\midrule
复位 & PC、模式、第一条指令字 & ROM 映射和取指可用 \\\tablerowrule
串口 & 时钟源、波特率、发送寄存器状态 & MMIO、外设时钟和基本输出可用 \\\tablerowrule
SRAM & 写读 walking bit、地址线测试 & 片内 RAM 和数据访问可用 \\\tablerowrule
DRAM & 初始化参数、训练结果、块读写校验 & 外部内存控制器和板级连线可用 \\\tablerowrule
异常 & 故障地址、异常号、返回 PC & 向量表和异常栈可用 \\\tablerowrule
cache/MMU & 页表摘要、属性、TLB flush 记录 & 地址转换和一致性边界可控 \\\tablerowrule
中断 & 控制器配置、pending/active 状态 & 异步事件链路可用 \\
\bottomrule
\end{longtable}

每条日志都应能回答“这条记录来自哪里”。例如串口能打印字符，不只证明 C 函数被调用，还证明 CPU 能取指、写 MMIO、外设时钟打开、引脚复用正确、波特率近似正确。若后续打开 cache 后打印乱码，就可以把变化集中到 cache 属性、写缓冲、MMIO ordering 或时钟切换，而不是重新怀疑所有层。

\topic{专题：RISC-V 与 x86 的两种指令集风格}
8086 到 Pentium 都属于 x86 家族。RISC-V 是定长的 load/store 架构：基本指令一律 32-bit 宽，只有加载和存储指令访问内存，运算指令的源和结果都在寄存器中，通用寄存器有 x1--x31 共 31 个（x0 恒为 0）。x86-64 则是变长的寄存器-内存混合架构：指令长度从 1 字节到 15 字节不等，一条运算指令可以直接拿内存操作数参与计算，通用寄存器只有 16 个。同一段数组求和循环，两种风格写出来差别很直观。

\begin{lstlisting}
# RISC-V (RV64): a0 = base, a1 = n, sum in a2
      mv   a2, zero           # sum = 0
      mv   t0, zero           # i = 0
loop: bge  t0, a1, done       # if i >= n, exit
      slli t1, t0, 3          # offset = i * 8
      add  t1, a0, t1         # address of a[i]
      ld   t2, 0(t1)          # load a[i]
      add  a2, a2, t2         # sum += a[i]
      addi t0, t0, 1          # i++
      j    loop
done:
\end{lstlisting}

\begin{lstlisting}
; x86-64 (Intel syntax): rdi = base, rsi = n, sum in rax
      xor  rax, rax           ; sum = 0
      xor  rcx, rcx           ; i = 0
loop: cmp  rcx, rsi
      jge  done               ; if i >= n, exit
      add  rax, [rdi + rcx*8] ; sum += a[i], memory operand
      inc  rcx                ; i++
      jmp  loop
done:
\end{lstlisting}

风格差别首先体现在译码器上。RISC-V 译码器拿到 4 字节就能确定操作码、寄存器字段和立即数位置，取指边界固定，多条指令可以成组并行译码；x86 译码器要先扫描前缀、确定指令边界，才知道操作码和操作数在哪里，变长边界识别本身就是一长串组合逻辑，现代 x86 靠预译码标记和微操作 cache 把这笔开销藏起来——Pentium Pro 以来“先译成微操作再乱序执行”的结构，很大程度上就是为了消化变长指令。其次体现在编译器上：load/store 架构把“算”和“访存”拆成两条指令，编译器要把加载尽量提前发射来隐藏内存等待，寄存器多则溢出到栈上的次数少；x86 的内存操作数把加载和加法折叠在一条指令里，代码更短、寄存器压力更小，但这条指令在片内仍要被拆成加载微操作和加法微操作，只是拆分者从编译器换成了硬件。最后是编码密度与译码功耗的取舍：定长指令浪费一些编码空间，换来简单的取指和译码路径；变长指令节省空间和取指带宽，换来更复杂的译码前端。两种风格都能做出高性能实现，差别在于复杂度放在译码硬件里，还是放在编译器调度里。

\topic{专题：ARM 把条件执行和移位操作数带进指令}
经典的 32-bit ARM 指令集（ARM 状态）走的是 x86 与 RISC-V 之间的另一条路：几乎每条数据处理指令都可以带一个条件码后缀，第二个源操作数可以先经过桶形移位器再参与运算。条件码占指令编码最高的 4 bit，取指和译码路径不变，执行时先检查程序状态寄存器的标志：条件不成立，这条指令直接变成空操作，不写结果、不改标志。短小的条件赋值因此不必跳来跳去，\code{moveq}、\code{addne} 这样的指令让短分支在流水线里原地消失。

\begin{lstlisting}
@ ARM: if (r1 == r2) r0 = 1 else r0 = r0 + 1
      cmp   r1, r2
      moveq r0, #1            @ executes only when Z = 1
      addne r0, r0, #1        @ executes only when Z = 0

@ ARM: r0 = r1 + (r2 << 2) in one instruction
      add   r0, r1, r2, lsl #2
\end{lstlisting}

第一个例子里，条件成立与否都会顺序取完这三条指令，没有分支方向要预测，短 if 的代价从“可能猜错分支、冲刷流水线”降为“多占一拍执行槽”；第二个例子里，桶形移位器放在 ALU 第二操作数的入口前，\code{lsl \#2} 不占额外指令，数组下标乘 4、乘 8 这类地址计算被折进单条运算指令。

和 x86-64 对照，能看出设计取向的异同。x86-64 用 \code{cmov} 实现条件赋值，思路与条件执行相同——都是避免短分支——但只有少数指令形式支持，其余场合要靠编译器权衡“分支可能猜错”还是“两条路径都算一遍”；x86 的运算指令不带移位操作数，但复杂寻址模式 \code{[base + index*scale + disp]} 把“下标乘 1/2/4/8、加基址、加偏移”折进了访存指令，与 ARM 的移位操作数解决的是同一类地址计算问题，只是一个挂在加载/存储的地址端，一个挂在 ALU 的数据端。ARM 自己的 64-bit 架构（AArch64）放弃了几乎全集的条件执行，只保留条件选择等少数形式：每条指令都带条件检查，执行控制更复杂，条件码字段还占去本可用于扩大寄存器编号或操作类型的编码位。这些取舍说明，指令集里没有免费的特性，每一项便利都要在别处付出代价。

\topic{专题：主板不是插槽集合，而是平台约定}
主板把 CPU、内存、固件、电源、时钟、PCIe 设备、USB/SATA/NVMe、显卡、网卡和低速管理控制器连成一个可启动平台。它的核心职责是保证复位后每个设备处于可枚举、可供电、可配置、可报告错误的状态。CPU 能跑程序只是平台的一部分；内存训练、PCIe 枚举、电源时序、时钟分配、固件配置和热管理同样决定整机是否可靠。

\begin{pdfdiagram}
  \node[hw state, minimum width=1.55cm] (cpu) at (0,0.8) {CPU};
  \node[hw memory, minimum width=2.2cm] (dram) at (2.6,1.65) {DDR 内存};
  \node[hw compute, minimum width=1.9cm] (pch) at (2.6,-0.35) {芯片组/PCH};
  \node[hw control, minimum width=1.7cm] (fw) at (0,-1.35) {UEFI/Flash};
  \node[hw control, minimum width=1.65cm] (pcie) at (5.2,0.8) {PCIe Root};
  \node[hw state, minimum width=1.55cm] (gpu) at (7.5,1.55) {显卡};
  \node[hw control, minimum width=1.55cm] (nvme) at (7.5,0.15) {NVMe};
  \node[hw control, minimum width=1.55cm] (usb) at (5.2,-1.35) {USB/SATA/网卡};
  \draw[hw bus] (cpu) -- node[above,hw label,fill=white]{内存控制器} (dram);
  \draw[hw bus] (cpu) -- node[left,hw label,fill=white]{DMI/片上互连} (pch);
  \draw[hw bus] (fw) -- node[left,hw label,fill=white]{启动固件} (pch);
  \draw[hw bus] (cpu) -- (pcie);
  \draw[hw bus] (pcie) -- (gpu);
  \draw[hw bus] (pcie) -- (nvme);
  \draw[hw bus] (pch) -- (usb);
\end{pdfdiagram}
\diagramnote{PC 主板的关键不是“很多线”，而是启动、枚举、供电、时钟和错误报告共同构成平台约定。}

\begin{longtable}{L{0.18\linewidth}L{0.36\linewidth}L{0.36\linewidth}}
\toprule
主板对象 & 负责的问题 & 典型失败表现 \\
\midrule
供电 VRM & 把电源转换成 CPU/GPU/内存所需电压，并按时序上电 & 复位不释放、负载变化时掉压、机器重启 \\\tablerowrule
时钟源 & 给 CPU、内存、PCIe、USB 等提供参考时钟 & 链路训练失败、设备枚举不稳定 \\\tablerowrule
固件 Flash & 保存 UEFI/BIOS、微码、平台配置 & 无法取第一阶段代码或配置错误 \\\tablerowrule
DDR 通道 & 连接内存颗粒或 DIMM，完成训练和 ECC & 内存训练失败、随机数据错误 \\\tablerowrule
PCIe Root & 枚举显卡、NVMe、网卡和扩展设备 & 设备消失、降速、AER 错误 \\\tablerowrule
低速管理 & Super I/O、EC、传感器、风扇、RTC & 键盘、风扇、温度或睡眠唤醒异常 \\
\bottomrule
\end{longtable}

\topic{专题：PCIe 是分组互连，不是并行总线升级版}
早期外部总线常暴露地址线、数据线、读写控制和 READY；PCIe 则把事务封装成包，在点对点串行链路上传输。CPU 或芯片组通过 root complex 发起配置、内存读写和消息事务；设备通过 BAR 暴露 MMIO 窗口，通过 MSI/MSI-X 发送中断，通过 DMA 访问主存。PCIe 的关键边界包括链路训练、配置空间、BAR 映射、事务层包、完成包和错误报告。

\begin{longtable}{L{0.18\linewidth}L{0.38\linewidth}L{0.34\linewidth}}
\toprule
PCIe 概念 & 含义 & 本书对应概念 \\
\midrule
配置空间 & 枚举设备、读取 ID、配置 BAR 和能力 & 地址地图和平台发现 \\\tablerowrule
BAR & 把设备寄存器映射到 CPU 地址空间 & MMIO 副作用和访问属性 \\\tablerowrule
Memory Read/Write TLP & 设备或 root 发起的内存事务包 & 总线事务和 completion \\\tablerowrule
MSI/MSI-X & 设备通过写内存样式消息触发中断 & 中断不一定是独立引脚 \\\tablerowrule
AER & 高级错误报告 & 错误路径可诊断 \\
\bottomrule
\end{longtable}

把 PCIe 讲清楚后，显卡、NVMe、网卡都可以归入同一模型：先由平台枚举设备并分配地址窗口；驱动把命令写入队列或 MMIO；设备通过 DMA 读写主存；完成后写 completion 或发 MSI；错误通过状态寄存器、completion 状态或 AER 上报。设备是能主动访问内存的总线主设备。

\topic{专题：显卡的核心不是显示器接口，而是并行处理器加显存系统}
显卡既是显示设备，也是大规模并行计算设备。GPU 内部有许多执行单元、调度器、寄存器文件、L1/L2 cache、纹理/采样单元、光栅化和显示引擎；外部或封装上有高带宽显存，例如 GDDR 或 HBM。CPU 不会逐像素把画面写到显示器，而是准备命令缓冲区、资源描述、纹理和顶点数据，GPU 读取这些对象并在显存中生成帧缓冲，显示引擎再按扫描时序把帧送到显示接口。

\begin{pdfdiagram}
  \node[hw state, minimum width=1.6cm] (cpu) at (0,0) {CPU};
  \node[hw memory, minimum width=2.0cm] (sysmem) at (2.4,0.8) {系统内存};
  \node[hw control, minimum width=2.0cm] (cmd) at (2.4,-0.8) {命令缓冲};
  \node[hw control, minimum width=1.85cm] (pcie) at (4.9,0) {PCIe};
  \node[hw compute, minimum width=2.0cm] (gpu) at (7.35,0.8) {GPU 执行};
  \node[hw memory, minimum width=1.75cm] (vram) at (9.7,0.8) {VRAM};
  \node[hw state, minimum width=1.75cm] (scan) at (7.35,-0.8) {显示引擎};
  \node[hw state, minimum width=1.55cm] (disp) at (9.7,-0.8) {显示器};
  \draw[hw bus] (cpu) -- (sysmem);
  \draw[hw bus] (cpu) -- (cmd);
  \draw[hw bus] (cmd) -- (pcie);
  \draw[hw bus] (pcie) -- (gpu);
  \draw[hw bus] (gpu) -- (vram);
  \draw[hw return] (vram) -- (scan);
  \draw[hw bus] (scan) -- (disp);
\end{pdfdiagram}
\diagramnote{GPU 数据路径：CPU 提交命令，GPU 经 PCIe/内存读取资源，在显存中生成帧，显示引擎独立扫描输出。}

\begin{longtable}{L{0.18\linewidth}L{0.38\linewidth}L{0.34\linewidth}}
\toprule
显卡对象 & 硬件含义 & 常见误解 \\
\midrule
命令缓冲 & CPU 写入、GPU 异步消费的工作队列 & 不是 CPU 直接调用 GPU 函数 \\\tablerowrule
VRAM & GPU 近处的大带宽存储系统 & 不是普通 DRAM 的别名，带宽和访问模式不同 \\\tablerowrule
着色/计算单元 & 大量相似线程并行执行 & 不是少数超快 CPU 核心 \\\tablerowrule
纹理/cache 单元 & 优化二维或三维局部访问和采样 & 不等于通用 cache 的简单复制 \\\tablerowrule
显示引擎 & 按固定时序读取帧缓冲并输出信号 & 渲染完成和显示扫描是两个边界 \\
\bottomrule
\end{longtable}

GPU 的性能问题也不同于 CPU。CPU 常追求低时延和强单线程控制流；GPU 常追求吞吐，用大量并行线程隐藏显存等待。分支发散、访存不合并、显存带宽不足、CPU/GPU 同步过多，都会让 GPU 性能掉下来。理解显卡时应问：命令何时提交，资源何时对 GPU 可见，GPU 何时完成，帧缓冲何时可显示，CPU 是否过早读回结果。

\topic{专题：整机案例：从打开程序到读取文件和显示图像}
一个真实用户动作会穿过多种硬件。用户打开程序，操作系统可能从 SSD 读取可执行文件和页面；CPU 触发块设备队列；NVMe 控制器 DMA 数据到主存；页表和 cache 让 CPU 取指执行；程序提交图形命令；GPU 读取资源、渲染到 VRAM；显示引擎扫描帧缓冲。这里没有哪个部件可以单独解释全程，必须把 CPU、内存、SSD、PCIe、GPU、显示和中断连成 trace。

\begin{longtable}{L{0.16\linewidth}L{0.44\linewidth}L{0.3\linewidth}}
\toprule
阶段 & 硬件动作 & 完成标志 \\
\midrule
缺页/读文件 & OS 提交 NVMe 读请求，设备 DMA 到主存 & 块 I/O completion 和页状态有效 \\\tablerowrule
执行代码 & CPU 取指、TLB/cache 查找、分支和写回 & 指令提交，异常路径为空 \\\tablerowrule
准备图形资源 & CPU 写命令缓冲和资源描述符 & 命令队列对 GPU 可见 \\\tablerowrule
GPU 渲染 & GPU 读取资源，执行 shader，写 VRAM & fence 或 completion 表示渲染完成 \\\tablerowrule
显示扫描 & 显示引擎按刷新时序读取帧缓冲 & vblank、翻页或显示控制器状态 \\
\bottomrule
\end{longtable}

把每一步的机器行为写清楚：谁拥有缓冲区，谁发起 DMA，哪个地址空间有效，cache/TLB 在哪里参与，哪个 completion 才说明完成，哪个错误会阻止下一阶段。

\topic{专题：用仿真器观察硬件行为}
纸面规格之外，QEMU 这类全系统仿真器给了读者另一双眼睛。它把目标机器的指令逐块翻译成宿主机能执行的操作：目标机的寄存器是宿主机内存里的数据结构，目标机的物理地址空间是一张查找表，读 RAM 变成读宿主机上的数组，写 MMIO 地址则变成对一段模拟设备代码的调用，由它更新设备状态、必要时触发模拟中断。于是“一台 PC”可以只是宿主机上的一个进程，但它照样从复位向量取指、跑 BIOS、读磁盘镜像。

接上调试器（QEMU 内置 monitor 或 GDB 调试桩）后，可以从复位后的第一条指令开始单步，看 BIOS 怎样建立段描述符和中断向量，看启动扇区怎样用 \code{int 0x13} 读盘；可以在任意时刻打印全部寄存器和任意一段物理内存，改一个字节再继续运行，直接观察“数据段复制错了”会长成什么样；还可以故意让代码访问未映射地址，看仿真器报告总线错误或目标机异常，而不必担心烧坏任何东西。

但仿真结果不能全部当成硬件结论。第一，仿真器按功能正确性建模，不按时间建模：一次访存在真实 DRAM 上要等几十纳秒，在仿真器里和一次寄存器加法几乎一样快，cache miss、等待状态、DRAM 刷新和总线仲裁的时序开销都不存在，性能结论无从谈起。第二，电气层面的问题整个不可见：复位时序、电源爬升、信号完整性、三态冲突、虚焊和接触不良，这些真实 bring-up 中最常见的故障在仿真器里没有对应物。第三，设备模型是简化的，作者没有建模的寄存器行为和时序敏感的错误路径，仿真器可能永远表现为“工作正常”。所以仿真器和真实硬件是互补关系：仿真器适合回答地址地图、启动顺序、异常路径这类“逻辑对不对”的问题；时序裕量、电气可靠性这类“电路稳不稳”的问题，仍要回到真实板级平台，用示波器和日志回答。

\topic{专题：纸面项目应替代不可执行的大实验}
如果没有硬件平台，仍然可以做高质量项目。项目形式应从“搭出来能跑”改成“规格写得可检查”。例如给出一台简化 PC：一个 CPU、两级 cache、DDR 控制器、PCIe root、NVMe、GPU、USB 控制器和固件 ROM。读者要写地址地图、启动顺序、设备枚举、一次文件读取、一次 GPU 命令提交、一次中断处理和一个错误路径。这样的纸面项目比空泛地说“了解主板”更有效。

\begin{longtable}{L{0.2\linewidth}L{0.46\linewidth}L{0.24\linewidth}}
\toprule
纸面项目 & 必须交付 & 核对要点 \\
\midrule
简化 PC 地址地图 & DRAM、固件、MMIO、PCIe BAR、未映射区 & 任意地址有明确目标或错误 \\\tablerowrule
NVMe 读路径 & submission、doorbell、DMA buffer、completion、中断 & CPU 不把门铃返回当成 I/O 完成 \\\tablerowrule
GPU 命令路径 & 命令缓冲、资源可见性、VRAM、fence、显示扫描 & 渲染完成和显示完成分开 \\\tablerowrule
主板启动路径 & 电源、时钟、复位、固件、内存训练、PCIe 枚举 & 每一步有失败后可诊断状态 \\\tablerowrule
错误注入 & TLB miss、cache miss、PCIe 错误、SSD 超时、GPU hang & 错误归属到具体层，不混成“程序卡住” \\
\bottomrule
\end{longtable}

\section*{章末练习}

\begin{chapterexercise}{ELF 节表读法}
对 \code{gcc -c} 生成的 \code{hello.o} 执行 \code{readelf -S}，找到 \code{.text}、\code{.rodata}、\code{.bss} 三节，记下各自的类型与大小，并解释 \code{.bss} 为什么在文件里一个字节也不占
\exercisecheck{学习目标}{分清 \code{PROGBITS} 与 \code{NOBITS}}
\end{chapterexercise}

\begin{chapterexercise}{重定位演算}
一条 \code{call} 指令的起始地址为 \code{0x401120}，其 32 位位移字段的地址 P 为 \code{0x401121}；目标符号地址 S 为 \code{0x401130}，重定位类型 \code{R\_X86\_64\_PC32}，加数 A 为 \code{-4}。算出写入该字段的值，并用执行时的 \code{rip} 验证结果
\exercisecheck{学习目标}{重定位就是解一次 S+A-P}
\end{chapterexercise}

\begin{chapterexercise}{静态库顺序}
\code{main.o} 用到 \code{libfoo.a} 里的函数 \code{foo}：比较 \code{gcc main.o -lfoo} 与 \code{gcc -lfoo main.o} 两条命令，指出哪一条会报未定义符号错误，并解释原因
\exercisecheck{学习目标}{库按需求出，顺序敏感}
\end{chapterexercise}

\begin{chapterexercise}{PLT/GOT 走读}
程序先后两次调用 \code{write}：分别写出两次调用经过 PLT 桩与 GOT 表项的路径差别，并指出 GOT 表项在哪一步、被谁改写
\exercisecheck{学习目标}{第一次解析，之后间跳直达}
\end{chapterexercise}

\begin{chapterexercise}{syscall 参数}
按 x86-64 系统调用约定，写出 \code{write(1, buf, 3)} 陷入内核前 \code{rax}、\code{rdi}、\code{rsi}、\code{rdx} 各放什么；再说明第四个参数为什么用 \code{r10} 而不是 \code{rcx}
\exercisecheck{学习目标}{编号走 \code{rax}，参数按固定寄存器序}
\end{chapterexercise}

\begin{chapterexercise}{编译链四阶段}
手头只有一个 \code{hello.c}：分别写出加 \code{-E}、\code{-S}、\code{-c} 与不加开关的四条 \code{gcc} 命令的产物文件名与性质，并说明哪个产物还不能直接执行、缺哪一步
\exercisecheck{学习目标}{四把开关各停一站}
\end{chapterexercise}

\section*{参考解答与要点}

\begin{chapteranswer}{ELF 节表读法}
\code{.text} 类型 \code{PROGBITS}，是指令字节本身；\code{.rodata} 同为 \code{PROGBITS}，本例装着字符串常量共 3 字节；\code{.bss} 类型 \code{NOBITS}，只有大小声明、没有内容——未初始化数据按约定在装载时清零，文件里存一串零纯属浪费，所以它只占节表中的一个条目，不占文件偏移
\answercheck{必须掌握的要点}{\code{NOBITS} 只声明大小，装载时补零}
\end{chapteranswer}

\begin{chapteranswer}{重定位演算}
写入值等于 S+A-P，即 \code{0x401130} 减 4 再减 \code{0x401121}，得 \code{0xb}，按小端写入四字节，依次为 \code{0b}、\code{00}、\code{00}、\code{00}。执行时 CPU 取完这条 5 字节指令后 \code{rip} 为 \code{0x401125}（即 P+4），跳转目标等于 \code{0x401125} 加 \code{0xb}，得 \code{0x401130}，正是 S，演算闭合
\answercheck{必须掌握的要点}{PC 相对的基准是下一条指令}
\end{chapteranswer}

\begin{chapteranswer}{静态库顺序}
\code{gcc -lfoo main.o} 报错。链接器从左向右扫描：处理 \code{-lfoo} 时未定义符号集合还是空的，库中任何成员都不会被取入；轮到 \code{main.o} 才引入未定义的 \code{foo}，可库已经扫过，于是报 undefined reference。正确写法是把库放在引用它的目标文件之后：\code{gcc main.o -lfoo}——扫到库时 \code{foo} 正被需要，定义它的成员才被取出
\answercheck{必须掌握的要点}{引用在前、库在后}
\end{chapteranswer}

\begin{chapteranswer}{PLT/GOT 走读}
第一次：\code{call} 进入 \code{write} 的 PLT 桩，桩首条指令经 GOT 表项间接跳转，而此时表项里存的是桩内下一条指令的地址——于是转回桩中，压入重定位序号，再跳公共解析桩；动态链接器算出 \code{write} 的真实地址、改写该 GOT 表项，最后才转入 \code{write} 本体。第二次：GOT 已是真实地址，桩里那一次间接跳转直达目标，解析路径被完全绕过。改写者是动态链接器的解析例程，时机是第一次调用期间
\answercheck{必须掌握的要点}{延迟绑定：第一次付费，以后一跳}
\end{chapteranswer}

\begin{chapteranswer}{syscall 参数}
\code{rax} 放编号，\code{write} 是 1；\code{rdi} 放 1（文件描述符），\code{rsi} 放 \code{buf} 的地址，\code{rdx} 放 3（长度）。\code{syscall} 指令本身把返回地址存进 \code{rcx}、把标志存进 \code{r11}，这两个寄存器被硬件占用，所以内核约定第四个参数改用 \code{r10}；完整顺序是 \code{rdi}、\code{rsi}、\code{rdx}、\code{r10}、\code{r8}、\code{r9}，返回值由 \code{rax} 带回
\answercheck{必须掌握的要点}{\code{rcx} 与 \code{r11} 被指令占用}
\end{chapteranswer}

\begin{chapteranswer}{编译链四阶段}
\code{-E} 得 \code{hello.i}：宏与头文件展开后的纯 C 文本；\code{-S} 得 \code{hello.s}：x86-64 汇编文本；\code{-c} 得 \code{hello.o}：可重定位目标文件，机器码已生成但外部符号只有名字没有地址；不加开关得可执行文件（\code{-o hello} 时名为 \code{hello}，缺省 \code{a.out}）。\code{hello.o} 不能直接执行：缺链接这一步——拼入启动文件、解析 \code{write}、填入口地址，三件事都在链接阶段完成
\answercheck{必须掌握的要点}{只有链接产物是给内核装载的}
\end{chapteranswer}
